% =====================================================================
%  Proof-of-Policy --- Supplementary Material (standalone IEEEtran)
%
%  Compile (two passes, quiet mode):
%      pdflatex "\def\SKELQUIET{}% =====================================================================
%  Proof-of-Policy --- Supplementary Material (standalone IEEEtran)
%
%  Compile (two passes, quiet mode):
%      pdflatex "\def\SKELQUIET{}% =====================================================================
%  Proof-of-Policy --- Supplementary Material (standalone IEEEtran)
%
%  Compile (two passes, quiet mode):
%      pdflatex "\def\SKELQUIET{}% =====================================================================
%  Proof-of-Policy --- Supplementary Material (standalone IEEEtran)
%
%  Compile (two passes, quiet mode):
%      pdflatex "\def\SKELQUIET{}\input{tdsc-supplementary.tex}"
%
%  This document carries the material removed from the 12--14 page main
%  paper to meet the TDSC length limit: full proofs, the detailed
%  evaluation tables, policy-private-mode detail, related-work tables,
%  and reproduction instructions. It is self-contained: it restates the
%  results it proves and can be read without the main PDF.
% =====================================================================
\documentclass[10pt,journal,compsoc]{IEEEtran}

\usepackage{amsmath,amsfonts,amssymb,amsthm}
\usepackage{algorithmic}
\usepackage{algorithm}
\usepackage{array}
\usepackage[caption=false,font=normalsize,labelfont=sf,textfont=sf]{subfig}
\usepackage{textcomp}
\usepackage{stfloats}
\usepackage{url}
\usepackage{verbatim}
\usepackage{graphicx}
\usepackage{cite}
\usepackage{booktabs}
\usepackage{xcolor}
\usepackage{microtype}
\usepackage{hyperref}

\hyphenation{Proof-of-Policy zero-knowl-edge re-exec-ution pol-icy
com-pli-ance bind-ing veri-fi-a-bil-ity cer-ti-fi-cate non-zero-knowl-edge}

\setlength{\emergencystretch}{2em}

\newcolumntype{L}[1]{>{\raggedright\arraybackslash}p{#1}}
\newcommand{\tabfont}{\footnotesize}

\newtheorem{theorem}{Theorem}
\newtheorem{lemma}[theorem]{Lemma}
\newtheorem{proposition}[theorem]{Proposition}
\newtheorem{corollary}[theorem]{Corollary}
\newtheorem{definition}[theorem]{Definition}
\newtheorem{assumption}[theorem]{Assumption}
\newtheorem{remark}[theorem]{Remark}

% Work markers are no-ops in the supplementary material.
\newcommand{\need}[1]{}
\newcommand{\fixme}[1]{}
\newcommand{\wpmark}[1]{}

% Supplementary numbering: S1, S2, ... and Table S1, S2, ...
\renewcommand{\thesection}{S\arabic{section}}
\renewcommand{\thetable}{S\arabic{table}}

\begin{document}

\title{Supplementary Material for\\``Proof-of-Policy: Zero-Knowledge Compliance
Certificates for LLM Agent Outputs''}
\author{Author Name(s)}
\maketitle

\begin{abstract}
This supplement contains the full proofs, the detailed evaluation tables, the
policy-private-mode tables, the related-work positioning, and the reproduction
instructions omitted from the main paper for length. Section and table numbers
are prefixed with ``S''. The main paper states every result; this document
supplies the arguments and the per-configuration numbers behind them.
\end{abstract}

% =====================================================================
\section{Full Proofs}
\label{app:proofs}

This section restates the results proved in the main paper and gives complete
proofs. We use the main paper's notation; for convenience, the two objects the
proofs depend on are repeated here.

\begin{definition}[Guessing entropy; restated from the main paper]
\label{def:k}
For a secret distribution $\mathcal{D}$ and adversary side information $Z$, let
\begin{multline*}
\mathsf{Guess}(X \mid Z) \;\triangleq\; \max_{\mathcal{A}} \Pr[\mathcal{A}(Z) = X]\\
\;=\; \mathbb{E}_z\!\left[\max_x \Pr[X = x \mid Z = z]\right],
\end{multline*}
\[
k_X \;\triangleq\; -\log_2 \mathsf{Guess}(X \mid Z),
\]
the maximum ranging over all time-unbounded adversaries.
\end{definition}

\begin{assumption}[Adversary side information; restated from the main paper]
\label{as:Z}
The adversary knows the prompt, the model (possibly its weights), and the
decoding parameters.
\end{assumption}

\begin{proposition}[Conditional trilemma; restated from the main paper]
\label{thm:t1}
Fix $\lambda,\varepsilon>0$ and $q\ge 1$. Let $X$ be a secret with guessing
entropy $k_X<\log q+\log(1/\varepsilon)$ whose binding is \emph{guess-tight}:
$\beta_\rho(X)=H(\mathsf{DOM}\|\rho\|X)$ for a random oracle $H$ and public
randomness $\rho$ independent of $X$ and $Z$. Then:
\begin{enumerate}
\item[(1)] \emph{Public binding.} For any adversary that evaluates $\beta$ at
most $q$ times,
\[
\Pr\!\left[\mathcal{A}^\beta(Z,\beta(X))=X\right]
\le q\cdot\mathsf{Guess}(X\mid Z)+\mathrm{negl}(\lambda)
= q\cdot 2^{-k_X}+\mathrm{negl}(\lambda).
\]
Hiding with advantage $\varepsilon$ is \emph{certified} by this bound when
$k_X\ge\log q+\log(1/\varepsilon)$; for a skewed distribution this is a
sufficient condition, not a matching attack. The integer $q$ counts evaluations
of $\beta$, not time: the adversary may be computationally unbounded.
\item[(2)] \emph{Keyed binding.} For $\beta_s(X)=H(s,X)$ with $s$ held only by a
designated party and absent from the public values,
$\Pr[\mathcal{A}(Z,b)=X]\le\mathsf{Guess}(X\mid Z)+\mathrm{negl}(\kappa)$.
What the keyed form buys is hiding when
$k_X<\log q+\log(1/\varepsilon)$; its price is PV-2.
\item[(3)] \emph{Conditional trilemma.} Consider schemes whose binding is
publicly recomputable (PV-2), whose hiding is exact recovery, and whose binding
is guess-tight. Within the scheme class defined above, when
$k_X<\log q+\log(1/\varepsilon)$ such a scheme cannot satisfy binding, hiding
and PV-2 simultaneously.
\end{enumerate}
\end{proposition}

\begin{proof}[Proof of Proposition~\ref{thm:t1}]
\emph{(1).} Fix $z$ and let the adversary's query set be
$Q = \{q_1, \dots, q_q\}$. Decompose by first hit,
\[
\Pr[\exists i : q_i = X] \;=\; \sum_i \Pr\!\left[q_i = X \wedge \neg E_{<i}\right],
\]
where $E_{<i} \triangleq \{\exists j < i : q_j = X\}$.
Conditioned on $\neg E_{<i}$, every oracle answer $a_1,\dots,a_{i-1}$ already
revealed is uniform and independent of $X$ (no queried point equals $X$, and $H$
is a fresh random function), and $\beta(X)$ itself is uniform and independent of
$X$. Hence $q_i$---a deterministic function of $(z, \beta(X), a_{<i})$---is
\emph{independent of $X$}, so
\begin{multline*}
\Pr[q_i = X \mid Z = z]
= \sum_x \Pr[q_i = x]\,\Pr[X = x \mid Z = z] \\
\le \max_x \Pr[X = x \mid Z = z] = 2^{-k_z},
\end{multline*}
giving $\Pr[\exists i : q_i = X \mid z] \le q\,2^{-k_z}$. Taking the expectation
over $z$ and using $\mathbb{E}_z[2^{-k_z}] = \mathsf{Guess}(X \mid Z) = 2^{-k_X}$
from Definition~\ref{def:k} yields $q \cdot 2^{-k_X}$. An adversary that outputs
an unqueried value wins with probability at most $2^{-k_z} \le q\,2^{-k_z}$,
which the same bound absorbs. No $q\cdot2^{-\lambda}$ term is needed in the
exact-recovery game: a query that lands on a different $x$ with
$\beta(x)=y$ is a second preimage of $H$, not a win. The explicit
$q\cdot2^{-\lambda}$ is kept only when the win condition is consistency with the
published $b$ rather than exact recovery, or when $q$ is unbounded.

\emph{(2).} Without an oracle for $\beta_s$, the value $b=H(s,X)$ is uniform and
independent of $X$ in the random-oracle model for an adversary that does not know
$s$, so the best strategy is the prior mode and succeeds with probability
$\mathsf{Guess}(X\mid Z)$. An attempted enumeration must also guess $s$, at cost
$2^{-\kappa}$ per query. The bound is therefore
$\mathsf{Guess}(X\mid Z)+\mathrm{negl}(\kappa)$, with no $q$-fold amplification.
Part (2) buys hiding only below the public-form threshold
$\log q+\log(1/\varepsilon)$; when $k_X=0$ the prior mode already recovers $X$
and no binding helps.

\emph{(3).} Public recomputability gives (1). If (1) exceeds $\varepsilon$, the
only way to avoid it is a binding that a verifier cannot recompute, i.e.\ a keyed
binding, which by part (2) forfeits PV-2. The third leg---hiding plus PV-2
without binding---is excluded from the class of compliance schemes because a
certificate whose binding does not pin $X$ certifies nothing about the delivered
artifact.
\end{proof}

\noindent\emph{Preimage-class bound.} For an arbitrary deterministic public
$\beta$, conditioning on $\beta(X)$ gives
$\Pr[\mathcal{A}^\beta(Z,\beta(X))=X]\le
\mathbb{E}_y[\max_x\Pr[X=x\mid Z,\beta(X)=y]]$, with no hypothesis on $\beta$.
For $\beta=\mathrm{id}$ or an injective $\mathsf{encode}$ this equals $1$; for a
lossy encoding it equals the inverse of the preimage-class size.

\begin{theorem}[D-privacy from three conditions; restated from the main paper]
\label{thm:dpriv}
The scheme is $D$-private if (C1) there is a PPT simulator $S_{\mathrm{rest}}$
such that, for every admissible $(\pi,T,n)$, the public values other than the
commitment satisfy
$(h_\pi,S_{\mathrm{rest}}(pp,h_\pi,D(\pi),\mathrm{pub}_T,n))\approx_c
L_{\mathrm{pub}}(\pi,T,n)$; (C2) the proof system is witness-hiding, i.e.\
$L_{\mathrm{proof}}$ has a simulator that takes only the public statement;
and (C3) the policy commitment is hiding for all $D$-equal policies,
$(pp,h_{\pi_0})\approx_c(pp,h_{\pi_1})$. The advantage is
$O(\varepsilon_{\mathrm{pub}}+\varepsilon_{\mathrm{wh}}+\varepsilon_{\mathrm{com}})$.
\end{theorem}

\begin{proof}[Proof of Theorem~\ref{thm:dpriv}]
Fix an admissible tuple $(\pi_0,\pi_1,T,n)$ with $D(\pi_0)=D(\pi_1)$ and
$R(\pi_0,T)=R(\pi_1,T)$. Write $h_i=h_{\pi_i}$,
$R_i=L_{\mathrm{rest}}(\pi_i,T,n)$, $P_i=L_{\mathrm{proof}}(\pi_i,T,n)$ and
$\mathsf{Stmt}_i=(h_i,R_i)$. Consider the hybrids
\begin{align*}
H_0&=(h_0,R_0,P_0),\\
H_1&=(h_0,R_0,S_{\mathrm{proof}}(\mathsf{Stmt}_0)),\\
H_2&=(h_0,S_{\mathrm{rest}}(h_0,D,\mathrm{pub}_T,n),S_{\mathrm{proof}}(\mathsf{Stmt}_0)),\\
H_3&=(h_1,S_{\mathrm{rest}}(h_1,D,\mathrm{pub}_T,n),S_{\mathrm{proof}}(\mathsf{Stmt}_1)),\\
H_4&=(h_1,R_1,S_{\mathrm{proof}}(\mathsf{Stmt}_1)),\\
H_5&=(h_1,R_1,P_1).
\end{align*}
$H_0\approx H_1$ and $H_4\approx H_5$ cost $\varepsilon_{\mathrm{wh}}$ each by
C2; $H_1\approx H_2$ and $H_3\approx H_4$ cost
$\varepsilon_{\mathrm{pub}}$ each by C1. For $H_2\approx H_3$, split into
$H_{2a}=(h_1,S_{\mathrm{rest}}(h_1,D,\mathrm{pub}_T,n),S_{\mathrm{proof}}(\mathsf{Stmt}_0))$
and $H_{2b}=H_3$: $H_2\approx H_{2a}$ costs $\varepsilon_{\mathrm{com}}$ by C3,
and $H_{2a}\approx H_{2b}$ follows because
$\mathsf{Stmt}_0\approx_c\mathsf{Stmt}_1$ (C1 and C3 by triangle inequality) and
$S_{\mathrm{proof}}$ is a PPT function of the statement. The triangle inequality
gives $4\varepsilon_{\mathrm{pub}}+2\varepsilon_{\mathrm{wh}}+2\varepsilon_{\mathrm{com}}$,
which is $O(\varepsilon_{\mathrm{pub}}+\varepsilon_{\mathrm{wh}}+\varepsilon_{\mathrm{com}})$.
The game fixes one $T$ because two different artifacts are distinguished by
$B(n,\cdot)$ itself, which is response-side leakage and not a policy advantage.
\end{proof}

\begin{theorem}[Leakage-profile decomposition; restated from the main paper]
\label{thm:leak}
Let $V=(L_{\mathrm{pub}},L_{\mathrm{proof}},L_{\mathrm{side}})$ be the verifier's
view, where $L_{\mathrm{pub}}$ is the certificate's public values,
$L_{\mathrm{proof}}$ is what the proof artefact reveals about the witness
$w=(\pi,T,\mathrm{trace})$, and $L_{\mathrm{side}}$ is the adversary's prior
information $Z$. For any $q\ge 1$ and any adversary that evaluates $\beta$ at
most $q$ times,
\[
\Pr\!\left[\mathcal{A}^\beta(V)=X\right]
\;\le\; q\cdot\mathsf{Guess}\!\left(X\mid L_{\mathrm{pub}},L_{\mathrm{proof}},L_{\mathrm{side}}\right)
+\mathrm{negl}(\lambda),
\]
with the final output charged as one evaluation.
\end{theorem}

\noindent\emph{Proof of Theorem~\ref{thm:leak}.} The proof is the first-hit
decomposition of Proposition~\ref{thm:t1}(1) with $(Z,\beta(X))$ replaced by the
verifier view $V$. The step that changes is the uniformity of answers to queries
that have not hit $X$: it holds conditionally on $L_{\mathrm{pub}}$ and fails when
$L_{\mathrm{proof}}$ is a function of the witness. A witness-hiding tier restores
$L_{\mathrm{proof}}=\bot$ and hence the public-value bound; a \texttt{core}
artefact does not, and the bound can reach $1$.

\begin{theorem}[Delivery binding; restated from the main paper]
\label{thm:t3}
Fix a policy $\pi$. A delivery adversary $\mathcal{B}$ may, \emph{before}
seeing the receiver's challenge, precompute a set $S$ of accepted transcripts,
where each transcript covers at most $c$ challenges. The receiver then
contributes $f$ fresh, uniformly random bits $r$ that enter the proven
statement. Under the certificate's soundness (D1) and challenge binding
(D2/B-bind), the divergence advantage satisfies
\[
\mathsf{Adv}_{\mathrm{del}}(\mathcal{B}) \;\le\; |S| \cdot c \cdot 2^{-f}
\;+\; \mathrm{negl}(\lambda).
\]
The replay-only probability---acceptance without requiring divergence---is at
most $|S|\cdot c\cdot 2^{-f}$. If the receiver contributes no fresh randomness
($f = 0$; e.g.\ an empty nonce), a precomputed transcript is accepted with
probability $1$ and \emph{does not decay in any parameter of the system}.
\end{theorem}

\begin{proof}[Proof of Theorem~\ref{thm:t3}]
The receiver's freshness $r$ enters the proven statement through
$b = H(\mathsf{dom} \,\|\, r \,\|\, T)$. For a certificate in $S$ to be accepted
on $r$, the value $r$ must fall inside that certificate's challenge set
$A(\mathrm{cert})$, with $|A(\mathrm{cert})|\le c$; a union bound over $S$ gives
$\sum_{\mathrm{cert}\in S}|A(\mathrm{cert})|/2^f\le |S|\cdot c\cdot 2^{-f}$.
A certificate minted \emph{after} $r$ is seen can diverge only by exhibiting two
distinct preimages of the same $b$ under $B(r,\cdot)$, which is exactly a
challenge-binding failure and costs $\mathrm{negl}(\lambda)$. When $f = 0$ every
element of $S$ is a permanently valid replay, so the replay probability is $1$
and does not decay in any parameter of the system. The $f = 0$ branch is not
hypothetical: \texttt{circuits/\allowbreak types/\allowbreak src/\allowbreak
lib.rs} (lines 950--995) documents an empty nonce as a legal input that
``provides no replay protection.'' The $f=0$ separation holds under stateless
public verification with no external mechanism (no revocation table, trusted
clock, or hardware). The $c$ factor is tight when the challenge sets of the
transcripts in $S$ are disjoint and of size $c$.
\end{proof}

\begin{proposition}[Quadratic cost of the naive matcher]
\label{prop:naive}
For the pattern \texttt{a+b} and input \texttt{a}$^n$, the naive backtracking
matcher performs
$\sum_{i=1}^{n}(n-i+1)=n(n+1)/2=\Theta(n^{2})$ NFA steps, whereas the Pike VM
performs $\Theta(n)$. The adversarial input family is admissible because the
response bytes are chosen by the party being checked. This bounds the cost of
\emph{that} baseline, not of policy-compliance matching: the Pike VM is linear
on the same family, and no lower bound for the problem is claimed.
\end{proposition}
\begin{proof}
No input position contains \texttt{b}, so no start position can fail or accept
before the end of the string. The start at $i$ therefore takes exactly
$n-i+1$ steps, and the sum is $n(n+1)/2$. The Pike VM keeps one state per NFA
position and scans once.
\end{proof}

\begin{proposition}[NFA contract]
\label{prop:nfa}
For the supported regular-expression subset $P$---ASCII, ordered and disjoint
character-class ranges, Unicode scalar values, and the same compiled policy on
both sides---the Thompson construction preserves the language, the Pike VM
implements unanchored substring search, and the Python reference and Rust
circuit matchers agree on the same compiled NFA specification. Empty character
classes are rejected at parse time (fail-closed), so ``non-empty character
class'' is a parser contract rather than an extra theorem assumption.
\end{proposition}
\begin{proof}[Proof sketch]
Language preservation is the standard Thompson construction over the restricted
syntax. The Pike VM simulates all NFA states at one input position and advances
once per character, so it decides unanchored search in one pass. The two
matchers consume the same \texttt{ConstraintSpec} bytes and implement the same
transition relation; agreement is checked by an algorithm-level bisimulation
plus differential fuzzing against Python's \texttt{re}, not by line-by-line
formal verification.
\end{proof}

% =====================================================================
\section{Evaluation Details}
\label{app:eval-details}

This section carries the per-configuration tables and the detailed readings
omitted from the main paper. Every number is the content of the checked-in
result file named in the caption, as in the main paper.

\subsection{Cycle sweep: length versus rule count}
\label{sec:eval-cycles}

Response length was pushed from 200 to 100{,}000 characters and the rule count
from 1 to 6, giving 40 measurement points (\texttt{bench/bench\_cycles.py}; the
full matrix is in \texttt{bench/results/cycles.md}). The response text comes from
an evaluation artefact the repository's own end-to-end demo produced (372
characters, three segments deduplicated by content); longer points are tiled from
it. What is measured is therefore the \emph{curve} of scanning cost, not ``a real
100{,}000-character document''. Every point was first confirmed to match no rule:
a matching rule short-circuits, and the cost of a short-circuit says nothing.

\begin{table}[t]
\centering
\caption{Cycle count against length, per rule set. Five rule sets, six length
points each; $R^2 \approx 1.0000$ throughout. Source:
\texttt{bench/results/cycles.md}.}
\label{tab:cycles}
\tabfont
\begin{tabular}{@{}rrrrr@{}}
\toprule
Rules & \texttt{keyword} & \texttt{pattern} & Slope (cyc/char) & Intercept \\
\midrule
1 & 0 & 0 & 73.9      & 50{,}137 \\
2 & 1 & 0 & 105.3     & 65{,}942 \\
3 & 1 & 1 & 4{,}289.8 & 387{,}501 \\
4 & 2 & 1 & 4{,}048.8 & 405{,}785 \\
6 & 3 & 2 & 5{,}403.5 & 993{,}257 \\
\bottomrule
\end{tabular}
\end{table}

Three findings, in decreasing order of how much they change what may be claimed.

\emph{(1) With the rule set fixed, cycle count is exactly linear in length.}
Five rule sets, six length points each, $R^2 \approx 1.0000$ throughout
(Table~\ref{tab:cycles}).

\emph{(2) Scanning cost is dominated by the NFA, not by the rule count.} One
keyword rule costs about $+31.6$ cycles per character at the margin; one NFA
regular expression (an email pattern) costs about $+3{,}634$---roughly
$115\times$ more. Equivalently, a rule set containing a regular expression has a
per-character slope ($4{,}289.8$) $58\times$ that of a rule set with only string
rules ($73.9$).

\emph{(3) Unit prices are not additive: a unit price is a property of a rule
\emph{in context}, not of the rule.} Fitting the slopes as
$(\text{keyword},\text{pattern})$ unit prices gives $R^2 = 0.9169$ and a
\emph{negative} keyword unit price ($-477.2$ cycles per character). This is not
noise; two probes give the mechanism directly:

\begin{itemize}
  \item \textbf{Ordering probe} (same rules, same text, $L = 10^5$): declared
  order costs $541{,}342{,}069$ cycles against $570{,}736{,}989$ in reverse,
  a difference of $+5.43\%$, while \texttt{passed} is identical in both. \emph{The
  semantics of the decision do not depend on order; the cost does.}
  \item \textbf{Marginal probe} ($L = 10^5$, filling the one matrix cell the
  sweep could not reach---regular expressions with no keyword rule):
  $[\text{len},\text{pat}]$ over $[\text{len}]$ has a marginal unit price of
  $3{,}633.7$, whereas ``add one more keyword rule'' has two different answers:
  $+31.6$ alone, $+585.6$ in a policy that already contains a regular expression
  ($-153.7$ if inserted before it). The same action changes sign between
  $-154$ and $+586$.
\end{itemize}

The mechanism lies in zkVM memory allocation: every \texttt{ascii\_lower} and
every Pike VM match allocates, and allocator state depends on previous
allocations. We therefore quote \emph{per-rule-set slopes} (the table above,
directly citable) rather than any ``unit price $\times$ count'' expression. The
simple model $\text{cycles} \approx a\cdot|T|\cdot\text{rules} + b$ that this
work anticipated is refuted by the data ($R^2 = 0.872$, negative intercept), and
we report that too: \emph{a model refuted by its own data is itself an evaluation
result.}

\subsection{Matcher ablation and proof-path cost}
\label{sec:eval-ablation}

\emph{Construction.} Pattern \texttt{a+b} against input
\texttt{a}$^{\times n}$ containing no \texttt{b}, so the match is necessarily
absent and both implementations must scan to the end. The naive implementation
reinvents and restarts at every start position, giving $\Theta(n^2)$; the Pike VM
decides in a single linear pass (\texttt{bench/bench\_ablation.py}).
Table~\ref{tab:ablation} reports the cycles. Read per-character cost rather than
absolute figures. Doubling $n$ leaves the Pike VM's per-character cost
essentially unchanged ($-11.1\%$, $-6.4\%$) while the naive path's nearly doubles
($+96.4\%$, $+98.6\%$). At $n = 400$ the naive path is already $90\times$ slower,
and that ratio doubles again with each doubling of $n$. \emph{This is an attack
surface}: the input length for policy compliance is chosen externally---an
agent's reply, a concatenated session---so an adversary who needs only to write a
long reply drives a naive implementation's cost quadratically. The comparison is
meaningful only because the two paths decide identically, which
\texttt{tests/test\_ablation.py} pins down (Python and Rust
\texttt{match\_search} $\equiv$ \texttt{match\_search\_naive}). In the systems we
surveyed, we are not aware of ablation evidence specific to policy-compliance
matching.

\begin{table}[t]
\centering
\caption{Matcher ablation for pattern \texttt{a+b} on input \texttt{a}$^n$.
Cycles are zkVM execution counts. Source:
\texttt{bench/results/ablation.md}.}
\label{tab:ablation}
\tabfont
\begin{tabular}{@{}rrrrr@{}}
\toprule
$n$ & Naive cycles & Pike cycles & Ratio \\
\midrule
100 & 12{,}382{,}720  & 643{,}894   & $19.2\times$ \\
200 & 48{,}638{,}946  & 1{,}145{,}209 & $42.5\times$ \\
400 & 193{,}207{,}130 & 2{,}143{,}617 & $90.1\times$ \\
\bottomrule
\end{tabular}
\end{table}

\paragraph{Proof-path cost.} Re-running the same construction through the
compliance driver (\texttt{bench/results/dos\_proof.md}) reproduces the ratios
on the execution path---$643{,}906$ against $12{,}382{,}732$ cycles at $n=100$,
rising to $90.1\times$ at $n=400$---and then prices the proving path. Proving is
where the quadratic term meets the prover's fixed memory floor: the Pike path at
$n=100$ succeeds at $122.2$\,s and $11{,}695$\,MiB; the naive path at $n=25$
succeeds at $153.2$\,s and $12{,}137$\,MiB; at $n=50$ it is
\textbf{OOM/SIGKILL} after $226.5$\,s at $13{,}772$\,MiB peak. We report the OOM
as a result rather than a discarded run: it is the point at which the
adversarial quadratic term exhausts a $14.6$\,GiB machine, and it is why the
$n\ge100$ naive points and all $n=400$ proving points are left to a cloud
machine. The cycle counts differ from these execution-path figures by about one
part in $10^5$, because they come from a separate driver build; the ratios
agree. These proving runs predate the NFA fix discussed in the main paper and
carry that binary's digest; the fix changes the guest and the verification key,
so the figures are not mixed with any post-fix run.

\begin{table}[t]
\centering
\caption{Proof-path denial of service on a 24-core / 14.6\,GiB machine. The
\texttt{naive} row at $n=50$ is an observed out-of-memory kill, retained as a
result. Source: \texttt{bench/results/dos\_proof.md}.}
\label{tab:dos}
\tabfont
\setlength{\tabcolsep}{3pt}
\begin{tabular}{@{}lrrrl@{}}
\toprule
Point & Cycles & Prove (s) & RSS (MiB) & Result \\
\midrule
pike $n{=}100$ & 643{,}906 & 122.2 & 11{,}695.1 & \checkmark \\
naive $n{=}25$ & --- & 153.2 & 12{,}137.3 & \checkmark \\
naive $n{=}50$ & --- & 226.5 & 13{,}771.6 & OOM/SIGKILL \\
\bottomrule
\end{tabular}
\end{table}

\subsection{Composition: four negative controls}
\label{sec:eval-compose}

The leakage-profile theorem of the main paper prices each half; it does not say
that conjoining two valid halves yields a valid composite. We built the four
cells of Table~\ref{tab:compose} against the composite verifier's binding layer,
with the cryptographic validity of the two sub-proofs covered separately by the
end-to-end tests. Only cell~(a) satisfies the strong form, and only because
$x_U=\mathsf{derive}(T_U)$; the implementation claims only the weak form of the
composition proposition in the main paper and its E-1 remark. Cell~(c) is the
one that matters: the current verifier binds the inference half to
$\mathsf{derive}(T_U)$, not to the input $x_U$ the user submitted, so it accepts
a certificate that proves the weak form
$\exists x: x=\mathsf{derive}(T_U)\wedge T_U=M(x)$ while the strong form
$x=x_U$ remains unproved. This is proof debt E-1, not a test artefact; it is the
concrete reason composition alone does not close the gap.

\begin{table}[t]
\centering
\caption{Composition negative controls. \texttt{ok} is the verifier's binding
check; \texttt{satisfied} is the policy outcome. Source:
\texttt{bench/results/composition\_controls.md}.}
\label{tab:compose}
\tabfont
\setlength{\tabcolsep}{3pt}
\begin{tabular}{@{}lL{0.34\columnwidth}ccL{0.23\columnwidth}@{}}
\toprule
Cell & Construction & ok & sat. & Reading \\
\midrule
(a) & same $n$; $x_U=\mathsf{derive}(T_U)$; \texttt{delegated} empty &
\checkmark & \checkmark & true composition \\
(b) & the two halves use different nonces & $\times$ & $\times$ &
rejected by the four-way binding check \\
(c) & $x_U\neq\mathsf{derive}(T_U)$; inference binds only the derived input &
\checkmark & \checkmark & \textbf{weak form only}; $x=x_U$ not landed (E-1) \\
(d) & \texttt{delegated} non-empty & \checkmark & $\times$ &
true but unsatisfied; fail closed \\
\bottomrule
\end{tabular}
\end{table}

\subsection{Gateway signing: measured cost, blocked circuit}
\label{sec:eval-gateway}

Table~\ref{tab:gateway} prices the trust boundary the main paper states rather
than hides. Off-circuit Ed25519 verification costs $0.094$, $0.844$ and
$10.165$\,ms per receipt batch at $1$, $10$ and $100$ receipts, with a
per-receipt marginal near $0.08$--$0.10$\,ms; the session-end \texttt{ToolSeal}
is $0.074$\,ms to sign, $0.184$\,ms to verify and $345$\,B. A concatenated
2-of-3 threshold prototype signs in $0.173$\,ms, merges in $0.001$\,ms, verifies
in $0.508$\,ms and produces a $981$\,B certificate ($759$\,B when only the two
required signatures are carried). The threshold semantics are machine-checked:
any two distinct signatures pass, one signature is rejected, and two signatures
under the same key id do not count as two.

\begin{table}[t]
\centering
\caption{Gateway signing cost on a 24-core / 14.6\,GiB machine ($n=300$ per
row). Rows are measured medians. The in-circuit row is \emph{blocked}: the
circuit contains no Ed25519 verifier by design, and no cycle, time or memory
figure is reported. Source: \texttt{bench/results/gateway\_sig.md}.}
\label{tab:gateway}
\tabfont
\begin{tabular}{@{}lrrrr@{}}
\toprule
Configuration & Sign (ms) & Verify (ms) & Cert (B) & $n$ \\
\midrule
Ed25519, 1 receipt & 0.058 & 0.094 & 237.0 & 300 \\
Ed25519, 10 receipts & 0.415 & 0.844 & 288.3 & 300 \\
Ed25519, 100 receipts & 4.885 & 10.165 & 295.2 & 300 \\
\texttt{ToolSeal} (10 receipts) & 0.074 & 0.184 & 345 & 300 \\
1-of-1 & 0.069 & 0.173 & 648 & 300 \\
2-of-3 (3 signatures) & 0.173 & 0.508 & 981 & 300 \\
\midrule
In-circuit Ed25519 & \multicolumn{3}{c}{blocked: no guest, no number} & --- \\
\bottomrule
\end{tabular}
\end{table}

The in-circuit row is the honest blank. Verifying Ed25519 inside the zkVM is not
in the circuit---\texttt{circuits/} deliberately keeps the signature field out
of the circuit view---and building a guest to measure it exceeds this
experiment's file ownership, so cycles, proving time and peak memory are
\emph{null}, not zero. We therefore do not infer a circuit cost from the
off-circuit cost, and the trust root still contains the gateway key. What the
threshold prototype buys is not a cheaper verification but a weaker single point
of trust: one key compromised is no longer enough, at roughly three times the
signing work and certificate size, and with a new obligation to argue that the
gateways agree on the chain tail.

Putting Ed25519 inside the guest is a design option, not a measured one. SP1
6.7.0 has no Ed25519 precompile; the guest change would rotate the policy vkey
and add public values to the outcome, invalidating existing certificates, and we
did not build or time it. The only gateway-side variant measured is the 2-of-3
threshold prototype. It removes the single-key trust but not gateway collusion
or tail truncation. An authenticated gateway roster, HSM key custody and a
trusted time source remain open.

\subsection{Cross-system baselines: 49 cells, 46 populated}
\label{sec:eval-baselines}

Table~\ref{tab:baselines} places PoP against six alternatives across seven
dimensions: 49 cells, of which 46 are populated (measured or cited) and 3 are
unavailable. Every populated cell is tagged $[m]$ (measured in this repository
or this study) or $[c]$ (cited from the named paper). The three measured arms are
PoP, the signed-Merkle log and a simulated trusted server; zkAgent, NiyamAI,
Zaguk and PCO are cited. Zaguk and PCO did not yield comparable performance
numbers in our environment---the former's proceedings PDF was unreachable and
the latter reports no performance table---so those cells are em dashes, not
estimates. The seven dimensions impose different proof obligations; the table is
a statement of \emph{what each system obliges the verifier to trust}, not a
single ranking.

\begin{table*}[t]
\centering
\caption{Cross-system comparison. $[m]$ = measured here; $[c]$ = cited from the
source named in \texttt{bench/results/baselines.md}; em dash = not available,
not guessed. The trust root is per obligation, not one column: PoP's response
and policy proofs trust the vkey, while tool receipts additionally trust the
gateway key.}
\label{tab:baselines}
\tabfont
\setlength{\tabcolsep}{3pt}
\scriptsize
\begin{tabular}{@{}L{0.11\textwidth}L{0.13\textwidth}L{0.13\textwidth}L{0.12\textwidth}L{0.11\textwidth}L{0.10\textwidth}L{0.10\textwidth}L{0.08\textwidth}@{}}
\toprule
Dimension & PoP & signed-Merkle & TEE (sim) & zkAgent & NiyamAI & Zaguk & PCO \\
\midrule
Trust root & vkey + crypto $[m]$; tool-receipt gateway key & Ed25519 gateway $[m]$ &
server / TEE $[m]$ & zkVM + provider $[c]$ & SNARK + judge $[c]$ &
ezkl + commitment $[c]$ & authority + ledger $[c]$ \\
Latency & 40.5--100.3\,s core; 279--344\,s \texttt{gr\-oth16}; 89.84\,ms verify $[m]$ &
0.132\,ms sign; 2.131\,ms append; 0.322\,ms check $[m]$ &
0.139\,ms sign; 0.102\,ms verify $[m]$ &
194.4\,s / 0.42\,s $[c]$ & 1.7\,s / 51\,ms $[c]$ &
not comparable $[c]$ & --- \\
Size & 2.79\,MB core; 1{,}874\,B \texttt{gr\-oth16} $[m]$ & 648\,B + 761\,B $[m]$ &
648\,B $[m]$ & 96\,MiB (Shout) $[c]$ & 18.6\,KB $[c]$ & --- & --- \\
Response privacy & committed $[m]$ & none $[m]$ & none $[m]$ &
committed $[c]$ & not reported $[c]$ & yes $[c]$ & yes $[c]$ \\
Policy privacy & public values salted; proof artefact open $[m]$ & none $[m]$ & none $[m]$ &
n/a $[c]$ & contract commitment $[c]$ & yes $[c]$ & partial $[c]$ \\
Needs $T$ plaintext & no $[m]$ & yes $[m]$ & no $[m]$ & no $[c]$ &
unclear $[c]$ & no $[c]$ & no $[c]$ \\
Needs $\pi$ plaintext & no $[m]$ & yes $[m]$ & no $[m]$ & n/a $[c]$ &
unclear $[c]$ & no $[c]$ & partial $[c]$ \\
\bottomrule
\end{tabular}
\end{table*}

The asymmetry the table exposes is the one the signature baseline already
showed: the cheap arms require the verifier to re-run the judgment on plaintext,
while the private arms move that work into a proof or into a trusted party. PoP
is on the proof side, with a measured cost four to five orders of magnitude
above the signature arm and a trust root that, unlike the TEE arm, contains no
hardware.

\subsection{Issuance: detection, not prevention}
\label{sec:eval-issuance}

A provider that never issues a certificate is indistinguishable, from the
verifier's side, from a provider with nothing to certify; no cryptographic
layer forces issuance. What a transparency log changes is the \emph{detection}
rate, and only under an external obligation to append before delivery. In a
simulation with $N=2000$ deliveries, $m=100$ withheld violations and a sampling
rate $p=0.05$, an audit that samples the \emph{delivered} artifacts detects
$0.9948$ of the withheld set in exact arithmetic ($0.9940$ by Monte Carlo over
500 draws). An audit that samples the \emph{log} detects $0.0000$, because the
withheld deliveries have no log entry at all: the sampling frame is wrong, not
unlucky.

\begin{table}[t]
\centering
\caption{Issuance detection under two sampling frames ($N=2000$, $m=100$,
$p=0.05$, seed 20261101). Exact is $1-\binom{N-m}{s}/\binom{N}{s}$ with $s=100$;
Monte Carlo uses 500 draws. Source: \texttt{bench/results/issuance.md}.}
\label{tab:issuance}
\tabfont
\begin{tabular}{@{}lrr@{}}
\toprule
Sampling frame & Exact & Monte Carlo \\
\midrule
Delivered artifacts (correct frame) & 0.9948 & 0.9940 \\
Log entries (wrong frame) & 0.0000 & 0.0000 \\
\bottomrule
\end{tabular}
\end{table}

Two boundaries travel with these numbers. First, the result quantifies
\emph{detection}, never \emph{prevention}: whether a detected violation is
punished depends on penalties and incentives outside the model. Second, the
sufficiency direction assumes an append-before-delivery obligation, a log that
is append-only and publicly readable, inclusion proofs, and an external
mechanism against split views. The prototype implements append-only semantics,
inclusion proofs, optional tree-head signatures and checkpoints; it does
\emph{not} implement consistency proofs or gossip, so split-view resistance is
an external assumption, not a delivered property. Three premises travel with the
result: (F-1) the append-before-deliver obligation is external to cryptography;
(F-3) the log entry must bind the delivered $T$ (or a receipt), not only
$H(\mathrm{cert})$; and (F-4) the audit sample frame is delivered artifacts, not
log entries. The $1-(1-p)^m$ form assumes independent sampling and
detection-on-hit.

\subsection{Chat-template rerun and refusal proxy}
\label{sec:eval-k-chat}

A chat-template rerun of 20 prompts at $N=10$ returned $k=3.32=\log_2 10$ at
temperature $0.7$ and $k=0$ at temperature $0$; its refusal proxy rose from
$15\%$ to $35\%$ ($7/10$ unsafe prompts) while the response-violation rate fell
from $100\%$ to $90\%$. The $90\%$ is dominated by keyword-policy false
positives---$9/10$ safe responses were flagged---and is not a model harm rate;
XSTest labels are prompt-level and the refusal proxy is a prefix heuristic.

% =====================================================================
\section{Policy-Private Mode Details}
\label{app:private-policy}

\paragraph{Three parties.} ``Three'' means \emph{authority / operator / third
party}; an auditor is a named instance of the third party, not a fourth trust
anchor. Table~\ref{tab:wp4-parties} states the model.

\begin{table}[t]
\centering
\caption{The three-party model of policy-private mode.}
\label{tab:wp4-parties}
\tabfont
\setlength{\tabcolsep}{3pt}
\begin{tabular}{@{}L{0.16\columnwidth}L{0.14\columnwidth}L{0.62\columnwidth}@{}}
\toprule
Party & Honest? & Holds; wants \\
\midrule
Authority $R$ & \textbf{Yes}, by (P-PP) & Holds $\pi$ and the issuing record;
wants $\pi$ actually enforced, and wants to be able to check that it was. \\
Operator $O$ & \textbf{No} (the adversary) & Holds $\pi$, $T$, the proving key;
wants $\pi$ hidden from third parties while still being able to certify. \\
Third party $V$ & Yes & Holds only public values---\texttt{policy\_\allowbreak hash},
\texttt{circuit\_\allowbreak hash}, \texttt{response\_\allowbreak commitment},
\texttt{passed}, the proof; wants to confirm that \emph{some issued} policy
passed, and must \emph{not} learn $\pi$. \\
\emph{(Auditor $V^{*}$)} & Yes & As $V$, plus $\pi$ once $O$ discloses it;
wants ex-post accountability. \\
\bottomrule
\end{tabular}
\end{table}

Two boundaries must be stated with the model. First, secrecy of $\pi$ comes from
a \emph{data-flow fact}---$\pi$ never enters the public values---and not from
zero-knowledge; any rule that writes a fragment of $\pi$ into the public values
breaks secrecy on that fragment. Second, accountability for $V^{*}$ is
\emph{voluntary}: $V^{*}$ can check only when $O$ chooses to disclose $\pi$ and
the salt. The mode offers no forced disclosure; that would require the
release-gating family discussed in \texttt{docs/tdsc-plan.md}~\S1.6(b), which we
do not claim.

\paragraph{What each verifier may conclude.} Table~\ref{tab:wp4-capability} is
the contribution of this subsection: the cryptography is downstream of it, and
the table survives even if the cryptography only partially lands.

\begin{table*}[t]
\centering
\caption{What each verifier may and may not conclude in policy-private mode.}
\label{tab:wp4-capability}
\tabfont
\setlength{\tabcolsep}{4pt}
\begin{tabular}{@{}L{0.13\textwidth}L{0.30\textwidth}L{0.30\textwidth}L{0.19\textwidth}@{}}
\toprule
Verifier & \textbf{May} conclude & \textbf{May not} conclude & Rests on \\
\midrule
Third party $V$ & (i) the commitment binds $T$
(Definition~1 of the main paper); (ii) \emph{some} policy with commitment $h_\pi$
passed; (iii) $h_\pi$ appears in the issuing record & (iv) the \emph{content} of
$\pi$; (v) how many rules $\pi$ has, or of which classes; (vi) ``the operator
enforces the regulation \emph{I} care about''---only ``some issued regulation'' &
(P-PP); Assumption~\ref{as:Z}; the $k$ of Proposition~\ref{thm:t1} \\
End user (also receiver) & As $V$, and additionally: the proved $T$ is the
string it received, since its nonce is bound (Theorem~\ref{thm:t3}, requires
$f>0$)$^{\dagger}$ & Only the string it received; \emph{not} what was delivered
to another downstream agent & $f>0$ \\
Authority $R$ & Everything; it holds $\pi$, the salt and can recompute $h_\pi$ &
--- & --- \\
Auditor $V^{*}$ (after disclosure) & (vii) $\mathsf{verify\_commit}(\pi,s)=
\texttt{policy\_\allowbreak hash}$, i.e.\ what ran was the issued $\pi$;
(viii) the disclosed violation record is consistent & (ix) the undisclosed part
of $T$ & voluntary disclosure \\
On-chain verifier & (i)--(iii) only; it hashes and verifies signatures, and does
not interpret policy semantics & (iv)--(ix) & (P-PP) \\
\bottomrule
\end{tabular}
\end{table*}

The $V$ row is the reason the table exists. A third party does not obtain
``compliance''; it obtains ``some issued policy was satisfied''. Conflating the
two is the sentence a referee would use to reject the mode, and preventing that
sentence is what the table is for. $^{\dagger}$\,The end-user binding is
currently a certificate-layer field; the circuit still computes the
verifier-domain binding, so the end-to-end form needs a circuit change and a new
vkey.

% =====================================================================
\section{Related-Work Positioning}
\label{app:related}

Table~\ref{tab:related} records the slot each lineage occupies---which is also,
read the other way, what the main paper does not have to claim.
Table~\ref{tab:zkagent} places PoP against zkAgent: the two prove different
obligations, so the rows are a statement of orthogonality, not a race.

\begin{table}[t]
\centering
\caption{Positioning against the five prior-art lineages. The right-hand column
names what each lineage already occupies, and therefore what we do not claim.}
\label{tab:related}
\tabfont
\setlength{\tabcolsep}{3pt}
\begin{tabular}{@{}L{0.26\columnwidth}L{0.42\columnwidth}L{0.26\columnwidth}@{}}
\toprule
Lineage & Representative work & Slot already occupied \\
\midrule
Policy $\to$ circuit & zkGen \cite{zkgen}, Circuitree, Prezta & DSL-to-circuit
translation \\
Agent/LLM compliance & NiyamAI, Zaguk, PCO \cite{pco}, WITNESS & policy-level
private-mode kernel \\
Model computation & DeepProve, zkGPT, zkAgent \cite{zkagent} & inference
integrity \\
Attacks / negative results & Hollow-LLM \cite{hollowllm}, zk-attach
\cite{zkattach} & naming the gaps \\
Non-ZK industry baseline & \texttt{@brivora/verify}, VAOL & signed audit logs \\
\bottomrule
\end{tabular}
\end{table}

\begin{table}[t]
\centering
\caption{PoP against zkAgent. The two prove different obligations, so the rows
are a statement of orthogonality, not a race.}
\label{tab:zkagent}
\tabfont
\setlength{\tabcolsep}{3pt}
\begin{tabular}{@{}L{0.22\columnwidth}L{0.38\columnwidth}L{0.30\columnwidth}@{}}
\toprule
& \textbf{zkAgent} \cite{zkagent} & \textbf{PoP} \\
\midrule
Obligation & provider executed the declared model and tools (inference integrity)
& response/trace satisfies policy $\pi$ (compliance) \\
Depends on the model & yes (weights, arithmetisation) & \textbf{no} \\
Cost grows with & model size, sequence length & response length; rule-class
composition (per-rule-set slope), not rule count \\
Content privacy & none & committed, selective disclosure \\
Hardware reported & 32 cores / 512\,GB & 16 cores / 90\,GiB \\
Proving & 194.4\,s (Shout, $T{=}512$) & 40.5--100.3\,s core; 279--344\,s \texttt{gr\-oth16} \\
Proof size & 96\,MiB (Shout); 3.1\,MiB (LogUp) & 2.79\,MB core; 1{,}874\,B \texttt{gr\-oth16} \\
Verification & 0.42\,s (Shout); 0.038\,s (LogUp) & $\approx$0.090\,s \\
\bottomrule
\end{tabular}
\end{table}

We deliberately do not claim that PoP is faster, and the temptation to do so is
worth naming because the raw numbers invite it. The correct statement is
narrower, and we make only it: \emph{for the compliance obligation}, PoP
completes in a proving time of the same order as zkAgent's, with a smaller memory
footprint and a comparable proof size, without touching the model---and it adds
committed content privacy that zkAgent does not provide. Two caveats must travel
with those figures. The hardware is not matched (16 cores and a 90\,GiB cgroup
against 32 cores and 512\,GB), and PoP's ceiling is memory-bound rather than
cycle-bound, so its feasible region is a non-rectangular set of measured points
rather than a curve. The comparison is of orders of magnitude under unmatched
hardware, not a race. Source: \texttt{bench/comparison\_zkagent.md}, which
itself warns against reading the rows as ``PoP is faster.''

% =====================================================================
\section{Reproduction}
\label{app:repro}

The full walkthrough---environment, a single compliance proof, cross-checking
against the reference layer, private mode, certificate issuance, the one-command
demo, and third-party verification---is maintained in \texttt{docs/reproduce.md}
together with the exact commands and their expected output. This section records
only what a reader needs in order to judge the claims of the main paper's
evaluation.

\paragraph{Environment.} Ubuntu or WSL2 on x86-64, no GPU. The Rust/SP1 side
needs Rust stable, \texttt{cargo-prove}~v6.7.0, Go $\ge$\,1.24, and the usual
native toolchain; the reference layer is pure Python~3.10+ and runs on its own.
Memory, not time, is the binding constraint: the CPU prover carries a fixed
floor of approximately 10.5\,GiB on the cloud cgroup (10.15\,GiB under
default local prover options; \texttt{NUM\_CORE\_WORKERS=1},
\texttt{CORE\_BUFFER\_SIZE=1} lowers it to $\approx$9.3\,GiB) regardless of the
trace; memory, not time, sets the ceiling in the proof-cost table of the main
paper.

\paragraph{What reproduces without the prover.} The reference layer, the policy
compiler, and the whole test suite run in seconds and need no Rust. This is
enough to reproduce the non-ZK baseline of the main paper, to inspect the
certificate and inclusion proofs, and to check the cross-layer contract---that
the Python reference implementation and the zkVM guest consume the same
compiled policy bytes.

\paragraph{What needs the prover.} Every timing, size, and memory figure in the
cycle, ablation, denial-of-service, and proof-cost evaluations, and the
end-to-end result, requires a real proof. The measurements were taken on a
16-core CPU cgroup with 90\,GiB of RAM, which is what places the ceiling
reported there: the successes and the out-of-memory configurations are both
properties of that machine.

\paragraph{Where each number comes from.} Every measurement in the main paper's
evaluation is the content of a checked-in result file with a machine-readable
companion: \texttt{bench/results/} holds \texttt{cycles}, \texttt{ablation},
\texttt{scale-cloud}, \texttt{scale-6x2000}--\texttt{scale-6x5000},
\texttt{eu\_ai\_act\_proof}, \texttt{privacy\_surface\_real},
\texttt{k\_entropy\_largeN}, \texttt{k\_entropy\_crossmodel},
\texttt{policy\_corpus\_infer}, \texttt{agent\_e2e\_deepseek},
\texttt{baseline}, \texttt{verify}, \texttt{semantic}, \texttt{streaming},
\texttt{compose}, \texttt{dos\_proof}, \texttt{gateway\_sig},
\texttt{baselines}, \texttt{issuance}, and \texttt{prover\_knobs}, each as a
\texttt{.json} alongside its \texttt{.md}. The reproduction guide records the
commit, the hardware, and the prover binary digest for each run, so that a
reported figure can be traced to one specific execution rather than to a state
of the tree.

\begin{thebibliography}{99}

\bibitem{hollowllm}
C.~Gong, B.~Liu, and M.~Li, ``Hollow-LLM Attack: Computationally Trivial
Weights in Zero-Knowledge Verification of LLM Inference,'' in \emph{IEEE
Symposium on Security and Privacy (S\&P)}, 2026, pp.~4339--4355,
DOI 10.1109/SP63933.2026.00258; arXiv:2607.28884.

\bibitem{zkattach}
A.~S.~Gopalakrishna, ``Zero-Knowledge Predicate Proofs Between AI Agents: A
Measured, Cross-Protocol Gateway and the Source-Integrity Gap,''
arXiv:2608.30083, 2026. (The \texttt{zk-attach/v0} envelope is the system
contribution of this paper.)

\bibitem{zkgen}
J.~Lauinger, J.~Ernstberger, and S.~Steinhorst, ``zkGen: Policy-to-Circuit
Transpiler,'' in \emph{IEEE International Conference on Blockchain and
Cryptocurrency (ICBC)}, 2024, pp.~620--624,
DOI 10.1109/ICBC59979.2024.10634440.

\bibitem{pco}
H.~Avni, S.~Dolev, A.~Yagudaev, and M.~Yung, ``Super-intelligence Survival
Guide: Verification via Proof-Carrying Output,'' IACR Cryptology ePrint
Archive, Report 2026/994, 2026.

\bibitem{zkagent}
L.~Wang, H.~Lou, C.~Li, Y.~Yu, and Y.~Hu, ``zkAgent: Verifiable LLM Agent
Execution via One-Shot Transcript Proofs,'' IACR Cryptology ePrint Archive,
Report 2026/199, 2026.

\bibitem{hiddencred}
J.~E.~Holt, R.~W.~Bradshaw, K.~E.~Seamons, and H.~Orman, ``Hidden
Credentials,'' in \emph{Proc.\ 2003 ACM Workshop on Privacy in the Electronic
Society (WPES)}, 2003, pp.~1--8, DOI 10.1145/1005140.1005142.

\bibitem{handshake}
D.~Balfanz, G.~Durfee, N.~Shankar, D.~K.~Smetters, J.~Staddon, and H.-C.~Wong,
``Secret Handshakes from Pairing-Based Key Agreements,'' in \emph{Proc.\ IEEE
Symposium on Security and Privacy (S\&P)}, 2003, pp.~180--196,
DOI 10.1109/SECPRI.2003.1199336.

\bibitem{xacmlpriv}
C.~A.~Ardagna, S.~De~Capitani~di~Vimercati, G.~Neven, S.~Paraboschi,
F.-S.~Preiss, P.~Samarati, and M.~Verdicchio, ``Enabling Privacy-Preserving
Credential-Based Access Control with XACML and SAML,'' in \emph{Proc.\ 10th
IEEE International Conference on Computer and Information Technology (CIT)},
2010, pp.~1090--1095, DOI 10.1109/CIT.2010.199.

\bibitem{sw05}
A.~Sahai and B.~Waters, ``Fuzzy Identity-Based Encryption,'' in \emph{Advances
in Cryptology -- EUROCRYPT 2005}, LNCS 3494, Springer, 2005, pp.~457--473,
DOI 10.1007/11426639\_27.

\bibitem{ksw08}
J.~Katz, A.~Sahai, and B.~Waters, ``Predicate Encryption Supporting
Disjunctions, Polynomial Equations, and Inner Products,'' in \emph{Advances in
Cryptology -- EUROCRYPT 2008}, Springer, 2008, pp.~146--162,
DOI 10.1007/978-3-540-78967-3\_9.

\bibitem{vdr}
H.~Chabanne, R.~Hugel, and J.~Keuffer, ``Verifiable Document Redacting,'' in
\emph{Computer Security -- ESORICS 2017}, LNCS 10492, Springer, 2017,
pp.~334--351, DOI 10.1007/978-3-319-66402-6\_20.

\end{thebibliography}

\end{document}
"
%
%  This document carries the material removed from the 12--14 page main
%  paper to meet the TDSC length limit: full proofs, the detailed
%  evaluation tables, policy-private-mode detail, related-work tables,
%  and reproduction instructions. It is self-contained: it restates the
%  results it proves and can be read without the main PDF.
% =====================================================================
\documentclass[10pt,journal,compsoc]{IEEEtran}

\usepackage{amsmath,amsfonts,amssymb,amsthm}
\usepackage{algorithmic}
\usepackage{algorithm}
\usepackage{array}
\usepackage[caption=false,font=normalsize,labelfont=sf,textfont=sf]{subfig}
\usepackage{textcomp}
\usepackage{stfloats}
\usepackage{url}
\usepackage{verbatim}
\usepackage{graphicx}
\usepackage{cite}
\usepackage{booktabs}
\usepackage{xcolor}
\usepackage{microtype}
\usepackage{hyperref}

\hyphenation{Proof-of-Policy zero-knowl-edge re-exec-ution pol-icy
com-pli-ance bind-ing veri-fi-a-bil-ity cer-ti-fi-cate non-zero-knowl-edge}

\setlength{\emergencystretch}{2em}

\newcolumntype{L}[1]{>{\raggedright\arraybackslash}p{#1}}
\newcommand{\tabfont}{\footnotesize}

\newtheorem{theorem}{Theorem}
\newtheorem{lemma}[theorem]{Lemma}
\newtheorem{proposition}[theorem]{Proposition}
\newtheorem{corollary}[theorem]{Corollary}
\newtheorem{definition}[theorem]{Definition}
\newtheorem{assumption}[theorem]{Assumption}
\newtheorem{remark}[theorem]{Remark}

% Work markers are no-ops in the supplementary material.
\newcommand{\need}[1]{}
\newcommand{\fixme}[1]{}
\newcommand{\wpmark}[1]{}

% Supplementary numbering: S1, S2, ... and Table S1, S2, ...
\renewcommand{\thesection}{S\arabic{section}}
\renewcommand{\thetable}{S\arabic{table}}

\begin{document}

\title{Supplementary Material for\\``Proof-of-Policy: Zero-Knowledge Compliance
Certificates for LLM Agent Outputs''}
\author{Author Name(s)}
\maketitle

\begin{abstract}
This supplement contains the full proofs, the detailed evaluation tables, the
policy-private-mode tables, the related-work positioning, and the reproduction
instructions omitted from the main paper for length. Section and table numbers
are prefixed with ``S''. The main paper states every result; this document
supplies the arguments and the per-configuration numbers behind them.
\end{abstract}

% =====================================================================
\section{Full Proofs}
\label{app:proofs}

This section restates the results proved in the main paper and gives complete
proofs. We use the main paper's notation; for convenience, the two objects the
proofs depend on are repeated here.

\begin{definition}[Guessing entropy; restated from the main paper]
\label{def:k}
For a secret distribution $\mathcal{D}$ and adversary side information $Z$, let
\begin{multline*}
\mathsf{Guess}(X \mid Z) \;\triangleq\; \max_{\mathcal{A}} \Pr[\mathcal{A}(Z) = X]\\
\;=\; \mathbb{E}_z\!\left[\max_x \Pr[X = x \mid Z = z]\right],
\end{multline*}
\[
k_X \;\triangleq\; -\log_2 \mathsf{Guess}(X \mid Z),
\]
the maximum ranging over all time-unbounded adversaries.
\end{definition}

\begin{assumption}[Adversary side information; restated from the main paper]
\label{as:Z}
The adversary knows the prompt, the model (possibly its weights), and the
decoding parameters.
\end{assumption}

\begin{proposition}[Conditional trilemma; restated from the main paper]
\label{thm:t1}
Fix $\lambda,\varepsilon>0$ and $q\ge 1$. Let $X$ be a secret with guessing
entropy $k_X<\log q+\log(1/\varepsilon)$ whose binding is \emph{guess-tight}:
$\beta_\rho(X)=H(\mathsf{DOM}\|\rho\|X)$ for a random oracle $H$ and public
randomness $\rho$ independent of $X$ and $Z$. Then:
\begin{enumerate}
\item[(1)] \emph{Public binding.} For any adversary that evaluates $\beta$ at
most $q$ times,
\[
\Pr\!\left[\mathcal{A}^\beta(Z,\beta(X))=X\right]
\le q\cdot\mathsf{Guess}(X\mid Z)+\mathrm{negl}(\lambda)
= q\cdot 2^{-k_X}+\mathrm{negl}(\lambda).
\]
Hiding with advantage $\varepsilon$ is \emph{certified} by this bound when
$k_X\ge\log q+\log(1/\varepsilon)$; for a skewed distribution this is a
sufficient condition, not a matching attack. The integer $q$ counts evaluations
of $\beta$, not time: the adversary may be computationally unbounded.
\item[(2)] \emph{Keyed binding.} For $\beta_s(X)=H(s,X)$ with $s$ held only by a
designated party and absent from the public values,
$\Pr[\mathcal{A}(Z,b)=X]\le\mathsf{Guess}(X\mid Z)+\mathrm{negl}(\kappa)$.
What the keyed form buys is hiding when
$k_X<\log q+\log(1/\varepsilon)$; its price is PV-2.
\item[(3)] \emph{Conditional trilemma.} Consider schemes whose binding is
publicly recomputable (PV-2), whose hiding is exact recovery, and whose binding
is guess-tight. Within the scheme class defined above, when
$k_X<\log q+\log(1/\varepsilon)$ such a scheme cannot satisfy binding, hiding
and PV-2 simultaneously.
\end{enumerate}
\end{proposition}

\begin{proof}[Proof of Proposition~\ref{thm:t1}]
\emph{(1).} Fix $z$ and let the adversary's query set be
$Q = \{q_1, \dots, q_q\}$. Decompose by first hit,
\[
\Pr[\exists i : q_i = X] \;=\; \sum_i \Pr\!\left[q_i = X \wedge \neg E_{<i}\right],
\]
where $E_{<i} \triangleq \{\exists j < i : q_j = X\}$.
Conditioned on $\neg E_{<i}$, every oracle answer $a_1,\dots,a_{i-1}$ already
revealed is uniform and independent of $X$ (no queried point equals $X$, and $H$
is a fresh random function), and $\beta(X)$ itself is uniform and independent of
$X$. Hence $q_i$---a deterministic function of $(z, \beta(X), a_{<i})$---is
\emph{independent of $X$}, so
\begin{multline*}
\Pr[q_i = X \mid Z = z]
= \sum_x \Pr[q_i = x]\,\Pr[X = x \mid Z = z] \\
\le \max_x \Pr[X = x \mid Z = z] = 2^{-k_z},
\end{multline*}
giving $\Pr[\exists i : q_i = X \mid z] \le q\,2^{-k_z}$. Taking the expectation
over $z$ and using $\mathbb{E}_z[2^{-k_z}] = \mathsf{Guess}(X \mid Z) = 2^{-k_X}$
from Definition~\ref{def:k} yields $q \cdot 2^{-k_X}$. An adversary that outputs
an unqueried value wins with probability at most $2^{-k_z} \le q\,2^{-k_z}$,
which the same bound absorbs. No $q\cdot2^{-\lambda}$ term is needed in the
exact-recovery game: a query that lands on a different $x$ with
$\beta(x)=y$ is a second preimage of $H$, not a win. The explicit
$q\cdot2^{-\lambda}$ is kept only when the win condition is consistency with the
published $b$ rather than exact recovery, or when $q$ is unbounded.

\emph{(2).} Without an oracle for $\beta_s$, the value $b=H(s,X)$ is uniform and
independent of $X$ in the random-oracle model for an adversary that does not know
$s$, so the best strategy is the prior mode and succeeds with probability
$\mathsf{Guess}(X\mid Z)$. An attempted enumeration must also guess $s$, at cost
$2^{-\kappa}$ per query. The bound is therefore
$\mathsf{Guess}(X\mid Z)+\mathrm{negl}(\kappa)$, with no $q$-fold amplification.
Part (2) buys hiding only below the public-form threshold
$\log q+\log(1/\varepsilon)$; when $k_X=0$ the prior mode already recovers $X$
and no binding helps.

\emph{(3).} Public recomputability gives (1). If (1) exceeds $\varepsilon$, the
only way to avoid it is a binding that a verifier cannot recompute, i.e.\ a keyed
binding, which by part (2) forfeits PV-2. The third leg---hiding plus PV-2
without binding---is excluded from the class of compliance schemes because a
certificate whose binding does not pin $X$ certifies nothing about the delivered
artifact.
\end{proof}

\noindent\emph{Preimage-class bound.} For an arbitrary deterministic public
$\beta$, conditioning on $\beta(X)$ gives
$\Pr[\mathcal{A}^\beta(Z,\beta(X))=X]\le
\mathbb{E}_y[\max_x\Pr[X=x\mid Z,\beta(X)=y]]$, with no hypothesis on $\beta$.
For $\beta=\mathrm{id}$ or an injective $\mathsf{encode}$ this equals $1$; for a
lossy encoding it equals the inverse of the preimage-class size.

\begin{theorem}[D-privacy from three conditions; restated from the main paper]
\label{thm:dpriv}
The scheme is $D$-private if (C1) there is a PPT simulator $S_{\mathrm{rest}}$
such that, for every admissible $(\pi,T,n)$, the public values other than the
commitment satisfy
$(h_\pi,S_{\mathrm{rest}}(pp,h_\pi,D(\pi),\mathrm{pub}_T,n))\approx_c
L_{\mathrm{pub}}(\pi,T,n)$; (C2) the proof system is witness-hiding, i.e.\
$L_{\mathrm{proof}}$ has a simulator that takes only the public statement;
and (C3) the policy commitment is hiding for all $D$-equal policies,
$(pp,h_{\pi_0})\approx_c(pp,h_{\pi_1})$. The advantage is
$O(\varepsilon_{\mathrm{pub}}+\varepsilon_{\mathrm{wh}}+\varepsilon_{\mathrm{com}})$.
\end{theorem}

\begin{proof}[Proof of Theorem~\ref{thm:dpriv}]
Fix an admissible tuple $(\pi_0,\pi_1,T,n)$ with $D(\pi_0)=D(\pi_1)$ and
$R(\pi_0,T)=R(\pi_1,T)$. Write $h_i=h_{\pi_i}$,
$R_i=L_{\mathrm{rest}}(\pi_i,T,n)$, $P_i=L_{\mathrm{proof}}(\pi_i,T,n)$ and
$\mathsf{Stmt}_i=(h_i,R_i)$. Consider the hybrids
\begin{align*}
H_0&=(h_0,R_0,P_0),\\
H_1&=(h_0,R_0,S_{\mathrm{proof}}(\mathsf{Stmt}_0)),\\
H_2&=(h_0,S_{\mathrm{rest}}(h_0,D,\mathrm{pub}_T,n),S_{\mathrm{proof}}(\mathsf{Stmt}_0)),\\
H_3&=(h_1,S_{\mathrm{rest}}(h_1,D,\mathrm{pub}_T,n),S_{\mathrm{proof}}(\mathsf{Stmt}_1)),\\
H_4&=(h_1,R_1,S_{\mathrm{proof}}(\mathsf{Stmt}_1)),\\
H_5&=(h_1,R_1,P_1).
\end{align*}
$H_0\approx H_1$ and $H_4\approx H_5$ cost $\varepsilon_{\mathrm{wh}}$ each by
C2; $H_1\approx H_2$ and $H_3\approx H_4$ cost
$\varepsilon_{\mathrm{pub}}$ each by C1. For $H_2\approx H_3$, split into
$H_{2a}=(h_1,S_{\mathrm{rest}}(h_1,D,\mathrm{pub}_T,n),S_{\mathrm{proof}}(\mathsf{Stmt}_0))$
and $H_{2b}=H_3$: $H_2\approx H_{2a}$ costs $\varepsilon_{\mathrm{com}}$ by C3,
and $H_{2a}\approx H_{2b}$ follows because
$\mathsf{Stmt}_0\approx_c\mathsf{Stmt}_1$ (C1 and C3 by triangle inequality) and
$S_{\mathrm{proof}}$ is a PPT function of the statement. The triangle inequality
gives $4\varepsilon_{\mathrm{pub}}+2\varepsilon_{\mathrm{wh}}+2\varepsilon_{\mathrm{com}}$,
which is $O(\varepsilon_{\mathrm{pub}}+\varepsilon_{\mathrm{wh}}+\varepsilon_{\mathrm{com}})$.
The game fixes one $T$ because two different artifacts are distinguished by
$B(n,\cdot)$ itself, which is response-side leakage and not a policy advantage.
\end{proof}

\begin{theorem}[Leakage-profile decomposition; restated from the main paper]
\label{thm:leak}
Let $V=(L_{\mathrm{pub}},L_{\mathrm{proof}},L_{\mathrm{side}})$ be the verifier's
view, where $L_{\mathrm{pub}}$ is the certificate's public values,
$L_{\mathrm{proof}}$ is what the proof artefact reveals about the witness
$w=(\pi,T,\mathrm{trace})$, and $L_{\mathrm{side}}$ is the adversary's prior
information $Z$. For any $q\ge 1$ and any adversary that evaluates $\beta$ at
most $q$ times,
\[
\Pr\!\left[\mathcal{A}^\beta(V)=X\right]
\;\le\; q\cdot\mathsf{Guess}\!\left(X\mid L_{\mathrm{pub}},L_{\mathrm{proof}},L_{\mathrm{side}}\right)
+\mathrm{negl}(\lambda),
\]
with the final output charged as one evaluation.
\end{theorem}

\noindent\emph{Proof of Theorem~\ref{thm:leak}.} The proof is the first-hit
decomposition of Proposition~\ref{thm:t1}(1) with $(Z,\beta(X))$ replaced by the
verifier view $V$. The step that changes is the uniformity of answers to queries
that have not hit $X$: it holds conditionally on $L_{\mathrm{pub}}$ and fails when
$L_{\mathrm{proof}}$ is a function of the witness. A witness-hiding tier restores
$L_{\mathrm{proof}}=\bot$ and hence the public-value bound; a \texttt{core}
artefact does not, and the bound can reach $1$.

\begin{theorem}[Delivery binding; restated from the main paper]
\label{thm:t3}
Fix a policy $\pi$. A delivery adversary $\mathcal{B}$ may, \emph{before}
seeing the receiver's challenge, precompute a set $S$ of accepted transcripts,
where each transcript covers at most $c$ challenges. The receiver then
contributes $f$ fresh, uniformly random bits $r$ that enter the proven
statement. Under the certificate's soundness (D1) and challenge binding
(D2/B-bind), the divergence advantage satisfies
\[
\mathsf{Adv}_{\mathrm{del}}(\mathcal{B}) \;\le\; |S| \cdot c \cdot 2^{-f}
\;+\; \mathrm{negl}(\lambda).
\]
The replay-only probability---acceptance without requiring divergence---is at
most $|S|\cdot c\cdot 2^{-f}$. If the receiver contributes no fresh randomness
($f = 0$; e.g.\ an empty nonce), a precomputed transcript is accepted with
probability $1$ and \emph{does not decay in any parameter of the system}.
\end{theorem}

\begin{proof}[Proof of Theorem~\ref{thm:t3}]
The receiver's freshness $r$ enters the proven statement through
$b = H(\mathsf{dom} \,\|\, r \,\|\, T)$. For a certificate in $S$ to be accepted
on $r$, the value $r$ must fall inside that certificate's challenge set
$A(\mathrm{cert})$, with $|A(\mathrm{cert})|\le c$; a union bound over $S$ gives
$\sum_{\mathrm{cert}\in S}|A(\mathrm{cert})|/2^f\le |S|\cdot c\cdot 2^{-f}$.
A certificate minted \emph{after} $r$ is seen can diverge only by exhibiting two
distinct preimages of the same $b$ under $B(r,\cdot)$, which is exactly a
challenge-binding failure and costs $\mathrm{negl}(\lambda)$. When $f = 0$ every
element of $S$ is a permanently valid replay, so the replay probability is $1$
and does not decay in any parameter of the system. The $f = 0$ branch is not
hypothetical: \texttt{circuits/\allowbreak types/\allowbreak src/\allowbreak
lib.rs} (lines 950--995) documents an empty nonce as a legal input that
``provides no replay protection.'' The $f=0$ separation holds under stateless
public verification with no external mechanism (no revocation table, trusted
clock, or hardware). The $c$ factor is tight when the challenge sets of the
transcripts in $S$ are disjoint and of size $c$.
\end{proof}

\begin{proposition}[Quadratic cost of the naive matcher]
\label{prop:naive}
For the pattern \texttt{a+b} and input \texttt{a}$^n$, the naive backtracking
matcher performs
$\sum_{i=1}^{n}(n-i+1)=n(n+1)/2=\Theta(n^{2})$ NFA steps, whereas the Pike VM
performs $\Theta(n)$. The adversarial input family is admissible because the
response bytes are chosen by the party being checked. This bounds the cost of
\emph{that} baseline, not of policy-compliance matching: the Pike VM is linear
on the same family, and no lower bound for the problem is claimed.
\end{proposition}
\begin{proof}
No input position contains \texttt{b}, so no start position can fail or accept
before the end of the string. The start at $i$ therefore takes exactly
$n-i+1$ steps, and the sum is $n(n+1)/2$. The Pike VM keeps one state per NFA
position and scans once.
\end{proof}

\begin{proposition}[NFA contract]
\label{prop:nfa}
For the supported regular-expression subset $P$---ASCII, ordered and disjoint
character-class ranges, Unicode scalar values, and the same compiled policy on
both sides---the Thompson construction preserves the language, the Pike VM
implements unanchored substring search, and the Python reference and Rust
circuit matchers agree on the same compiled NFA specification. Empty character
classes are rejected at parse time (fail-closed), so ``non-empty character
class'' is a parser contract rather than an extra theorem assumption.
\end{proposition}
\begin{proof}[Proof sketch]
Language preservation is the standard Thompson construction over the restricted
syntax. The Pike VM simulates all NFA states at one input position and advances
once per character, so it decides unanchored search in one pass. The two
matchers consume the same \texttt{ConstraintSpec} bytes and implement the same
transition relation; agreement is checked by an algorithm-level bisimulation
plus differential fuzzing against Python's \texttt{re}, not by line-by-line
formal verification.
\end{proof}

% =====================================================================
\section{Evaluation Details}
\label{app:eval-details}

This section carries the per-configuration tables and the detailed readings
omitted from the main paper. Every number is the content of the checked-in
result file named in the caption, as in the main paper.

\subsection{Cycle sweep: length versus rule count}
\label{sec:eval-cycles}

Response length was pushed from 200 to 100{,}000 characters and the rule count
from 1 to 6, giving 40 measurement points (\texttt{bench/bench\_cycles.py}; the
full matrix is in \texttt{bench/results/cycles.md}). The response text comes from
an evaluation artefact the repository's own end-to-end demo produced (372
characters, three segments deduplicated by content); longer points are tiled from
it. What is measured is therefore the \emph{curve} of scanning cost, not ``a real
100{,}000-character document''. Every point was first confirmed to match no rule:
a matching rule short-circuits, and the cost of a short-circuit says nothing.

\begin{table}[t]
\centering
\caption{Cycle count against length, per rule set. Five rule sets, six length
points each; $R^2 \approx 1.0000$ throughout. Source:
\texttt{bench/results/cycles.md}.}
\label{tab:cycles}
\tabfont
\begin{tabular}{@{}rrrrr@{}}
\toprule
Rules & \texttt{keyword} & \texttt{pattern} & Slope (cyc/char) & Intercept \\
\midrule
1 & 0 & 0 & 73.9      & 50{,}137 \\
2 & 1 & 0 & 105.3     & 65{,}942 \\
3 & 1 & 1 & 4{,}289.8 & 387{,}501 \\
4 & 2 & 1 & 4{,}048.8 & 405{,}785 \\
6 & 3 & 2 & 5{,}403.5 & 993{,}257 \\
\bottomrule
\end{tabular}
\end{table}

Three findings, in decreasing order of how much they change what may be claimed.

\emph{(1) With the rule set fixed, cycle count is exactly linear in length.}
Five rule sets, six length points each, $R^2 \approx 1.0000$ throughout
(Table~\ref{tab:cycles}).

\emph{(2) Scanning cost is dominated by the NFA, not by the rule count.} One
keyword rule costs about $+31.6$ cycles per character at the margin; one NFA
regular expression (an email pattern) costs about $+3{,}634$---roughly
$115\times$ more. Equivalently, a rule set containing a regular expression has a
per-character slope ($4{,}289.8$) $58\times$ that of a rule set with only string
rules ($73.9$).

\emph{(3) Unit prices are not additive: a unit price is a property of a rule
\emph{in context}, not of the rule.} Fitting the slopes as
$(\text{keyword},\text{pattern})$ unit prices gives $R^2 = 0.9169$ and a
\emph{negative} keyword unit price ($-477.2$ cycles per character). This is not
noise; two probes give the mechanism directly:

\begin{itemize}
  \item \textbf{Ordering probe} (same rules, same text, $L = 10^5$): declared
  order costs $541{,}342{,}069$ cycles against $570{,}736{,}989$ in reverse,
  a difference of $+5.43\%$, while \texttt{passed} is identical in both. \emph{The
  semantics of the decision do not depend on order; the cost does.}
  \item \textbf{Marginal probe} ($L = 10^5$, filling the one matrix cell the
  sweep could not reach---regular expressions with no keyword rule):
  $[\text{len},\text{pat}]$ over $[\text{len}]$ has a marginal unit price of
  $3{,}633.7$, whereas ``add one more keyword rule'' has two different answers:
  $+31.6$ alone, $+585.6$ in a policy that already contains a regular expression
  ($-153.7$ if inserted before it). The same action changes sign between
  $-154$ and $+586$.
\end{itemize}

The mechanism lies in zkVM memory allocation: every \texttt{ascii\_lower} and
every Pike VM match allocates, and allocator state depends on previous
allocations. We therefore quote \emph{per-rule-set slopes} (the table above,
directly citable) rather than any ``unit price $\times$ count'' expression. The
simple model $\text{cycles} \approx a\cdot|T|\cdot\text{rules} + b$ that this
work anticipated is refuted by the data ($R^2 = 0.872$, negative intercept), and
we report that too: \emph{a model refuted by its own data is itself an evaluation
result.}

\subsection{Matcher ablation and proof-path cost}
\label{sec:eval-ablation}

\emph{Construction.} Pattern \texttt{a+b} against input
\texttt{a}$^{\times n}$ containing no \texttt{b}, so the match is necessarily
absent and both implementations must scan to the end. The naive implementation
reinvents and restarts at every start position, giving $\Theta(n^2)$; the Pike VM
decides in a single linear pass (\texttt{bench/bench\_ablation.py}).
Table~\ref{tab:ablation} reports the cycles. Read per-character cost rather than
absolute figures. Doubling $n$ leaves the Pike VM's per-character cost
essentially unchanged ($-11.1\%$, $-6.4\%$) while the naive path's nearly doubles
($+96.4\%$, $+98.6\%$). At $n = 400$ the naive path is already $90\times$ slower,
and that ratio doubles again with each doubling of $n$. \emph{This is an attack
surface}: the input length for policy compliance is chosen externally---an
agent's reply, a concatenated session---so an adversary who needs only to write a
long reply drives a naive implementation's cost quadratically. The comparison is
meaningful only because the two paths decide identically, which
\texttt{tests/test\_ablation.py} pins down (Python and Rust
\texttt{match\_search} $\equiv$ \texttt{match\_search\_naive}). In the systems we
surveyed, we are not aware of ablation evidence specific to policy-compliance
matching.

\begin{table}[t]
\centering
\caption{Matcher ablation for pattern \texttt{a+b} on input \texttt{a}$^n$.
Cycles are zkVM execution counts. Source:
\texttt{bench/results/ablation.md}.}
\label{tab:ablation}
\tabfont
\begin{tabular}{@{}rrrrr@{}}
\toprule
$n$ & Naive cycles & Pike cycles & Ratio \\
\midrule
100 & 12{,}382{,}720  & 643{,}894   & $19.2\times$ \\
200 & 48{,}638{,}946  & 1{,}145{,}209 & $42.5\times$ \\
400 & 193{,}207{,}130 & 2{,}143{,}617 & $90.1\times$ \\
\bottomrule
\end{tabular}
\end{table}

\paragraph{Proof-path cost.} Re-running the same construction through the
compliance driver (\texttt{bench/results/dos\_proof.md}) reproduces the ratios
on the execution path---$643{,}906$ against $12{,}382{,}732$ cycles at $n=100$,
rising to $90.1\times$ at $n=400$---and then prices the proving path. Proving is
where the quadratic term meets the prover's fixed memory floor: the Pike path at
$n=100$ succeeds at $122.2$\,s and $11{,}695$\,MiB; the naive path at $n=25$
succeeds at $153.2$\,s and $12{,}137$\,MiB; at $n=50$ it is
\textbf{OOM/SIGKILL} after $226.5$\,s at $13{,}772$\,MiB peak. We report the OOM
as a result rather than a discarded run: it is the point at which the
adversarial quadratic term exhausts a $14.6$\,GiB machine, and it is why the
$n\ge100$ naive points and all $n=400$ proving points are left to a cloud
machine. The cycle counts differ from these execution-path figures by about one
part in $10^5$, because they come from a separate driver build; the ratios
agree. These proving runs predate the NFA fix discussed in the main paper and
carry that binary's digest; the fix changes the guest and the verification key,
so the figures are not mixed with any post-fix run.

\begin{table}[t]
\centering
\caption{Proof-path denial of service on a 24-core / 14.6\,GiB machine. The
\texttt{naive} row at $n=50$ is an observed out-of-memory kill, retained as a
result. Source: \texttt{bench/results/dos\_proof.md}.}
\label{tab:dos}
\tabfont
\setlength{\tabcolsep}{3pt}
\begin{tabular}{@{}lrrrl@{}}
\toprule
Point & Cycles & Prove (s) & RSS (MiB) & Result \\
\midrule
pike $n{=}100$ & 643{,}906 & 122.2 & 11{,}695.1 & \checkmark \\
naive $n{=}25$ & --- & 153.2 & 12{,}137.3 & \checkmark \\
naive $n{=}50$ & --- & 226.5 & 13{,}771.6 & OOM/SIGKILL \\
\bottomrule
\end{tabular}
\end{table}

\subsection{Composition: four negative controls}
\label{sec:eval-compose}

The leakage-profile theorem of the main paper prices each half; it does not say
that conjoining two valid halves yields a valid composite. We built the four
cells of Table~\ref{tab:compose} against the composite verifier's binding layer,
with the cryptographic validity of the two sub-proofs covered separately by the
end-to-end tests. Only cell~(a) satisfies the strong form, and only because
$x_U=\mathsf{derive}(T_U)$; the implementation claims only the weak form of the
composition proposition in the main paper and its E-1 remark. Cell~(c) is the
one that matters: the current verifier binds the inference half to
$\mathsf{derive}(T_U)$, not to the input $x_U$ the user submitted, so it accepts
a certificate that proves the weak form
$\exists x: x=\mathsf{derive}(T_U)\wedge T_U=M(x)$ while the strong form
$x=x_U$ remains unproved. This is proof debt E-1, not a test artefact; it is the
concrete reason composition alone does not close the gap.

\begin{table}[t]
\centering
\caption{Composition negative controls. \texttt{ok} is the verifier's binding
check; \texttt{satisfied} is the policy outcome. Source:
\texttt{bench/results/composition\_controls.md}.}
\label{tab:compose}
\tabfont
\setlength{\tabcolsep}{3pt}
\begin{tabular}{@{}lL{0.34\columnwidth}ccL{0.23\columnwidth}@{}}
\toprule
Cell & Construction & ok & sat. & Reading \\
\midrule
(a) & same $n$; $x_U=\mathsf{derive}(T_U)$; \texttt{delegated} empty &
\checkmark & \checkmark & true composition \\
(b) & the two halves use different nonces & $\times$ & $\times$ &
rejected by the four-way binding check \\
(c) & $x_U\neq\mathsf{derive}(T_U)$; inference binds only the derived input &
\checkmark & \checkmark & \textbf{weak form only}; $x=x_U$ not landed (E-1) \\
(d) & \texttt{delegated} non-empty & \checkmark & $\times$ &
true but unsatisfied; fail closed \\
\bottomrule
\end{tabular}
\end{table}

\subsection{Gateway signing: measured cost, blocked circuit}
\label{sec:eval-gateway}

Table~\ref{tab:gateway} prices the trust boundary the main paper states rather
than hides. Off-circuit Ed25519 verification costs $0.094$, $0.844$ and
$10.165$\,ms per receipt batch at $1$, $10$ and $100$ receipts, with a
per-receipt marginal near $0.08$--$0.10$\,ms; the session-end \texttt{ToolSeal}
is $0.074$\,ms to sign, $0.184$\,ms to verify and $345$\,B. A concatenated
2-of-3 threshold prototype signs in $0.173$\,ms, merges in $0.001$\,ms, verifies
in $0.508$\,ms and produces a $981$\,B certificate ($759$\,B when only the two
required signatures are carried). The threshold semantics are machine-checked:
any two distinct signatures pass, one signature is rejected, and two signatures
under the same key id do not count as two.

\begin{table}[t]
\centering
\caption{Gateway signing cost on a 24-core / 14.6\,GiB machine ($n=300$ per
row). Rows are measured medians. The in-circuit row is \emph{blocked}: the
circuit contains no Ed25519 verifier by design, and no cycle, time or memory
figure is reported. Source: \texttt{bench/results/gateway\_sig.md}.}
\label{tab:gateway}
\tabfont
\begin{tabular}{@{}lrrrr@{}}
\toprule
Configuration & Sign (ms) & Verify (ms) & Cert (B) & $n$ \\
\midrule
Ed25519, 1 receipt & 0.058 & 0.094 & 237.0 & 300 \\
Ed25519, 10 receipts & 0.415 & 0.844 & 288.3 & 300 \\
Ed25519, 100 receipts & 4.885 & 10.165 & 295.2 & 300 \\
\texttt{ToolSeal} (10 receipts) & 0.074 & 0.184 & 345 & 300 \\
1-of-1 & 0.069 & 0.173 & 648 & 300 \\
2-of-3 (3 signatures) & 0.173 & 0.508 & 981 & 300 \\
\midrule
In-circuit Ed25519 & \multicolumn{3}{c}{blocked: no guest, no number} & --- \\
\bottomrule
\end{tabular}
\end{table}

The in-circuit row is the honest blank. Verifying Ed25519 inside the zkVM is not
in the circuit---\texttt{circuits/} deliberately keeps the signature field out
of the circuit view---and building a guest to measure it exceeds this
experiment's file ownership, so cycles, proving time and peak memory are
\emph{null}, not zero. We therefore do not infer a circuit cost from the
off-circuit cost, and the trust root still contains the gateway key. What the
threshold prototype buys is not a cheaper verification but a weaker single point
of trust: one key compromised is no longer enough, at roughly three times the
signing work and certificate size, and with a new obligation to argue that the
gateways agree on the chain tail.

Putting Ed25519 inside the guest is a design option, not a measured one. SP1
6.7.0 has no Ed25519 precompile; the guest change would rotate the policy vkey
and add public values to the outcome, invalidating existing certificates, and we
did not build or time it. The only gateway-side variant measured is the 2-of-3
threshold prototype. It removes the single-key trust but not gateway collusion
or tail truncation. An authenticated gateway roster, HSM key custody and a
trusted time source remain open.

\subsection{Cross-system baselines: 49 cells, 46 populated}
\label{sec:eval-baselines}

Table~\ref{tab:baselines} places PoP against six alternatives across seven
dimensions: 49 cells, of which 46 are populated (measured or cited) and 3 are
unavailable. Every populated cell is tagged $[m]$ (measured in this repository
or this study) or $[c]$ (cited from the named paper). The three measured arms are
PoP, the signed-Merkle log and a simulated trusted server; zkAgent, NiyamAI,
Zaguk and PCO are cited. Zaguk and PCO did not yield comparable performance
numbers in our environment---the former's proceedings PDF was unreachable and
the latter reports no performance table---so those cells are em dashes, not
estimates. The seven dimensions impose different proof obligations; the table is
a statement of \emph{what each system obliges the verifier to trust}, not a
single ranking.

\begin{table*}[t]
\centering
\caption{Cross-system comparison. $[m]$ = measured here; $[c]$ = cited from the
source named in \texttt{bench/results/baselines.md}; em dash = not available,
not guessed. The trust root is per obligation, not one column: PoP's response
and policy proofs trust the vkey, while tool receipts additionally trust the
gateway key.}
\label{tab:baselines}
\tabfont
\setlength{\tabcolsep}{3pt}
\scriptsize
\begin{tabular}{@{}L{0.11\textwidth}L{0.13\textwidth}L{0.13\textwidth}L{0.12\textwidth}L{0.11\textwidth}L{0.10\textwidth}L{0.10\textwidth}L{0.08\textwidth}@{}}
\toprule
Dimension & PoP & signed-Merkle & TEE (sim) & zkAgent & NiyamAI & Zaguk & PCO \\
\midrule
Trust root & vkey + crypto $[m]$; tool-receipt gateway key & Ed25519 gateway $[m]$ &
server / TEE $[m]$ & zkVM + provider $[c]$ & SNARK + judge $[c]$ &
ezkl + commitment $[c]$ & authority + ledger $[c]$ \\
Latency & 40.5--100.3\,s core; 279--344\,s \texttt{gr\-oth16}; 89.84\,ms verify $[m]$ &
0.132\,ms sign; 2.131\,ms append; 0.322\,ms check $[m]$ &
0.139\,ms sign; 0.102\,ms verify $[m]$ &
194.4\,s / 0.42\,s $[c]$ & 1.7\,s / 51\,ms $[c]$ &
not comparable $[c]$ & --- \\
Size & 2.79\,MB core; 1{,}874\,B \texttt{gr\-oth16} $[m]$ & 648\,B + 761\,B $[m]$ &
648\,B $[m]$ & 96\,MiB (Shout) $[c]$ & 18.6\,KB $[c]$ & --- & --- \\
Response privacy & committed $[m]$ & none $[m]$ & none $[m]$ &
committed $[c]$ & not reported $[c]$ & yes $[c]$ & yes $[c]$ \\
Policy privacy & public values salted; proof artefact open $[m]$ & none $[m]$ & none $[m]$ &
n/a $[c]$ & contract commitment $[c]$ & yes $[c]$ & partial $[c]$ \\
Needs $T$ plaintext & no $[m]$ & yes $[m]$ & no $[m]$ & no $[c]$ &
unclear $[c]$ & no $[c]$ & no $[c]$ \\
Needs $\pi$ plaintext & no $[m]$ & yes $[m]$ & no $[m]$ & n/a $[c]$ &
unclear $[c]$ & no $[c]$ & partial $[c]$ \\
\bottomrule
\end{tabular}
\end{table*}

The asymmetry the table exposes is the one the signature baseline already
showed: the cheap arms require the verifier to re-run the judgment on plaintext,
while the private arms move that work into a proof or into a trusted party. PoP
is on the proof side, with a measured cost four to five orders of magnitude
above the signature arm and a trust root that, unlike the TEE arm, contains no
hardware.

\subsection{Issuance: detection, not prevention}
\label{sec:eval-issuance}

A provider that never issues a certificate is indistinguishable, from the
verifier's side, from a provider with nothing to certify; no cryptographic
layer forces issuance. What a transparency log changes is the \emph{detection}
rate, and only under an external obligation to append before delivery. In a
simulation with $N=2000$ deliveries, $m=100$ withheld violations and a sampling
rate $p=0.05$, an audit that samples the \emph{delivered} artifacts detects
$0.9948$ of the withheld set in exact arithmetic ($0.9940$ by Monte Carlo over
500 draws). An audit that samples the \emph{log} detects $0.0000$, because the
withheld deliveries have no log entry at all: the sampling frame is wrong, not
unlucky.

\begin{table}[t]
\centering
\caption{Issuance detection under two sampling frames ($N=2000$, $m=100$,
$p=0.05$, seed 20261101). Exact is $1-\binom{N-m}{s}/\binom{N}{s}$ with $s=100$;
Monte Carlo uses 500 draws. Source: \texttt{bench/results/issuance.md}.}
\label{tab:issuance}
\tabfont
\begin{tabular}{@{}lrr@{}}
\toprule
Sampling frame & Exact & Monte Carlo \\
\midrule
Delivered artifacts (correct frame) & 0.9948 & 0.9940 \\
Log entries (wrong frame) & 0.0000 & 0.0000 \\
\bottomrule
\end{tabular}
\end{table}

Two boundaries travel with these numbers. First, the result quantifies
\emph{detection}, never \emph{prevention}: whether a detected violation is
punished depends on penalties and incentives outside the model. Second, the
sufficiency direction assumes an append-before-delivery obligation, a log that
is append-only and publicly readable, inclusion proofs, and an external
mechanism against split views. The prototype implements append-only semantics,
inclusion proofs, optional tree-head signatures and checkpoints; it does
\emph{not} implement consistency proofs or gossip, so split-view resistance is
an external assumption, not a delivered property. Three premises travel with the
result: (F-1) the append-before-deliver obligation is external to cryptography;
(F-3) the log entry must bind the delivered $T$ (or a receipt), not only
$H(\mathrm{cert})$; and (F-4) the audit sample frame is delivered artifacts, not
log entries. The $1-(1-p)^m$ form assumes independent sampling and
detection-on-hit.

\subsection{Chat-template rerun and refusal proxy}
\label{sec:eval-k-chat}

A chat-template rerun of 20 prompts at $N=10$ returned $k=3.32=\log_2 10$ at
temperature $0.7$ and $k=0$ at temperature $0$; its refusal proxy rose from
$15\%$ to $35\%$ ($7/10$ unsafe prompts) while the response-violation rate fell
from $100\%$ to $90\%$. The $90\%$ is dominated by keyword-policy false
positives---$9/10$ safe responses were flagged---and is not a model harm rate;
XSTest labels are prompt-level and the refusal proxy is a prefix heuristic.

% =====================================================================
\section{Policy-Private Mode Details}
\label{app:private-policy}

\paragraph{Three parties.} ``Three'' means \emph{authority / operator / third
party}; an auditor is a named instance of the third party, not a fourth trust
anchor. Table~\ref{tab:wp4-parties} states the model.

\begin{table}[t]
\centering
\caption{The three-party model of policy-private mode.}
\label{tab:wp4-parties}
\tabfont
\setlength{\tabcolsep}{3pt}
\begin{tabular}{@{}L{0.16\columnwidth}L{0.14\columnwidth}L{0.62\columnwidth}@{}}
\toprule
Party & Honest? & Holds; wants \\
\midrule
Authority $R$ & \textbf{Yes}, by (P-PP) & Holds $\pi$ and the issuing record;
wants $\pi$ actually enforced, and wants to be able to check that it was. \\
Operator $O$ & \textbf{No} (the adversary) & Holds $\pi$, $T$, the proving key;
wants $\pi$ hidden from third parties while still being able to certify. \\
Third party $V$ & Yes & Holds only public values---\texttt{policy\_\allowbreak hash},
\texttt{circuit\_\allowbreak hash}, \texttt{response\_\allowbreak commitment},
\texttt{passed}, the proof; wants to confirm that \emph{some issued} policy
passed, and must \emph{not} learn $\pi$. \\
\emph{(Auditor $V^{*}$)} & Yes & As $V$, plus $\pi$ once $O$ discloses it;
wants ex-post accountability. \\
\bottomrule
\end{tabular}
\end{table}

Two boundaries must be stated with the model. First, secrecy of $\pi$ comes from
a \emph{data-flow fact}---$\pi$ never enters the public values---and not from
zero-knowledge; any rule that writes a fragment of $\pi$ into the public values
breaks secrecy on that fragment. Second, accountability for $V^{*}$ is
\emph{voluntary}: $V^{*}$ can check only when $O$ chooses to disclose $\pi$ and
the salt. The mode offers no forced disclosure; that would require the
release-gating family discussed in \texttt{docs/tdsc-plan.md}~\S1.6(b), which we
do not claim.

\paragraph{What each verifier may conclude.} Table~\ref{tab:wp4-capability} is
the contribution of this subsection: the cryptography is downstream of it, and
the table survives even if the cryptography only partially lands.

\begin{table*}[t]
\centering
\caption{What each verifier may and may not conclude in policy-private mode.}
\label{tab:wp4-capability}
\tabfont
\setlength{\tabcolsep}{4pt}
\begin{tabular}{@{}L{0.13\textwidth}L{0.30\textwidth}L{0.30\textwidth}L{0.19\textwidth}@{}}
\toprule
Verifier & \textbf{May} conclude & \textbf{May not} conclude & Rests on \\
\midrule
Third party $V$ & (i) the commitment binds $T$
(Definition~1 of the main paper); (ii) \emph{some} policy with commitment $h_\pi$
passed; (iii) $h_\pi$ appears in the issuing record & (iv) the \emph{content} of
$\pi$; (v) how many rules $\pi$ has, or of which classes; (vi) ``the operator
enforces the regulation \emph{I} care about''---only ``some issued regulation'' &
(P-PP); Assumption~\ref{as:Z}; the $k$ of Proposition~\ref{thm:t1} \\
End user (also receiver) & As $V$, and additionally: the proved $T$ is the
string it received, since its nonce is bound (Theorem~\ref{thm:t3}, requires
$f>0$)$^{\dagger}$ & Only the string it received; \emph{not} what was delivered
to another downstream agent & $f>0$ \\
Authority $R$ & Everything; it holds $\pi$, the salt and can recompute $h_\pi$ &
--- & --- \\
Auditor $V^{*}$ (after disclosure) & (vii) $\mathsf{verify\_commit}(\pi,s)=
\texttt{policy\_\allowbreak hash}$, i.e.\ what ran was the issued $\pi$;
(viii) the disclosed violation record is consistent & (ix) the undisclosed part
of $T$ & voluntary disclosure \\
On-chain verifier & (i)--(iii) only; it hashes and verifies signatures, and does
not interpret policy semantics & (iv)--(ix) & (P-PP) \\
\bottomrule
\end{tabular}
\end{table*}

The $V$ row is the reason the table exists. A third party does not obtain
``compliance''; it obtains ``some issued policy was satisfied''. Conflating the
two is the sentence a referee would use to reject the mode, and preventing that
sentence is what the table is for. $^{\dagger}$\,The end-user binding is
currently a certificate-layer field; the circuit still computes the
verifier-domain binding, so the end-to-end form needs a circuit change and a new
vkey.

% =====================================================================
\section{Related-Work Positioning}
\label{app:related}

Table~\ref{tab:related} records the slot each lineage occupies---which is also,
read the other way, what the main paper does not have to claim.
Table~\ref{tab:zkagent} places PoP against zkAgent: the two prove different
obligations, so the rows are a statement of orthogonality, not a race.

\begin{table}[t]
\centering
\caption{Positioning against the five prior-art lineages. The right-hand column
names what each lineage already occupies, and therefore what we do not claim.}
\label{tab:related}
\tabfont
\setlength{\tabcolsep}{3pt}
\begin{tabular}{@{}L{0.26\columnwidth}L{0.42\columnwidth}L{0.26\columnwidth}@{}}
\toprule
Lineage & Representative work & Slot already occupied \\
\midrule
Policy $\to$ circuit & zkGen \cite{zkgen}, Circuitree, Prezta & DSL-to-circuit
translation \\
Agent/LLM compliance & NiyamAI, Zaguk, PCO \cite{pco}, WITNESS & policy-level
private-mode kernel \\
Model computation & DeepProve, zkGPT, zkAgent \cite{zkagent} & inference
integrity \\
Attacks / negative results & Hollow-LLM \cite{hollowllm}, zk-attach
\cite{zkattach} & naming the gaps \\
Non-ZK industry baseline & \texttt{@brivora/verify}, VAOL & signed audit logs \\
\bottomrule
\end{tabular}
\end{table}

\begin{table}[t]
\centering
\caption{PoP against zkAgent. The two prove different obligations, so the rows
are a statement of orthogonality, not a race.}
\label{tab:zkagent}
\tabfont
\setlength{\tabcolsep}{3pt}
\begin{tabular}{@{}L{0.22\columnwidth}L{0.38\columnwidth}L{0.30\columnwidth}@{}}
\toprule
& \textbf{zkAgent} \cite{zkagent} & \textbf{PoP} \\
\midrule
Obligation & provider executed the declared model and tools (inference integrity)
& response/trace satisfies policy $\pi$ (compliance) \\
Depends on the model & yes (weights, arithmetisation) & \textbf{no} \\
Cost grows with & model size, sequence length & response length; rule-class
composition (per-rule-set slope), not rule count \\
Content privacy & none & committed, selective disclosure \\
Hardware reported & 32 cores / 512\,GB & 16 cores / 90\,GiB \\
Proving & 194.4\,s (Shout, $T{=}512$) & 40.5--100.3\,s core; 279--344\,s \texttt{gr\-oth16} \\
Proof size & 96\,MiB (Shout); 3.1\,MiB (LogUp) & 2.79\,MB core; 1{,}874\,B \texttt{gr\-oth16} \\
Verification & 0.42\,s (Shout); 0.038\,s (LogUp) & $\approx$0.090\,s \\
\bottomrule
\end{tabular}
\end{table}

We deliberately do not claim that PoP is faster, and the temptation to do so is
worth naming because the raw numbers invite it. The correct statement is
narrower, and we make only it: \emph{for the compliance obligation}, PoP
completes in a proving time of the same order as zkAgent's, with a smaller memory
footprint and a comparable proof size, without touching the model---and it adds
committed content privacy that zkAgent does not provide. Two caveats must travel
with those figures. The hardware is not matched (16 cores and a 90\,GiB cgroup
against 32 cores and 512\,GB), and PoP's ceiling is memory-bound rather than
cycle-bound, so its feasible region is a non-rectangular set of measured points
rather than a curve. The comparison is of orders of magnitude under unmatched
hardware, not a race. Source: \texttt{bench/comparison\_zkagent.md}, which
itself warns against reading the rows as ``PoP is faster.''

% =====================================================================
\section{Reproduction}
\label{app:repro}

The full walkthrough---environment, a single compliance proof, cross-checking
against the reference layer, private mode, certificate issuance, the one-command
demo, and third-party verification---is maintained in \texttt{docs/reproduce.md}
together with the exact commands and their expected output. This section records
only what a reader needs in order to judge the claims of the main paper's
evaluation.

\paragraph{Environment.} Ubuntu or WSL2 on x86-64, no GPU. The Rust/SP1 side
needs Rust stable, \texttt{cargo-prove}~v6.7.0, Go $\ge$\,1.24, and the usual
native toolchain; the reference layer is pure Python~3.10+ and runs on its own.
Memory, not time, is the binding constraint: the CPU prover carries a fixed
floor of approximately 10.5\,GiB on the cloud cgroup (10.15\,GiB under
default local prover options; \texttt{NUM\_CORE\_WORKERS=1},
\texttt{CORE\_BUFFER\_SIZE=1} lowers it to $\approx$9.3\,GiB) regardless of the
trace; memory, not time, sets the ceiling in the proof-cost table of the main
paper.

\paragraph{What reproduces without the prover.} The reference layer, the policy
compiler, and the whole test suite run in seconds and need no Rust. This is
enough to reproduce the non-ZK baseline of the main paper, to inspect the
certificate and inclusion proofs, and to check the cross-layer contract---that
the Python reference implementation and the zkVM guest consume the same
compiled policy bytes.

\paragraph{What needs the prover.} Every timing, size, and memory figure in the
cycle, ablation, denial-of-service, and proof-cost evaluations, and the
end-to-end result, requires a real proof. The measurements were taken on a
16-core CPU cgroup with 90\,GiB of RAM, which is what places the ceiling
reported there: the successes and the out-of-memory configurations are both
properties of that machine.

\paragraph{Where each number comes from.} Every measurement in the main paper's
evaluation is the content of a checked-in result file with a machine-readable
companion: \texttt{bench/results/} holds \texttt{cycles}, \texttt{ablation},
\texttt{scale-cloud}, \texttt{scale-6x2000}--\texttt{scale-6x5000},
\texttt{eu\_ai\_act\_proof}, \texttt{privacy\_surface\_real},
\texttt{k\_entropy\_largeN}, \texttt{k\_entropy\_crossmodel},
\texttt{policy\_corpus\_infer}, \texttt{agent\_e2e\_deepseek},
\texttt{baseline}, \texttt{verify}, \texttt{semantic}, \texttt{streaming},
\texttt{compose}, \texttt{dos\_proof}, \texttt{gateway\_sig},
\texttt{baselines}, \texttt{issuance}, and \texttt{prover\_knobs}, each as a
\texttt{.json} alongside its \texttt{.md}. The reproduction guide records the
commit, the hardware, and the prover binary digest for each run, so that a
reported figure can be traced to one specific execution rather than to a state
of the tree.

\begin{thebibliography}{99}

\bibitem{hollowllm}
C.~Gong, B.~Liu, and M.~Li, ``Hollow-LLM Attack: Computationally Trivial
Weights in Zero-Knowledge Verification of LLM Inference,'' in \emph{IEEE
Symposium on Security and Privacy (S\&P)}, 2026, pp.~4339--4355,
DOI 10.1109/SP63933.2026.00258; arXiv:2607.28884.

\bibitem{zkattach}
A.~S.~Gopalakrishna, ``Zero-Knowledge Predicate Proofs Between AI Agents: A
Measured, Cross-Protocol Gateway and the Source-Integrity Gap,''
arXiv:2608.30083, 2026. (The \texttt{zk-attach/v0} envelope is the system
contribution of this paper.)

\bibitem{zkgen}
J.~Lauinger, J.~Ernstberger, and S.~Steinhorst, ``zkGen: Policy-to-Circuit
Transpiler,'' in \emph{IEEE International Conference on Blockchain and
Cryptocurrency (ICBC)}, 2024, pp.~620--624,
DOI 10.1109/ICBC59979.2024.10634440.

\bibitem{pco}
H.~Avni, S.~Dolev, A.~Yagudaev, and M.~Yung, ``Super-intelligence Survival
Guide: Verification via Proof-Carrying Output,'' IACR Cryptology ePrint
Archive, Report 2026/994, 2026.

\bibitem{zkagent}
L.~Wang, H.~Lou, C.~Li, Y.~Yu, and Y.~Hu, ``zkAgent: Verifiable LLM Agent
Execution via One-Shot Transcript Proofs,'' IACR Cryptology ePrint Archive,
Report 2026/199, 2026.

\bibitem{hiddencred}
J.~E.~Holt, R.~W.~Bradshaw, K.~E.~Seamons, and H.~Orman, ``Hidden
Credentials,'' in \emph{Proc.\ 2003 ACM Workshop on Privacy in the Electronic
Society (WPES)}, 2003, pp.~1--8, DOI 10.1145/1005140.1005142.

\bibitem{handshake}
D.~Balfanz, G.~Durfee, N.~Shankar, D.~K.~Smetters, J.~Staddon, and H.-C.~Wong,
``Secret Handshakes from Pairing-Based Key Agreements,'' in \emph{Proc.\ IEEE
Symposium on Security and Privacy (S\&P)}, 2003, pp.~180--196,
DOI 10.1109/SECPRI.2003.1199336.

\bibitem{xacmlpriv}
C.~A.~Ardagna, S.~De~Capitani~di~Vimercati, G.~Neven, S.~Paraboschi,
F.-S.~Preiss, P.~Samarati, and M.~Verdicchio, ``Enabling Privacy-Preserving
Credential-Based Access Control with XACML and SAML,'' in \emph{Proc.\ 10th
IEEE International Conference on Computer and Information Technology (CIT)},
2010, pp.~1090--1095, DOI 10.1109/CIT.2010.199.

\bibitem{sw05}
A.~Sahai and B.~Waters, ``Fuzzy Identity-Based Encryption,'' in \emph{Advances
in Cryptology -- EUROCRYPT 2005}, LNCS 3494, Springer, 2005, pp.~457--473,
DOI 10.1007/11426639\_27.

\bibitem{ksw08}
J.~Katz, A.~Sahai, and B.~Waters, ``Predicate Encryption Supporting
Disjunctions, Polynomial Equations, and Inner Products,'' in \emph{Advances in
Cryptology -- EUROCRYPT 2008}, Springer, 2008, pp.~146--162,
DOI 10.1007/978-3-540-78967-3\_9.

\bibitem{vdr}
H.~Chabanne, R.~Hugel, and J.~Keuffer, ``Verifiable Document Redacting,'' in
\emph{Computer Security -- ESORICS 2017}, LNCS 10492, Springer, 2017,
pp.~334--351, DOI 10.1007/978-3-319-66402-6\_20.

\end{thebibliography}

\end{document}
"
%
%  This document carries the material removed from the 12--14 page main
%  paper to meet the TDSC length limit: full proofs, the detailed
%  evaluation tables, policy-private-mode detail, related-work tables,
%  and reproduction instructions. It is self-contained: it restates the
%  results it proves and can be read without the main PDF.
% =====================================================================
\documentclass[10pt,journal,compsoc]{IEEEtran}

\usepackage{amsmath,amsfonts,amssymb,amsthm}
\usepackage{algorithmic}
\usepackage{algorithm}
\usepackage{array}
\usepackage[caption=false,font=normalsize,labelfont=sf,textfont=sf]{subfig}
\usepackage{textcomp}
\usepackage{stfloats}
\usepackage{url}
\usepackage{verbatim}
\usepackage{graphicx}
\usepackage{cite}
\usepackage{booktabs}
\usepackage{xcolor}
\usepackage{microtype}
\usepackage{hyperref}

\hyphenation{Proof-of-Policy zero-knowl-edge re-exec-ution pol-icy
com-pli-ance bind-ing veri-fi-a-bil-ity cer-ti-fi-cate non-zero-knowl-edge}

\setlength{\emergencystretch}{2em}

\newcolumntype{L}[1]{>{\raggedright\arraybackslash}p{#1}}
\newcommand{\tabfont}{\footnotesize}

\newtheorem{theorem}{Theorem}
\newtheorem{lemma}[theorem]{Lemma}
\newtheorem{proposition}[theorem]{Proposition}
\newtheorem{corollary}[theorem]{Corollary}
\newtheorem{definition}[theorem]{Definition}
\newtheorem{assumption}[theorem]{Assumption}
\newtheorem{remark}[theorem]{Remark}

% Work markers are no-ops in the supplementary material.
\newcommand{\need}[1]{}
\newcommand{\fixme}[1]{}
\newcommand{\wpmark}[1]{}

% Supplementary numbering: S1, S2, ... and Table S1, S2, ...
\renewcommand{\thesection}{S\arabic{section}}
\renewcommand{\thetable}{S\arabic{table}}

\begin{document}

\title{Supplementary Material for\\``Proof-of-Policy: Zero-Knowledge Compliance
Certificates for LLM Agent Outputs''}
\author{Author Name(s)}
\maketitle

\begin{abstract}
This supplement contains the full proofs, the detailed evaluation tables, the
policy-private-mode tables, the related-work positioning, and the reproduction
instructions omitted from the main paper for length. Section and table numbers
are prefixed with ``S''. The main paper states every result; this document
supplies the arguments and the per-configuration numbers behind them.
\end{abstract}

% =====================================================================
\section{Full Proofs}
\label{app:proofs}

This section restates the results proved in the main paper and gives complete
proofs. We use the main paper's notation; for convenience, the two objects the
proofs depend on are repeated here.

\begin{definition}[Guessing entropy; restated from the main paper]
\label{def:k}
For a secret distribution $\mathcal{D}$ and adversary side information $Z$, let
\begin{multline*}
\mathsf{Guess}(X \mid Z) \;\triangleq\; \max_{\mathcal{A}} \Pr[\mathcal{A}(Z) = X]\\
\;=\; \mathbb{E}_z\!\left[\max_x \Pr[X = x \mid Z = z]\right],
\end{multline*}
\[
k_X \;\triangleq\; -\log_2 \mathsf{Guess}(X \mid Z),
\]
the maximum ranging over all time-unbounded adversaries.
\end{definition}

\begin{assumption}[Adversary side information; restated from the main paper]
\label{as:Z}
The adversary knows the prompt, the model (possibly its weights), and the
decoding parameters.
\end{assumption}

\begin{proposition}[Conditional trilemma; restated from the main paper]
\label{thm:t1}
Fix $\lambda,\varepsilon>0$ and $q\ge 1$. Let $X$ be a secret with guessing
entropy $k_X<\log q+\log(1/\varepsilon)$ whose binding is \emph{guess-tight}:
$\beta_\rho(X)=H(\mathsf{DOM}\|\rho\|X)$ for a random oracle $H$ and public
randomness $\rho$ independent of $X$ and $Z$. Then:
\begin{enumerate}
\item[(1)] \emph{Public binding.} For any adversary that evaluates $\beta$ at
most $q$ times,
\[
\Pr\!\left[\mathcal{A}^\beta(Z,\beta(X))=X\right]
\le q\cdot\mathsf{Guess}(X\mid Z)+\mathrm{negl}(\lambda)
= q\cdot 2^{-k_X}+\mathrm{negl}(\lambda).
\]
Hiding with advantage $\varepsilon$ is \emph{certified} by this bound when
$k_X\ge\log q+\log(1/\varepsilon)$; for a skewed distribution this is a
sufficient condition, not a matching attack. The integer $q$ counts evaluations
of $\beta$, not time: the adversary may be computationally unbounded.
\item[(2)] \emph{Keyed binding.} For $\beta_s(X)=H(s,X)$ with $s$ held only by a
designated party and absent from the public values,
$\Pr[\mathcal{A}(Z,b)=X]\le\mathsf{Guess}(X\mid Z)+\mathrm{negl}(\kappa)$.
What the keyed form buys is hiding when
$k_X<\log q+\log(1/\varepsilon)$; its price is PV-2.
\item[(3)] \emph{Conditional trilemma.} Consider schemes whose binding is
publicly recomputable (PV-2), whose hiding is exact recovery, and whose binding
is guess-tight. Within the scheme class defined above, when
$k_X<\log q+\log(1/\varepsilon)$ such a scheme cannot satisfy binding, hiding
and PV-2 simultaneously.
\end{enumerate}
\end{proposition}

\begin{proof}[Proof of Proposition~\ref{thm:t1}]
\emph{(1).} Fix $z$ and let the adversary's query set be
$Q = \{q_1, \dots, q_q\}$. Decompose by first hit,
\[
\Pr[\exists i : q_i = X] \;=\; \sum_i \Pr\!\left[q_i = X \wedge \neg E_{<i}\right],
\]
where $E_{<i} \triangleq \{\exists j < i : q_j = X\}$.
Conditioned on $\neg E_{<i}$, every oracle answer $a_1,\dots,a_{i-1}$ already
revealed is uniform and independent of $X$ (no queried point equals $X$, and $H$
is a fresh random function), and $\beta(X)$ itself is uniform and independent of
$X$. Hence $q_i$---a deterministic function of $(z, \beta(X), a_{<i})$---is
\emph{independent of $X$}, so
\begin{multline*}
\Pr[q_i = X \mid Z = z]
= \sum_x \Pr[q_i = x]\,\Pr[X = x \mid Z = z] \\
\le \max_x \Pr[X = x \mid Z = z] = 2^{-k_z},
\end{multline*}
giving $\Pr[\exists i : q_i = X \mid z] \le q\,2^{-k_z}$. Taking the expectation
over $z$ and using $\mathbb{E}_z[2^{-k_z}] = \mathsf{Guess}(X \mid Z) = 2^{-k_X}$
from Definition~\ref{def:k} yields $q \cdot 2^{-k_X}$. An adversary that outputs
an unqueried value wins with probability at most $2^{-k_z} \le q\,2^{-k_z}$,
which the same bound absorbs. No $q\cdot2^{-\lambda}$ term is needed in the
exact-recovery game: a query that lands on a different $x$ with
$\beta(x)=y$ is a second preimage of $H$, not a win. The explicit
$q\cdot2^{-\lambda}$ is kept only when the win condition is consistency with the
published $b$ rather than exact recovery, or when $q$ is unbounded.

\emph{(2).} Without an oracle for $\beta_s$, the value $b=H(s,X)$ is uniform and
independent of $X$ in the random-oracle model for an adversary that does not know
$s$, so the best strategy is the prior mode and succeeds with probability
$\mathsf{Guess}(X\mid Z)$. An attempted enumeration must also guess $s$, at cost
$2^{-\kappa}$ per query. The bound is therefore
$\mathsf{Guess}(X\mid Z)+\mathrm{negl}(\kappa)$, with no $q$-fold amplification.
Part (2) buys hiding only below the public-form threshold
$\log q+\log(1/\varepsilon)$; when $k_X=0$ the prior mode already recovers $X$
and no binding helps.

\emph{(3).} Public recomputability gives (1). If (1) exceeds $\varepsilon$, the
only way to avoid it is a binding that a verifier cannot recompute, i.e.\ a keyed
binding, which by part (2) forfeits PV-2. The third leg---hiding plus PV-2
without binding---is excluded from the class of compliance schemes because a
certificate whose binding does not pin $X$ certifies nothing about the delivered
artifact.
\end{proof}

\noindent\emph{Preimage-class bound.} For an arbitrary deterministic public
$\beta$, conditioning on $\beta(X)$ gives
$\Pr[\mathcal{A}^\beta(Z,\beta(X))=X]\le
\mathbb{E}_y[\max_x\Pr[X=x\mid Z,\beta(X)=y]]$, with no hypothesis on $\beta$.
For $\beta=\mathrm{id}$ or an injective $\mathsf{encode}$ this equals $1$; for a
lossy encoding it equals the inverse of the preimage-class size.

\begin{theorem}[D-privacy from three conditions; restated from the main paper]
\label{thm:dpriv}
The scheme is $D$-private if (C1) there is a PPT simulator $S_{\mathrm{rest}}$
such that, for every admissible $(\pi,T,n)$, the public values other than the
commitment satisfy
$(h_\pi,S_{\mathrm{rest}}(pp,h_\pi,D(\pi),\mathrm{pub}_T,n))\approx_c
L_{\mathrm{pub}}(\pi,T,n)$; (C2) the proof system is witness-hiding, i.e.\
$L_{\mathrm{proof}}$ has a simulator that takes only the public statement;
and (C3) the policy commitment is hiding for all $D$-equal policies,
$(pp,h_{\pi_0})\approx_c(pp,h_{\pi_1})$. The advantage is
$O(\varepsilon_{\mathrm{pub}}+\varepsilon_{\mathrm{wh}}+\varepsilon_{\mathrm{com}})$.
\end{theorem}

\begin{proof}[Proof of Theorem~\ref{thm:dpriv}]
Fix an admissible tuple $(\pi_0,\pi_1,T,n)$ with $D(\pi_0)=D(\pi_1)$ and
$R(\pi_0,T)=R(\pi_1,T)$. Write $h_i=h_{\pi_i}$,
$R_i=L_{\mathrm{rest}}(\pi_i,T,n)$, $P_i=L_{\mathrm{proof}}(\pi_i,T,n)$ and
$\mathsf{Stmt}_i=(h_i,R_i)$. Consider the hybrids
\begin{align*}
H_0&=(h_0,R_0,P_0),\\
H_1&=(h_0,R_0,S_{\mathrm{proof}}(\mathsf{Stmt}_0)),\\
H_2&=(h_0,S_{\mathrm{rest}}(h_0,D,\mathrm{pub}_T,n),S_{\mathrm{proof}}(\mathsf{Stmt}_0)),\\
H_3&=(h_1,S_{\mathrm{rest}}(h_1,D,\mathrm{pub}_T,n),S_{\mathrm{proof}}(\mathsf{Stmt}_1)),\\
H_4&=(h_1,R_1,S_{\mathrm{proof}}(\mathsf{Stmt}_1)),\\
H_5&=(h_1,R_1,P_1).
\end{align*}
$H_0\approx H_1$ and $H_4\approx H_5$ cost $\varepsilon_{\mathrm{wh}}$ each by
C2; $H_1\approx H_2$ and $H_3\approx H_4$ cost
$\varepsilon_{\mathrm{pub}}$ each by C1. For $H_2\approx H_3$, split into
$H_{2a}=(h_1,S_{\mathrm{rest}}(h_1,D,\mathrm{pub}_T,n),S_{\mathrm{proof}}(\mathsf{Stmt}_0))$
and $H_{2b}=H_3$: $H_2\approx H_{2a}$ costs $\varepsilon_{\mathrm{com}}$ by C3,
and $H_{2a}\approx H_{2b}$ follows because
$\mathsf{Stmt}_0\approx_c\mathsf{Stmt}_1$ (C1 and C3 by triangle inequality) and
$S_{\mathrm{proof}}$ is a PPT function of the statement. The triangle inequality
gives $4\varepsilon_{\mathrm{pub}}+2\varepsilon_{\mathrm{wh}}+2\varepsilon_{\mathrm{com}}$,
which is $O(\varepsilon_{\mathrm{pub}}+\varepsilon_{\mathrm{wh}}+\varepsilon_{\mathrm{com}})$.
The game fixes one $T$ because two different artifacts are distinguished by
$B(n,\cdot)$ itself, which is response-side leakage and not a policy advantage.
\end{proof}

\begin{theorem}[Leakage-profile decomposition; restated from the main paper]
\label{thm:leak}
Let $V=(L_{\mathrm{pub}},L_{\mathrm{proof}},L_{\mathrm{side}})$ be the verifier's
view, where $L_{\mathrm{pub}}$ is the certificate's public values,
$L_{\mathrm{proof}}$ is what the proof artefact reveals about the witness
$w=(\pi,T,\mathrm{trace})$, and $L_{\mathrm{side}}$ is the adversary's prior
information $Z$. For any $q\ge 1$ and any adversary that evaluates $\beta$ at
most $q$ times,
\[
\Pr\!\left[\mathcal{A}^\beta(V)=X\right]
\;\le\; q\cdot\mathsf{Guess}\!\left(X\mid L_{\mathrm{pub}},L_{\mathrm{proof}},L_{\mathrm{side}}\right)
+\mathrm{negl}(\lambda),
\]
with the final output charged as one evaluation.
\end{theorem}

\noindent\emph{Proof of Theorem~\ref{thm:leak}.} The proof is the first-hit
decomposition of Proposition~\ref{thm:t1}(1) with $(Z,\beta(X))$ replaced by the
verifier view $V$. The step that changes is the uniformity of answers to queries
that have not hit $X$: it holds conditionally on $L_{\mathrm{pub}}$ and fails when
$L_{\mathrm{proof}}$ is a function of the witness. A witness-hiding tier restores
$L_{\mathrm{proof}}=\bot$ and hence the public-value bound; a \texttt{core}
artefact does not, and the bound can reach $1$.

\begin{theorem}[Delivery binding; restated from the main paper]
\label{thm:t3}
Fix a policy $\pi$. A delivery adversary $\mathcal{B}$ may, \emph{before}
seeing the receiver's challenge, precompute a set $S$ of accepted transcripts,
where each transcript covers at most $c$ challenges. The receiver then
contributes $f$ fresh, uniformly random bits $r$ that enter the proven
statement. Under the certificate's soundness (D1) and challenge binding
(D2/B-bind), the divergence advantage satisfies
\[
\mathsf{Adv}_{\mathrm{del}}(\mathcal{B}) \;\le\; |S| \cdot c \cdot 2^{-f}
\;+\; \mathrm{negl}(\lambda).
\]
The replay-only probability---acceptance without requiring divergence---is at
most $|S|\cdot c\cdot 2^{-f}$. If the receiver contributes no fresh randomness
($f = 0$; e.g.\ an empty nonce), a precomputed transcript is accepted with
probability $1$ and \emph{does not decay in any parameter of the system}.
\end{theorem}

\begin{proof}[Proof of Theorem~\ref{thm:t3}]
The receiver's freshness $r$ enters the proven statement through
$b = H(\mathsf{dom} \,\|\, r \,\|\, T)$. For a certificate in $S$ to be accepted
on $r$, the value $r$ must fall inside that certificate's challenge set
$A(\mathrm{cert})$, with $|A(\mathrm{cert})|\le c$; a union bound over $S$ gives
$\sum_{\mathrm{cert}\in S}|A(\mathrm{cert})|/2^f\le |S|\cdot c\cdot 2^{-f}$.
A certificate minted \emph{after} $r$ is seen can diverge only by exhibiting two
distinct preimages of the same $b$ under $B(r,\cdot)$, which is exactly a
challenge-binding failure and costs $\mathrm{negl}(\lambda)$. When $f = 0$ every
element of $S$ is a permanently valid replay, so the replay probability is $1$
and does not decay in any parameter of the system. The $f = 0$ branch is not
hypothetical: \texttt{circuits/\allowbreak types/\allowbreak src/\allowbreak
lib.rs} (lines 950--995) documents an empty nonce as a legal input that
``provides no replay protection.'' The $f=0$ separation holds under stateless
public verification with no external mechanism (no revocation table, trusted
clock, or hardware). The $c$ factor is tight when the challenge sets of the
transcripts in $S$ are disjoint and of size $c$.
\end{proof}

\begin{proposition}[Quadratic cost of the naive matcher]
\label{prop:naive}
For the pattern \texttt{a+b} and input \texttt{a}$^n$, the naive backtracking
matcher performs
$\sum_{i=1}^{n}(n-i+1)=n(n+1)/2=\Theta(n^{2})$ NFA steps, whereas the Pike VM
performs $\Theta(n)$. The adversarial input family is admissible because the
response bytes are chosen by the party being checked. This bounds the cost of
\emph{that} baseline, not of policy-compliance matching: the Pike VM is linear
on the same family, and no lower bound for the problem is claimed.
\end{proposition}
\begin{proof}
No input position contains \texttt{b}, so no start position can fail or accept
before the end of the string. The start at $i$ therefore takes exactly
$n-i+1$ steps, and the sum is $n(n+1)/2$. The Pike VM keeps one state per NFA
position and scans once.
\end{proof}

\begin{proposition}[NFA contract]
\label{prop:nfa}
For the supported regular-expression subset $P$---ASCII, ordered and disjoint
character-class ranges, Unicode scalar values, and the same compiled policy on
both sides---the Thompson construction preserves the language, the Pike VM
implements unanchored substring search, and the Python reference and Rust
circuit matchers agree on the same compiled NFA specification. Empty character
classes are rejected at parse time (fail-closed), so ``non-empty character
class'' is a parser contract rather than an extra theorem assumption.
\end{proposition}
\begin{proof}[Proof sketch]
Language preservation is the standard Thompson construction over the restricted
syntax. The Pike VM simulates all NFA states at one input position and advances
once per character, so it decides unanchored search in one pass. The two
matchers consume the same \texttt{ConstraintSpec} bytes and implement the same
transition relation; agreement is checked by an algorithm-level bisimulation
plus differential fuzzing against Python's \texttt{re}, not by line-by-line
formal verification.
\end{proof}

% =====================================================================
\section{Evaluation Details}
\label{app:eval-details}

This section carries the per-configuration tables and the detailed readings
omitted from the main paper. Every number is the content of the checked-in
result file named in the caption, as in the main paper.

\subsection{Cycle sweep: length versus rule count}
\label{sec:eval-cycles}

Response length was pushed from 200 to 100{,}000 characters and the rule count
from 1 to 6, giving 40 measurement points (\texttt{bench/bench\_cycles.py}; the
full matrix is in \texttt{bench/results/cycles.md}). The response text comes from
an evaluation artefact the repository's own end-to-end demo produced (372
characters, three segments deduplicated by content); longer points are tiled from
it. What is measured is therefore the \emph{curve} of scanning cost, not ``a real
100{,}000-character document''. Every point was first confirmed to match no rule:
a matching rule short-circuits, and the cost of a short-circuit says nothing.

\begin{table}[t]
\centering
\caption{Cycle count against length, per rule set. Five rule sets, six length
points each; $R^2 \approx 1.0000$ throughout. Source:
\texttt{bench/results/cycles.md}.}
\label{tab:cycles}
\tabfont
\begin{tabular}{@{}rrrrr@{}}
\toprule
Rules & \texttt{keyword} & \texttt{pattern} & Slope (cyc/char) & Intercept \\
\midrule
1 & 0 & 0 & 73.9      & 50{,}137 \\
2 & 1 & 0 & 105.3     & 65{,}942 \\
3 & 1 & 1 & 4{,}289.8 & 387{,}501 \\
4 & 2 & 1 & 4{,}048.8 & 405{,}785 \\
6 & 3 & 2 & 5{,}403.5 & 993{,}257 \\
\bottomrule
\end{tabular}
\end{table}

Three findings, in decreasing order of how much they change what may be claimed.

\emph{(1) With the rule set fixed, cycle count is exactly linear in length.}
Five rule sets, six length points each, $R^2 \approx 1.0000$ throughout
(Table~\ref{tab:cycles}).

\emph{(2) Scanning cost is dominated by the NFA, not by the rule count.} One
keyword rule costs about $+31.6$ cycles per character at the margin; one NFA
regular expression (an email pattern) costs about $+3{,}634$---roughly
$115\times$ more. Equivalently, a rule set containing a regular expression has a
per-character slope ($4{,}289.8$) $58\times$ that of a rule set with only string
rules ($73.9$).

\emph{(3) Unit prices are not additive: a unit price is a property of a rule
\emph{in context}, not of the rule.} Fitting the slopes as
$(\text{keyword},\text{pattern})$ unit prices gives $R^2 = 0.9169$ and a
\emph{negative} keyword unit price ($-477.2$ cycles per character). This is not
noise; two probes give the mechanism directly:

\begin{itemize}
  \item \textbf{Ordering probe} (same rules, same text, $L = 10^5$): declared
  order costs $541{,}342{,}069$ cycles against $570{,}736{,}989$ in reverse,
  a difference of $+5.43\%$, while \texttt{passed} is identical in both. \emph{The
  semantics of the decision do not depend on order; the cost does.}
  \item \textbf{Marginal probe} ($L = 10^5$, filling the one matrix cell the
  sweep could not reach---regular expressions with no keyword rule):
  $[\text{len},\text{pat}]$ over $[\text{len}]$ has a marginal unit price of
  $3{,}633.7$, whereas ``add one more keyword rule'' has two different answers:
  $+31.6$ alone, $+585.6$ in a policy that already contains a regular expression
  ($-153.7$ if inserted before it). The same action changes sign between
  $-154$ and $+586$.
\end{itemize}

The mechanism lies in zkVM memory allocation: every \texttt{ascii\_lower} and
every Pike VM match allocates, and allocator state depends on previous
allocations. We therefore quote \emph{per-rule-set slopes} (the table above,
directly citable) rather than any ``unit price $\times$ count'' expression. The
simple model $\text{cycles} \approx a\cdot|T|\cdot\text{rules} + b$ that this
work anticipated is refuted by the data ($R^2 = 0.872$, negative intercept), and
we report that too: \emph{a model refuted by its own data is itself an evaluation
result.}

\subsection{Matcher ablation and proof-path cost}
\label{sec:eval-ablation}

\emph{Construction.} Pattern \texttt{a+b} against input
\texttt{a}$^{\times n}$ containing no \texttt{b}, so the match is necessarily
absent and both implementations must scan to the end. The naive implementation
reinvents and restarts at every start position, giving $\Theta(n^2)$; the Pike VM
decides in a single linear pass (\texttt{bench/bench\_ablation.py}).
Table~\ref{tab:ablation} reports the cycles. Read per-character cost rather than
absolute figures. Doubling $n$ leaves the Pike VM's per-character cost
essentially unchanged ($-11.1\%$, $-6.4\%$) while the naive path's nearly doubles
($+96.4\%$, $+98.6\%$). At $n = 400$ the naive path is already $90\times$ slower,
and that ratio doubles again with each doubling of $n$. \emph{This is an attack
surface}: the input length for policy compliance is chosen externally---an
agent's reply, a concatenated session---so an adversary who needs only to write a
long reply drives a naive implementation's cost quadratically. The comparison is
meaningful only because the two paths decide identically, which
\texttt{tests/test\_ablation.py} pins down (Python and Rust
\texttt{match\_search} $\equiv$ \texttt{match\_search\_naive}). In the systems we
surveyed, we are not aware of ablation evidence specific to policy-compliance
matching.

\begin{table}[t]
\centering
\caption{Matcher ablation for pattern \texttt{a+b} on input \texttt{a}$^n$.
Cycles are zkVM execution counts. Source:
\texttt{bench/results/ablation.md}.}
\label{tab:ablation}
\tabfont
\begin{tabular}{@{}rrrrr@{}}
\toprule
$n$ & Naive cycles & Pike cycles & Ratio \\
\midrule
100 & 12{,}382{,}720  & 643{,}894   & $19.2\times$ \\
200 & 48{,}638{,}946  & 1{,}145{,}209 & $42.5\times$ \\
400 & 193{,}207{,}130 & 2{,}143{,}617 & $90.1\times$ \\
\bottomrule
\end{tabular}
\end{table}

\paragraph{Proof-path cost.} Re-running the same construction through the
compliance driver (\texttt{bench/results/dos\_proof.md}) reproduces the ratios
on the execution path---$643{,}906$ against $12{,}382{,}732$ cycles at $n=100$,
rising to $90.1\times$ at $n=400$---and then prices the proving path. Proving is
where the quadratic term meets the prover's fixed memory floor: the Pike path at
$n=100$ succeeds at $122.2$\,s and $11{,}695$\,MiB; the naive path at $n=25$
succeeds at $153.2$\,s and $12{,}137$\,MiB; at $n=50$ it is
\textbf{OOM/SIGKILL} after $226.5$\,s at $13{,}772$\,MiB peak. We report the OOM
as a result rather than a discarded run: it is the point at which the
adversarial quadratic term exhausts a $14.6$\,GiB machine, and it is why the
$n\ge100$ naive points and all $n=400$ proving points are left to a cloud
machine. The cycle counts differ from these execution-path figures by about one
part in $10^5$, because they come from a separate driver build; the ratios
agree. These proving runs predate the NFA fix discussed in the main paper and
carry that binary's digest; the fix changes the guest and the verification key,
so the figures are not mixed with any post-fix run.

\begin{table}[t]
\centering
\caption{Proof-path denial of service on a 24-core / 14.6\,GiB machine. The
\texttt{naive} row at $n=50$ is an observed out-of-memory kill, retained as a
result. Source: \texttt{bench/results/dos\_proof.md}.}
\label{tab:dos}
\tabfont
\setlength{\tabcolsep}{3pt}
\begin{tabular}{@{}lrrrl@{}}
\toprule
Point & Cycles & Prove (s) & RSS (MiB) & Result \\
\midrule
pike $n{=}100$ & 643{,}906 & 122.2 & 11{,}695.1 & \checkmark \\
naive $n{=}25$ & --- & 153.2 & 12{,}137.3 & \checkmark \\
naive $n{=}50$ & --- & 226.5 & 13{,}771.6 & OOM/SIGKILL \\
\bottomrule
\end{tabular}
\end{table}

\subsection{Composition: four negative controls}
\label{sec:eval-compose}

The leakage-profile theorem of the main paper prices each half; it does not say
that conjoining two valid halves yields a valid composite. We built the four
cells of Table~\ref{tab:compose} against the composite verifier's binding layer,
with the cryptographic validity of the two sub-proofs covered separately by the
end-to-end tests. Only cell~(a) satisfies the strong form, and only because
$x_U=\mathsf{derive}(T_U)$; the implementation claims only the weak form of the
composition proposition in the main paper and its E-1 remark. Cell~(c) is the
one that matters: the current verifier binds the inference half to
$\mathsf{derive}(T_U)$, not to the input $x_U$ the user submitted, so it accepts
a certificate that proves the weak form
$\exists x: x=\mathsf{derive}(T_U)\wedge T_U=M(x)$ while the strong form
$x=x_U$ remains unproved. This is proof debt E-1, not a test artefact; it is the
concrete reason composition alone does not close the gap.

\begin{table}[t]
\centering
\caption{Composition negative controls. \texttt{ok} is the verifier's binding
check; \texttt{satisfied} is the policy outcome. Source:
\texttt{bench/results/composition\_controls.md}.}
\label{tab:compose}
\tabfont
\setlength{\tabcolsep}{3pt}
\begin{tabular}{@{}lL{0.34\columnwidth}ccL{0.23\columnwidth}@{}}
\toprule
Cell & Construction & ok & sat. & Reading \\
\midrule
(a) & same $n$; $x_U=\mathsf{derive}(T_U)$; \texttt{delegated} empty &
\checkmark & \checkmark & true composition \\
(b) & the two halves use different nonces & $\times$ & $\times$ &
rejected by the four-way binding check \\
(c) & $x_U\neq\mathsf{derive}(T_U)$; inference binds only the derived input &
\checkmark & \checkmark & \textbf{weak form only}; $x=x_U$ not landed (E-1) \\
(d) & \texttt{delegated} non-empty & \checkmark & $\times$ &
true but unsatisfied; fail closed \\
\bottomrule
\end{tabular}
\end{table}

\subsection{Gateway signing: measured cost, blocked circuit}
\label{sec:eval-gateway}

Table~\ref{tab:gateway} prices the trust boundary the main paper states rather
than hides. Off-circuit Ed25519 verification costs $0.094$, $0.844$ and
$10.165$\,ms per receipt batch at $1$, $10$ and $100$ receipts, with a
per-receipt marginal near $0.08$--$0.10$\,ms; the session-end \texttt{ToolSeal}
is $0.074$\,ms to sign, $0.184$\,ms to verify and $345$\,B. A concatenated
2-of-3 threshold prototype signs in $0.173$\,ms, merges in $0.001$\,ms, verifies
in $0.508$\,ms and produces a $981$\,B certificate ($759$\,B when only the two
required signatures are carried). The threshold semantics are machine-checked:
any two distinct signatures pass, one signature is rejected, and two signatures
under the same key id do not count as two.

\begin{table}[t]
\centering
\caption{Gateway signing cost on a 24-core / 14.6\,GiB machine ($n=300$ per
row). Rows are measured medians. The in-circuit row is \emph{blocked}: the
circuit contains no Ed25519 verifier by design, and no cycle, time or memory
figure is reported. Source: \texttt{bench/results/gateway\_sig.md}.}
\label{tab:gateway}
\tabfont
\begin{tabular}{@{}lrrrr@{}}
\toprule
Configuration & Sign (ms) & Verify (ms) & Cert (B) & $n$ \\
\midrule
Ed25519, 1 receipt & 0.058 & 0.094 & 237.0 & 300 \\
Ed25519, 10 receipts & 0.415 & 0.844 & 288.3 & 300 \\
Ed25519, 100 receipts & 4.885 & 10.165 & 295.2 & 300 \\
\texttt{ToolSeal} (10 receipts) & 0.074 & 0.184 & 345 & 300 \\
1-of-1 & 0.069 & 0.173 & 648 & 300 \\
2-of-3 (3 signatures) & 0.173 & 0.508 & 981 & 300 \\
\midrule
In-circuit Ed25519 & \multicolumn{3}{c}{blocked: no guest, no number} & --- \\
\bottomrule
\end{tabular}
\end{table}

The in-circuit row is the honest blank. Verifying Ed25519 inside the zkVM is not
in the circuit---\texttt{circuits/} deliberately keeps the signature field out
of the circuit view---and building a guest to measure it exceeds this
experiment's file ownership, so cycles, proving time and peak memory are
\emph{null}, not zero. We therefore do not infer a circuit cost from the
off-circuit cost, and the trust root still contains the gateway key. What the
threshold prototype buys is not a cheaper verification but a weaker single point
of trust: one key compromised is no longer enough, at roughly three times the
signing work and certificate size, and with a new obligation to argue that the
gateways agree on the chain tail.

Putting Ed25519 inside the guest is a design option, not a measured one. SP1
6.7.0 has no Ed25519 precompile; the guest change would rotate the policy vkey
and add public values to the outcome, invalidating existing certificates, and we
did not build or time it. The only gateway-side variant measured is the 2-of-3
threshold prototype. It removes the single-key trust but not gateway collusion
or tail truncation. An authenticated gateway roster, HSM key custody and a
trusted time source remain open.

\subsection{Cross-system baselines: 49 cells, 46 populated}
\label{sec:eval-baselines}

Table~\ref{tab:baselines} places PoP against six alternatives across seven
dimensions: 49 cells, of which 46 are populated (measured or cited) and 3 are
unavailable. Every populated cell is tagged $[m]$ (measured in this repository
or this study) or $[c]$ (cited from the named paper). The three measured arms are
PoP, the signed-Merkle log and a simulated trusted server; zkAgent, NiyamAI,
Zaguk and PCO are cited. Zaguk and PCO did not yield comparable performance
numbers in our environment---the former's proceedings PDF was unreachable and
the latter reports no performance table---so those cells are em dashes, not
estimates. The seven dimensions impose different proof obligations; the table is
a statement of \emph{what each system obliges the verifier to trust}, not a
single ranking.

\begin{table*}[t]
\centering
\caption{Cross-system comparison. $[m]$ = measured here; $[c]$ = cited from the
source named in \texttt{bench/results/baselines.md}; em dash = not available,
not guessed. The trust root is per obligation, not one column: PoP's response
and policy proofs trust the vkey, while tool receipts additionally trust the
gateway key.}
\label{tab:baselines}
\tabfont
\setlength{\tabcolsep}{3pt}
\scriptsize
\begin{tabular}{@{}L{0.11\textwidth}L{0.13\textwidth}L{0.13\textwidth}L{0.12\textwidth}L{0.11\textwidth}L{0.10\textwidth}L{0.10\textwidth}L{0.08\textwidth}@{}}
\toprule
Dimension & PoP & signed-Merkle & TEE (sim) & zkAgent & NiyamAI & Zaguk & PCO \\
\midrule
Trust root & vkey + crypto $[m]$; tool-receipt gateway key & Ed25519 gateway $[m]$ &
server / TEE $[m]$ & zkVM + provider $[c]$ & SNARK + judge $[c]$ &
ezkl + commitment $[c]$ & authority + ledger $[c]$ \\
Latency & 40.5--100.3\,s core; 279--344\,s \texttt{gr\-oth16}; 89.84\,ms verify $[m]$ &
0.132\,ms sign; 2.131\,ms append; 0.322\,ms check $[m]$ &
0.139\,ms sign; 0.102\,ms verify $[m]$ &
194.4\,s / 0.42\,s $[c]$ & 1.7\,s / 51\,ms $[c]$ &
not comparable $[c]$ & --- \\
Size & 2.79\,MB core; 1{,}874\,B \texttt{gr\-oth16} $[m]$ & 648\,B + 761\,B $[m]$ &
648\,B $[m]$ & 96\,MiB (Shout) $[c]$ & 18.6\,KB $[c]$ & --- & --- \\
Response privacy & committed $[m]$ & none $[m]$ & none $[m]$ &
committed $[c]$ & not reported $[c]$ & yes $[c]$ & yes $[c]$ \\
Policy privacy & public values salted; proof artefact open $[m]$ & none $[m]$ & none $[m]$ &
n/a $[c]$ & contract commitment $[c]$ & yes $[c]$ & partial $[c]$ \\
Needs $T$ plaintext & no $[m]$ & yes $[m]$ & no $[m]$ & no $[c]$ &
unclear $[c]$ & no $[c]$ & no $[c]$ \\
Needs $\pi$ plaintext & no $[m]$ & yes $[m]$ & no $[m]$ & n/a $[c]$ &
unclear $[c]$ & no $[c]$ & partial $[c]$ \\
\bottomrule
\end{tabular}
\end{table*}

The asymmetry the table exposes is the one the signature baseline already
showed: the cheap arms require the verifier to re-run the judgment on plaintext,
while the private arms move that work into a proof or into a trusted party. PoP
is on the proof side, with a measured cost four to five orders of magnitude
above the signature arm and a trust root that, unlike the TEE arm, contains no
hardware.

\subsection{Issuance: detection, not prevention}
\label{sec:eval-issuance}

A provider that never issues a certificate is indistinguishable, from the
verifier's side, from a provider with nothing to certify; no cryptographic
layer forces issuance. What a transparency log changes is the \emph{detection}
rate, and only under an external obligation to append before delivery. In a
simulation with $N=2000$ deliveries, $m=100$ withheld violations and a sampling
rate $p=0.05$, an audit that samples the \emph{delivered} artifacts detects
$0.9948$ of the withheld set in exact arithmetic ($0.9940$ by Monte Carlo over
500 draws). An audit that samples the \emph{log} detects $0.0000$, because the
withheld deliveries have no log entry at all: the sampling frame is wrong, not
unlucky.

\begin{table}[t]
\centering
\caption{Issuance detection under two sampling frames ($N=2000$, $m=100$,
$p=0.05$, seed 20261101). Exact is $1-\binom{N-m}{s}/\binom{N}{s}$ with $s=100$;
Monte Carlo uses 500 draws. Source: \texttt{bench/results/issuance.md}.}
\label{tab:issuance}
\tabfont
\begin{tabular}{@{}lrr@{}}
\toprule
Sampling frame & Exact & Monte Carlo \\
\midrule
Delivered artifacts (correct frame) & 0.9948 & 0.9940 \\
Log entries (wrong frame) & 0.0000 & 0.0000 \\
\bottomrule
\end{tabular}
\end{table}

Two boundaries travel with these numbers. First, the result quantifies
\emph{detection}, never \emph{prevention}: whether a detected violation is
punished depends on penalties and incentives outside the model. Second, the
sufficiency direction assumes an append-before-delivery obligation, a log that
is append-only and publicly readable, inclusion proofs, and an external
mechanism against split views. The prototype implements append-only semantics,
inclusion proofs, optional tree-head signatures and checkpoints; it does
\emph{not} implement consistency proofs or gossip, so split-view resistance is
an external assumption, not a delivered property. Three premises travel with the
result: (F-1) the append-before-deliver obligation is external to cryptography;
(F-3) the log entry must bind the delivered $T$ (or a receipt), not only
$H(\mathrm{cert})$; and (F-4) the audit sample frame is delivered artifacts, not
log entries. The $1-(1-p)^m$ form assumes independent sampling and
detection-on-hit.

\subsection{Chat-template rerun and refusal proxy}
\label{sec:eval-k-chat}

A chat-template rerun of 20 prompts at $N=10$ returned $k=3.32=\log_2 10$ at
temperature $0.7$ and $k=0$ at temperature $0$; its refusal proxy rose from
$15\%$ to $35\%$ ($7/10$ unsafe prompts) while the response-violation rate fell
from $100\%$ to $90\%$. The $90\%$ is dominated by keyword-policy false
positives---$9/10$ safe responses were flagged---and is not a model harm rate;
XSTest labels are prompt-level and the refusal proxy is a prefix heuristic.

% =====================================================================
\section{Policy-Private Mode Details}
\label{app:private-policy}

\paragraph{Three parties.} ``Three'' means \emph{authority / operator / third
party}; an auditor is a named instance of the third party, not a fourth trust
anchor. Table~\ref{tab:wp4-parties} states the model.

\begin{table}[t]
\centering
\caption{The three-party model of policy-private mode.}
\label{tab:wp4-parties}
\tabfont
\setlength{\tabcolsep}{3pt}
\begin{tabular}{@{}L{0.16\columnwidth}L{0.14\columnwidth}L{0.62\columnwidth}@{}}
\toprule
Party & Honest? & Holds; wants \\
\midrule
Authority $R$ & \textbf{Yes}, by (P-PP) & Holds $\pi$ and the issuing record;
wants $\pi$ actually enforced, and wants to be able to check that it was. \\
Operator $O$ & \textbf{No} (the adversary) & Holds $\pi$, $T$, the proving key;
wants $\pi$ hidden from third parties while still being able to certify. \\
Third party $V$ & Yes & Holds only public values---\texttt{policy\_\allowbreak hash},
\texttt{circuit\_\allowbreak hash}, \texttt{response\_\allowbreak commitment},
\texttt{passed}, the proof; wants to confirm that \emph{some issued} policy
passed, and must \emph{not} learn $\pi$. \\
\emph{(Auditor $V^{*}$)} & Yes & As $V$, plus $\pi$ once $O$ discloses it;
wants ex-post accountability. \\
\bottomrule
\end{tabular}
\end{table}

Two boundaries must be stated with the model. First, secrecy of $\pi$ comes from
a \emph{data-flow fact}---$\pi$ never enters the public values---and not from
zero-knowledge; any rule that writes a fragment of $\pi$ into the public values
breaks secrecy on that fragment. Second, accountability for $V^{*}$ is
\emph{voluntary}: $V^{*}$ can check only when $O$ chooses to disclose $\pi$ and
the salt. The mode offers no forced disclosure; that would require the
release-gating family discussed in \texttt{docs/tdsc-plan.md}~\S1.6(b), which we
do not claim.

\paragraph{What each verifier may conclude.} Table~\ref{tab:wp4-capability} is
the contribution of this subsection: the cryptography is downstream of it, and
the table survives even if the cryptography only partially lands.

\begin{table*}[t]
\centering
\caption{What each verifier may and may not conclude in policy-private mode.}
\label{tab:wp4-capability}
\tabfont
\setlength{\tabcolsep}{4pt}
\begin{tabular}{@{}L{0.13\textwidth}L{0.30\textwidth}L{0.30\textwidth}L{0.19\textwidth}@{}}
\toprule
Verifier & \textbf{May} conclude & \textbf{May not} conclude & Rests on \\
\midrule
Third party $V$ & (i) the commitment binds $T$
(Definition~1 of the main paper); (ii) \emph{some} policy with commitment $h_\pi$
passed; (iii) $h_\pi$ appears in the issuing record & (iv) the \emph{content} of
$\pi$; (v) how many rules $\pi$ has, or of which classes; (vi) ``the operator
enforces the regulation \emph{I} care about''---only ``some issued regulation'' &
(P-PP); Assumption~\ref{as:Z}; the $k$ of Proposition~\ref{thm:t1} \\
End user (also receiver) & As $V$, and additionally: the proved $T$ is the
string it received, since its nonce is bound (Theorem~\ref{thm:t3}, requires
$f>0$)$^{\dagger}$ & Only the string it received; \emph{not} what was delivered
to another downstream agent & $f>0$ \\
Authority $R$ & Everything; it holds $\pi$, the salt and can recompute $h_\pi$ &
--- & --- \\
Auditor $V^{*}$ (after disclosure) & (vii) $\mathsf{verify\_commit}(\pi,s)=
\texttt{policy\_\allowbreak hash}$, i.e.\ what ran was the issued $\pi$;
(viii) the disclosed violation record is consistent & (ix) the undisclosed part
of $T$ & voluntary disclosure \\
On-chain verifier & (i)--(iii) only; it hashes and verifies signatures, and does
not interpret policy semantics & (iv)--(ix) & (P-PP) \\
\bottomrule
\end{tabular}
\end{table*}

The $V$ row is the reason the table exists. A third party does not obtain
``compliance''; it obtains ``some issued policy was satisfied''. Conflating the
two is the sentence a referee would use to reject the mode, and preventing that
sentence is what the table is for. $^{\dagger}$\,The end-user binding is
currently a certificate-layer field; the circuit still computes the
verifier-domain binding, so the end-to-end form needs a circuit change and a new
vkey.

% =====================================================================
\section{Related-Work Positioning}
\label{app:related}

Table~\ref{tab:related} records the slot each lineage occupies---which is also,
read the other way, what the main paper does not have to claim.
Table~\ref{tab:zkagent} places PoP against zkAgent: the two prove different
obligations, so the rows are a statement of orthogonality, not a race.

\begin{table}[t]
\centering
\caption{Positioning against the five prior-art lineages. The right-hand column
names what each lineage already occupies, and therefore what we do not claim.}
\label{tab:related}
\tabfont
\setlength{\tabcolsep}{3pt}
\begin{tabular}{@{}L{0.26\columnwidth}L{0.42\columnwidth}L{0.26\columnwidth}@{}}
\toprule
Lineage & Representative work & Slot already occupied \\
\midrule
Policy $\to$ circuit & zkGen \cite{zkgen}, Circuitree, Prezta & DSL-to-circuit
translation \\
Agent/LLM compliance & NiyamAI, Zaguk, PCO \cite{pco}, WITNESS & policy-level
private-mode kernel \\
Model computation & DeepProve, zkGPT, zkAgent \cite{zkagent} & inference
integrity \\
Attacks / negative results & Hollow-LLM \cite{hollowllm}, zk-attach
\cite{zkattach} & naming the gaps \\
Non-ZK industry baseline & \texttt{@brivora/verify}, VAOL & signed audit logs \\
\bottomrule
\end{tabular}
\end{table}

\begin{table}[t]
\centering
\caption{PoP against zkAgent. The two prove different obligations, so the rows
are a statement of orthogonality, not a race.}
\label{tab:zkagent}
\tabfont
\setlength{\tabcolsep}{3pt}
\begin{tabular}{@{}L{0.22\columnwidth}L{0.38\columnwidth}L{0.30\columnwidth}@{}}
\toprule
& \textbf{zkAgent} \cite{zkagent} & \textbf{PoP} \\
\midrule
Obligation & provider executed the declared model and tools (inference integrity)
& response/trace satisfies policy $\pi$ (compliance) \\
Depends on the model & yes (weights, arithmetisation) & \textbf{no} \\
Cost grows with & model size, sequence length & response length; rule-class
composition (per-rule-set slope), not rule count \\
Content privacy & none & committed, selective disclosure \\
Hardware reported & 32 cores / 512\,GB & 16 cores / 90\,GiB \\
Proving & 194.4\,s (Shout, $T{=}512$) & 40.5--100.3\,s core; 279--344\,s \texttt{gr\-oth16} \\
Proof size & 96\,MiB (Shout); 3.1\,MiB (LogUp) & 2.79\,MB core; 1{,}874\,B \texttt{gr\-oth16} \\
Verification & 0.42\,s (Shout); 0.038\,s (LogUp) & $\approx$0.090\,s \\
\bottomrule
\end{tabular}
\end{table}

We deliberately do not claim that PoP is faster, and the temptation to do so is
worth naming because the raw numbers invite it. The correct statement is
narrower, and we make only it: \emph{for the compliance obligation}, PoP
completes in a proving time of the same order as zkAgent's, with a smaller memory
footprint and a comparable proof size, without touching the model---and it adds
committed content privacy that zkAgent does not provide. Two caveats must travel
with those figures. The hardware is not matched (16 cores and a 90\,GiB cgroup
against 32 cores and 512\,GB), and PoP's ceiling is memory-bound rather than
cycle-bound, so its feasible region is a non-rectangular set of measured points
rather than a curve. The comparison is of orders of magnitude under unmatched
hardware, not a race. Source: \texttt{bench/comparison\_zkagent.md}, which
itself warns against reading the rows as ``PoP is faster.''

% =====================================================================
\section{Reproduction}
\label{app:repro}

The full walkthrough---environment, a single compliance proof, cross-checking
against the reference layer, private mode, certificate issuance, the one-command
demo, and third-party verification---is maintained in \texttt{docs/reproduce.md}
together with the exact commands and their expected output. This section records
only what a reader needs in order to judge the claims of the main paper's
evaluation.

\paragraph{Environment.} Ubuntu or WSL2 on x86-64, no GPU. The Rust/SP1 side
needs Rust stable, \texttt{cargo-prove}~v6.7.0, Go $\ge$\,1.24, and the usual
native toolchain; the reference layer is pure Python~3.10+ and runs on its own.
Memory, not time, is the binding constraint: the CPU prover carries a fixed
floor of approximately 10.5\,GiB on the cloud cgroup (10.15\,GiB under
default local prover options; \texttt{NUM\_CORE\_WORKERS=1},
\texttt{CORE\_BUFFER\_SIZE=1} lowers it to $\approx$9.3\,GiB) regardless of the
trace; memory, not time, sets the ceiling in the proof-cost table of the main
paper.

\paragraph{What reproduces without the prover.} The reference layer, the policy
compiler, and the whole test suite run in seconds and need no Rust. This is
enough to reproduce the non-ZK baseline of the main paper, to inspect the
certificate and inclusion proofs, and to check the cross-layer contract---that
the Python reference implementation and the zkVM guest consume the same
compiled policy bytes.

\paragraph{What needs the prover.} Every timing, size, and memory figure in the
cycle, ablation, denial-of-service, and proof-cost evaluations, and the
end-to-end result, requires a real proof. The measurements were taken on a
16-core CPU cgroup with 90\,GiB of RAM, which is what places the ceiling
reported there: the successes and the out-of-memory configurations are both
properties of that machine.

\paragraph{Where each number comes from.} Every measurement in the main paper's
evaluation is the content of a checked-in result file with a machine-readable
companion: \texttt{bench/results/} holds \texttt{cycles}, \texttt{ablation},
\texttt{scale-cloud}, \texttt{scale-6x2000}--\texttt{scale-6x5000},
\texttt{eu\_ai\_act\_proof}, \texttt{privacy\_surface\_real},
\texttt{k\_entropy\_largeN}, \texttt{k\_entropy\_crossmodel},
\texttt{policy\_corpus\_infer}, \texttt{agent\_e2e\_deepseek},
\texttt{baseline}, \texttt{verify}, \texttt{semantic}, \texttt{streaming},
\texttt{compose}, \texttt{dos\_proof}, \texttt{gateway\_sig},
\texttt{baselines}, \texttt{issuance}, and \texttt{prover\_knobs}, each as a
\texttt{.json} alongside its \texttt{.md}. The reproduction guide records the
commit, the hardware, and the prover binary digest for each run, so that a
reported figure can be traced to one specific execution rather than to a state
of the tree.

\begin{thebibliography}{99}

\bibitem{hollowllm}
C.~Gong, B.~Liu, and M.~Li, ``Hollow-LLM Attack: Computationally Trivial
Weights in Zero-Knowledge Verification of LLM Inference,'' in \emph{IEEE
Symposium on Security and Privacy (S\&P)}, 2026, pp.~4339--4355,
DOI 10.1109/SP63933.2026.00258; arXiv:2607.28884.

\bibitem{zkattach}
A.~S.~Gopalakrishna, ``Zero-Knowledge Predicate Proofs Between AI Agents: A
Measured, Cross-Protocol Gateway and the Source-Integrity Gap,''
arXiv:2608.30083, 2026. (The \texttt{zk-attach/v0} envelope is the system
contribution of this paper.)

\bibitem{zkgen}
J.~Lauinger, J.~Ernstberger, and S.~Steinhorst, ``zkGen: Policy-to-Circuit
Transpiler,'' in \emph{IEEE International Conference on Blockchain and
Cryptocurrency (ICBC)}, 2024, pp.~620--624,
DOI 10.1109/ICBC59979.2024.10634440.

\bibitem{pco}
H.~Avni, S.~Dolev, A.~Yagudaev, and M.~Yung, ``Super-intelligence Survival
Guide: Verification via Proof-Carrying Output,'' IACR Cryptology ePrint
Archive, Report 2026/994, 2026.

\bibitem{zkagent}
L.~Wang, H.~Lou, C.~Li, Y.~Yu, and Y.~Hu, ``zkAgent: Verifiable LLM Agent
Execution via One-Shot Transcript Proofs,'' IACR Cryptology ePrint Archive,
Report 2026/199, 2026.

\bibitem{hiddencred}
J.~E.~Holt, R.~W.~Bradshaw, K.~E.~Seamons, and H.~Orman, ``Hidden
Credentials,'' in \emph{Proc.\ 2003 ACM Workshop on Privacy in the Electronic
Society (WPES)}, 2003, pp.~1--8, DOI 10.1145/1005140.1005142.

\bibitem{handshake}
D.~Balfanz, G.~Durfee, N.~Shankar, D.~K.~Smetters, J.~Staddon, and H.-C.~Wong,
``Secret Handshakes from Pairing-Based Key Agreements,'' in \emph{Proc.\ IEEE
Symposium on Security and Privacy (S\&P)}, 2003, pp.~180--196,
DOI 10.1109/SECPRI.2003.1199336.

\bibitem{xacmlpriv}
C.~A.~Ardagna, S.~De~Capitani~di~Vimercati, G.~Neven, S.~Paraboschi,
F.-S.~Preiss, P.~Samarati, and M.~Verdicchio, ``Enabling Privacy-Preserving
Credential-Based Access Control with XACML and SAML,'' in \emph{Proc.\ 10th
IEEE International Conference on Computer and Information Technology (CIT)},
2010, pp.~1090--1095, DOI 10.1109/CIT.2010.199.

\bibitem{sw05}
A.~Sahai and B.~Waters, ``Fuzzy Identity-Based Encryption,'' in \emph{Advances
in Cryptology -- EUROCRYPT 2005}, LNCS 3494, Springer, 2005, pp.~457--473,
DOI 10.1007/11426639\_27.

\bibitem{ksw08}
J.~Katz, A.~Sahai, and B.~Waters, ``Predicate Encryption Supporting
Disjunctions, Polynomial Equations, and Inner Products,'' in \emph{Advances in
Cryptology -- EUROCRYPT 2008}, Springer, 2008, pp.~146--162,
DOI 10.1007/978-3-540-78967-3\_9.

\bibitem{vdr}
H.~Chabanne, R.~Hugel, and J.~Keuffer, ``Verifiable Document Redacting,'' in
\emph{Computer Security -- ESORICS 2017}, LNCS 10492, Springer, 2017,
pp.~334--351, DOI 10.1007/978-3-319-66402-6\_20.

\end{thebibliography}

\end{document}
"
%
%  This document carries the material removed from the 12--14 page main
%  paper to meet the TDSC length limit: full proofs, the detailed
%  evaluation tables, policy-private-mode detail, related-work tables,
%  and reproduction instructions. It is self-contained: it restates the
%  results it proves and can be read without the main PDF.
% =====================================================================
\documentclass[10pt,journal,compsoc]{IEEEtran}

\usepackage{amsmath,amsfonts,amssymb,amsthm}
\usepackage{algorithmic}
\usepackage{algorithm}
\usepackage{array}
\usepackage[caption=false,font=normalsize,labelfont=sf,textfont=sf]{subfig}
\usepackage{textcomp}
\usepackage{stfloats}
\usepackage{url}
\usepackage{verbatim}
\usepackage{graphicx}
\usepackage{cite}
\usepackage{booktabs}
\usepackage{xcolor}
\usepackage{microtype}
\usepackage{hyperref}

\hyphenation{Proof-of-Policy zero-knowl-edge re-exec-ution pol-icy
com-pli-ance bind-ing veri-fi-a-bil-ity cer-ti-fi-cate non-zero-knowl-edge}

\setlength{\emergencystretch}{2em}

\newcolumntype{L}[1]{>{\raggedright\arraybackslash}p{#1}}
\newcommand{\tabfont}{\footnotesize}

\newtheorem{theorem}{Theorem}
\newtheorem{lemma}[theorem]{Lemma}
\newtheorem{proposition}[theorem]{Proposition}
\newtheorem{corollary}[theorem]{Corollary}
\newtheorem{definition}[theorem]{Definition}
\newtheorem{assumption}[theorem]{Assumption}
\newtheorem{remark}[theorem]{Remark}

% Work markers are no-ops in the supplementary material.
\newcommand{\need}[1]{}
\newcommand{\fixme}[1]{}
\newcommand{\wpmark}[1]{}

% Supplementary numbering: S1, S2, ... and Table S1, S2, ...
\renewcommand{\thesection}{S\arabic{section}}
\renewcommand{\thetable}{S\arabic{table}}

\begin{document}

\title{Supplementary Material for\\``Proof-of-Policy: Zero-Knowledge Compliance
Certificates for LLM Agent Outputs''}
\author{Author Name(s)}
\maketitle

\begin{abstract}
This supplement contains the full proofs, the detailed evaluation tables, the
policy-private-mode tables, the related-work positioning, and the reproduction
instructions omitted from the main paper for length. Section and table numbers
are prefixed with ``S''. The main paper states every result; this document
supplies the arguments and the per-configuration numbers behind them.
\end{abstract}

% =====================================================================
\section{Full Proofs}
\label{app:proofs}

This section restates the results proved in the main paper and gives complete
proofs. We use the main paper's notation; for convenience, the two objects the
proofs depend on are repeated here.

\begin{definition}[Guessing entropy; restated from the main paper]
\label{def:k}
For a secret distribution $\mathcal{D}$ and adversary side information $Z$, let
\begin{multline*}
\mathsf{Guess}(X \mid Z) \;\triangleq\; \max_{\mathcal{A}} \Pr[\mathcal{A}(Z) = X]\\
\;=\; \mathbb{E}_z\!\left[\max_x \Pr[X = x \mid Z = z]\right],
\end{multline*}
\[
k_X \;\triangleq\; -\log_2 \mathsf{Guess}(X \mid Z),
\]
the maximum ranging over all time-unbounded adversaries.
\end{definition}

\begin{assumption}[Adversary side information; restated from the main paper]
\label{as:Z}
The adversary knows the prompt, the model (possibly its weights), and the
decoding parameters.
\end{assumption}

\begin{proposition}[Conditional trilemma; restated from the main paper]
\label{thm:t1}
Fix $\lambda,\varepsilon>0$ and $q\ge 1$. Let $X$ be a secret with guessing
entropy $k_X<\log q+\log(1/\varepsilon)$ whose binding is \emph{guess-tight}:
$\beta_\rho(X)=H(\mathsf{DOM}\|\rho\|X)$ for a random oracle $H$ and public
randomness $\rho$ independent of $X$ and $Z$. Then:
\begin{enumerate}
\item[(1)] \emph{Public binding.} For any adversary that evaluates $\beta$ at
most $q$ times,
\[
\Pr\!\left[\mathcal{A}^\beta(Z,\beta(X))=X\right]
\le q\cdot\mathsf{Guess}(X\mid Z)+\mathrm{negl}(\lambda)
= q\cdot 2^{-k_X}+\mathrm{negl}(\lambda).
\]
Hiding with advantage $\varepsilon$ is \emph{certified} by this bound when
$k_X\ge\log q+\log(1/\varepsilon)$; for a skewed distribution this is a
sufficient condition, not a matching attack. The integer $q$ counts evaluations
of $\beta$, not time: the adversary may be computationally unbounded.
\item[(2)] \emph{Keyed binding.} For $\beta_s(X)=H(s,X)$ with $s$ held only by a
designated party and absent from the public values,
$\Pr[\mathcal{A}(Z,b)=X]\le\mathsf{Guess}(X\mid Z)+\mathrm{negl}(\kappa)$.
What the keyed form buys is hiding when
$k_X<\log q+\log(1/\varepsilon)$; its price is PV-2.
\item[(3)] \emph{Conditional trilemma.} Consider schemes whose binding is
publicly recomputable (PV-2), whose hiding is exact recovery, and whose binding
is guess-tight. Within the scheme class defined above, when
$k_X<\log q+\log(1/\varepsilon)$ such a scheme cannot satisfy binding, hiding
and PV-2 simultaneously.
\end{enumerate}
\end{proposition}

\begin{proof}[Proof of Proposition~\ref{thm:t1}]
\emph{(1).} Fix $z$ and let the adversary's query set be
$Q = \{q_1, \dots, q_q\}$. Decompose by first hit,
\[
\Pr[\exists i : q_i = X] \;=\; \sum_i \Pr\!\left[q_i = X \wedge \neg E_{<i}\right],
\]
where $E_{<i} \triangleq \{\exists j < i : q_j = X\}$.
Conditioned on $\neg E_{<i}$, every oracle answer $a_1,\dots,a_{i-1}$ already
revealed is uniform and independent of $X$ (no queried point equals $X$, and $H$
is a fresh random function), and $\beta(X)$ itself is uniform and independent of
$X$. Hence $q_i$---a deterministic function of $(z, \beta(X), a_{<i})$---is
\emph{independent of $X$}, so
\begin{multline*}
\Pr[q_i = X \mid Z = z]
= \sum_x \Pr[q_i = x]\,\Pr[X = x \mid Z = z] \\
\le \max_x \Pr[X = x \mid Z = z] = 2^{-k_z},
\end{multline*}
giving $\Pr[\exists i : q_i = X \mid z] \le q\,2^{-k_z}$. Taking the expectation
over $z$ and using $\mathbb{E}_z[2^{-k_z}] = \mathsf{Guess}(X \mid Z) = 2^{-k_X}$
from Definition~\ref{def:k} yields $q \cdot 2^{-k_X}$. An adversary that outputs
an unqueried value wins with probability at most $2^{-k_z} \le q\,2^{-k_z}$,
which the same bound absorbs. No $q\cdot2^{-\lambda}$ term is needed in the
exact-recovery game: a query that lands on a different $x$ with
$\beta(x)=y$ is a second preimage of $H$, not a win. The explicit
$q\cdot2^{-\lambda}$ is kept only when the win condition is consistency with the
published $b$ rather than exact recovery, or when $q$ is unbounded.

\emph{(2).} Without an oracle for $\beta_s$, the value $b=H(s,X)$ is uniform and
independent of $X$ in the random-oracle model for an adversary that does not know
$s$, so the best strategy is the prior mode and succeeds with probability
$\mathsf{Guess}(X\mid Z)$. An attempted enumeration must also guess $s$, at cost
$2^{-\kappa}$ per query. The bound is therefore
$\mathsf{Guess}(X\mid Z)+\mathrm{negl}(\kappa)$, with no $q$-fold amplification.
Part (2) buys hiding only below the public-form threshold
$\log q+\log(1/\varepsilon)$; when $k_X=0$ the prior mode already recovers $X$
and no binding helps.

\emph{(3).} Public recomputability gives (1). If (1) exceeds $\varepsilon$, the
only way to avoid it is a binding that a verifier cannot recompute, i.e.\ a keyed
binding, which by part (2) forfeits PV-2. The third leg---hiding plus PV-2
without binding---is excluded from the class of compliance schemes because a
certificate whose binding does not pin $X$ certifies nothing about the delivered
artifact.
\end{proof}

\noindent\emph{Preimage-class bound.} For an arbitrary deterministic public
$\beta$, conditioning on $\beta(X)$ gives
$\Pr[\mathcal{A}^\beta(Z,\beta(X))=X]\le
\mathbb{E}_y[\max_x\Pr[X=x\mid Z,\beta(X)=y]]$, with no hypothesis on $\beta$.
For $\beta=\mathrm{id}$ or an injective $\mathsf{encode}$ this equals $1$; for a
lossy encoding it equals the inverse of the preimage-class size.

\begin{theorem}[D-privacy from three conditions; restated from the main paper]
\label{thm:dpriv}
The scheme is $D$-private if (C1) there is a PPT simulator $S_{\mathrm{rest}}$
such that, for every admissible $(\pi,T,n)$, the public values other than the
commitment satisfy
$(h_\pi,S_{\mathrm{rest}}(pp,h_\pi,D(\pi),\mathrm{pub}_T,n))\approx_c
L_{\mathrm{pub}}(\pi,T,n)$; (C2) the proof system is witness-hiding, i.e.\
$L_{\mathrm{proof}}$ has a simulator that takes only the public statement;
and (C3) the policy commitment is hiding for all $D$-equal policies,
$(pp,h_{\pi_0})\approx_c(pp,h_{\pi_1})$. The advantage is
$O(\varepsilon_{\mathrm{pub}}+\varepsilon_{\mathrm{wh}}+\varepsilon_{\mathrm{com}})$.
\end{theorem}

\begin{proof}[Proof of Theorem~\ref{thm:dpriv}]
Fix an admissible tuple $(\pi_0,\pi_1,T,n)$ with $D(\pi_0)=D(\pi_1)$ and
$R(\pi_0,T)=R(\pi_1,T)$. Write $h_i=h_{\pi_i}$,
$R_i=L_{\mathrm{rest}}(\pi_i,T,n)$, $P_i=L_{\mathrm{proof}}(\pi_i,T,n)$ and
$\mathsf{Stmt}_i=(h_i,R_i)$. Consider the hybrids
\begin{align*}
H_0&=(h_0,R_0,P_0),\\
H_1&=(h_0,R_0,S_{\mathrm{proof}}(\mathsf{Stmt}_0)),\\
H_2&=(h_0,S_{\mathrm{rest}}(h_0,D,\mathrm{pub}_T,n),S_{\mathrm{proof}}(\mathsf{Stmt}_0)),\\
H_3&=(h_1,S_{\mathrm{rest}}(h_1,D,\mathrm{pub}_T,n),S_{\mathrm{proof}}(\mathsf{Stmt}_1)),\\
H_4&=(h_1,R_1,S_{\mathrm{proof}}(\mathsf{Stmt}_1)),\\
H_5&=(h_1,R_1,P_1).
\end{align*}
$H_0\approx H_1$ and $H_4\approx H_5$ cost $\varepsilon_{\mathrm{wh}}$ each by
C2; $H_1\approx H_2$ and $H_3\approx H_4$ cost
$\varepsilon_{\mathrm{pub}}$ each by C1. For $H_2\approx H_3$, split into
$H_{2a}=(h_1,S_{\mathrm{rest}}(h_1,D,\mathrm{pub}_T,n),S_{\mathrm{proof}}(\mathsf{Stmt}_0))$
and $H_{2b}=H_3$: $H_2\approx H_{2a}$ costs $\varepsilon_{\mathrm{com}}$ by C3,
and $H_{2a}\approx H_{2b}$ follows because
$\mathsf{Stmt}_0\approx_c\mathsf{Stmt}_1$ (C1 and C3 by triangle inequality) and
$S_{\mathrm{proof}}$ is a PPT function of the statement. The triangle inequality
gives $4\varepsilon_{\mathrm{pub}}+2\varepsilon_{\mathrm{wh}}+2\varepsilon_{\mathrm{com}}$,
which is $O(\varepsilon_{\mathrm{pub}}+\varepsilon_{\mathrm{wh}}+\varepsilon_{\mathrm{com}})$.
The game fixes one $T$ because two different artifacts are distinguished by
$B(n,\cdot)$ itself, which is response-side leakage and not a policy advantage.
\end{proof}

\begin{theorem}[Leakage-profile decomposition; restated from the main paper]
\label{thm:leak}
Let $V=(L_{\mathrm{pub}},L_{\mathrm{proof}},L_{\mathrm{side}})$ be the verifier's
view, where $L_{\mathrm{pub}}$ is the certificate's public values,
$L_{\mathrm{proof}}$ is what the proof artefact reveals about the witness
$w=(\pi,T,\mathrm{trace})$, and $L_{\mathrm{side}}$ is the adversary's prior
information $Z$. For any $q\ge 1$ and any adversary that evaluates $\beta$ at
most $q$ times,
\[
\Pr\!\left[\mathcal{A}^\beta(V)=X\right]
\;\le\; q\cdot\mathsf{Guess}\!\left(X\mid L_{\mathrm{pub}},L_{\mathrm{proof}},L_{\mathrm{side}}\right)
+\mathrm{negl}(\lambda),
\]
with the final output charged as one evaluation.
\end{theorem}

\noindent\emph{Proof of Theorem~\ref{thm:leak}.} The proof is the first-hit
decomposition of Proposition~\ref{thm:t1}(1) with $(Z,\beta(X))$ replaced by the
verifier view $V$. The step that changes is the uniformity of answers to queries
that have not hit $X$: it holds conditionally on $L_{\mathrm{pub}}$ and fails when
$L_{\mathrm{proof}}$ is a function of the witness. A witness-hiding tier restores
$L_{\mathrm{proof}}=\bot$ and hence the public-value bound; a \texttt{core}
artefact does not, and the bound can reach $1$.

\begin{theorem}[Delivery binding; restated from the main paper]
\label{thm:t3}
Fix a policy $\pi$. A delivery adversary $\mathcal{B}$ may, \emph{before}
seeing the receiver's challenge, precompute a set $S$ of accepted transcripts,
where each transcript covers at most $c$ challenges. The receiver then
contributes $f$ fresh, uniformly random bits $r$ that enter the proven
statement. Under the certificate's soundness (D1) and challenge binding
(D2/B-bind), the divergence advantage satisfies
\[
\mathsf{Adv}_{\mathrm{del}}(\mathcal{B}) \;\le\; |S| \cdot c \cdot 2^{-f}
\;+\; \mathrm{negl}(\lambda).
\]
The replay-only probability---acceptance without requiring divergence---is at
most $|S|\cdot c\cdot 2^{-f}$. If the receiver contributes no fresh randomness
($f = 0$; e.g.\ an empty nonce), a precomputed transcript is accepted with
probability $1$ and \emph{does not decay in any parameter of the system}.
\end{theorem}

\begin{proof}[Proof of Theorem~\ref{thm:t3}]
The receiver's freshness $r$ enters the proven statement through
$b = H(\mathsf{dom} \,\|\, r \,\|\, T)$. For a certificate in $S$ to be accepted
on $r$, the value $r$ must fall inside that certificate's challenge set
$A(\mathrm{cert})$, with $|A(\mathrm{cert})|\le c$; a union bound over $S$ gives
$\sum_{\mathrm{cert}\in S}|A(\mathrm{cert})|/2^f\le |S|\cdot c\cdot 2^{-f}$.
A certificate minted \emph{after} $r$ is seen can diverge only by exhibiting two
distinct preimages of the same $b$ under $B(r,\cdot)$, which is exactly a
challenge-binding failure and costs $\mathrm{negl}(\lambda)$. When $f = 0$ every
element of $S$ is a permanently valid replay, so the replay probability is $1$
and does not decay in any parameter of the system. The $f = 0$ branch is not
hypothetical: \texttt{circuits/\allowbreak types/\allowbreak src/\allowbreak
lib.rs} (lines 950--995) documents an empty nonce as a legal input that
``provides no replay protection.'' The $f=0$ separation holds under stateless
public verification with no external mechanism (no revocation table, trusted
clock, or hardware). The $c$ factor is tight when the challenge sets of the
transcripts in $S$ are disjoint and of size $c$.
\end{proof}

\begin{proposition}[Quadratic cost of the naive matcher]
\label{prop:naive}
For the pattern \texttt{a+b} and input \texttt{a}$^n$, the naive backtracking
matcher performs
$\sum_{i=1}^{n}(n-i+1)=n(n+1)/2=\Theta(n^{2})$ NFA steps, whereas the Pike VM
performs $\Theta(n)$. The adversarial input family is admissible because the
response bytes are chosen by the party being checked. This bounds the cost of
\emph{that} baseline, not of policy-compliance matching: the Pike VM is linear
on the same family, and no lower bound for the problem is claimed.
\end{proposition}
\begin{proof}
No input position contains \texttt{b}, so no start position can fail or accept
before the end of the string. The start at $i$ therefore takes exactly
$n-i+1$ steps, and the sum is $n(n+1)/2$. The Pike VM keeps one state per NFA
position and scans once.
\end{proof}

\begin{proposition}[NFA contract]
\label{prop:nfa}
For the supported regular-expression subset $P$---ASCII, ordered and disjoint
character-class ranges, Unicode scalar values, and the same compiled policy on
both sides---the Thompson construction preserves the language, the Pike VM
implements unanchored substring search, and the Python reference and Rust
circuit matchers agree on the same compiled NFA specification. Empty character
classes are rejected at parse time (fail-closed), so ``non-empty character
class'' is a parser contract rather than an extra theorem assumption.
\end{proposition}
\begin{proof}[Proof sketch]
Language preservation is the standard Thompson construction over the restricted
syntax. The Pike VM simulates all NFA states at one input position and advances
once per character, so it decides unanchored search in one pass. The two
matchers consume the same \texttt{ConstraintSpec} bytes and implement the same
transition relation; agreement is checked by an algorithm-level bisimulation
plus differential fuzzing against Python's \texttt{re}, not by line-by-line
formal verification.
\end{proof}

% =====================================================================
\section{Evaluation Details}
\label{app:eval-details}

This section carries the per-configuration tables and the detailed readings
omitted from the main paper. Every number is the content of the checked-in
result file named in the caption, as in the main paper.

\subsection{Cycle sweep: length versus rule count}
\label{sec:eval-cycles}

Response length was pushed from 200 to 100{,}000 characters and the rule count
from 1 to 6, giving 40 measurement points (\texttt{bench/bench\_cycles.py}; the
full matrix is in \texttt{bench/results/cycles.md}). The response text comes from
an evaluation artefact the repository's own end-to-end demo produced (372
characters, three segments deduplicated by content); longer points are tiled from
it. What is measured is therefore the \emph{curve} of scanning cost, not ``a real
100{,}000-character document''. Every point was first confirmed to match no rule:
a matching rule short-circuits, and the cost of a short-circuit says nothing.

\begin{table}[t]
\centering
\caption{Cycle count against length, per rule set. Five rule sets, six length
points each; $R^2 \approx 1.0000$ throughout. Source:
\texttt{bench/results/cycles.md}.}
\label{tab:cycles}
\tabfont
\begin{tabular}{@{}rrrrr@{}}
\toprule
Rules & \texttt{keyword} & \texttt{pattern} & Slope (cyc/char) & Intercept \\
\midrule
1 & 0 & 0 & 73.9      & 50{,}137 \\
2 & 1 & 0 & 105.3     & 65{,}942 \\
3 & 1 & 1 & 4{,}289.8 & 387{,}501 \\
4 & 2 & 1 & 4{,}048.8 & 405{,}785 \\
6 & 3 & 2 & 5{,}403.5 & 993{,}257 \\
\bottomrule
\end{tabular}
\end{table}

Three findings, in decreasing order of how much they change what may be claimed.

\emph{(1) With the rule set fixed, cycle count is exactly linear in length.}
Five rule sets, six length points each, $R^2 \approx 1.0000$ throughout
(Table~\ref{tab:cycles}).

\emph{(2) Scanning cost is dominated by the NFA, not by the rule count.} One
keyword rule costs about $+31.6$ cycles per character at the margin; one NFA
regular expression (an email pattern) costs about $+3{,}634$---roughly
$115\times$ more. Equivalently, a rule set containing a regular expression has a
per-character slope ($4{,}289.8$) $58\times$ that of a rule set with only string
rules ($73.9$).

\emph{(3) Unit prices are not additive: a unit price is a property of a rule
\emph{in context}, not of the rule.} Fitting the slopes as
$(\text{keyword},\text{pattern})$ unit prices gives $R^2 = 0.9169$ and a
\emph{negative} keyword unit price ($-477.2$ cycles per character). This is not
noise; two probes give the mechanism directly:

\begin{itemize}
  \item \textbf{Ordering probe} (same rules, same text, $L = 10^5$): declared
  order costs $541{,}342{,}069$ cycles against $570{,}736{,}989$ in reverse,
  a difference of $+5.43\%$, while \texttt{passed} is identical in both. \emph{The
  semantics of the decision do not depend on order; the cost does.}
  \item \textbf{Marginal probe} ($L = 10^5$, filling the one matrix cell the
  sweep could not reach---regular expressions with no keyword rule):
  $[\text{len},\text{pat}]$ over $[\text{len}]$ has a marginal unit price of
  $3{,}633.7$, whereas ``add one more keyword rule'' has two different answers:
  $+31.6$ alone, $+585.6$ in a policy that already contains a regular expression
  ($-153.7$ if inserted before it). The same action changes sign between
  $-154$ and $+586$.
\end{itemize}

The mechanism lies in zkVM memory allocation: every \texttt{ascii\_lower} and
every Pike VM match allocates, and allocator state depends on previous
allocations. We therefore quote \emph{per-rule-set slopes} (the table above,
directly citable) rather than any ``unit price $\times$ count'' expression. The
simple model $\text{cycles} \approx a\cdot|T|\cdot\text{rules} + b$ that this
work anticipated is refuted by the data ($R^2 = 0.872$, negative intercept), and
we report that too: \emph{a model refuted by its own data is itself an evaluation
result.}

\subsection{Matcher ablation and proof-path cost}
\label{sec:eval-ablation}

\emph{Construction.} Pattern \texttt{a+b} against input
\texttt{a}$^{\times n}$ containing no \texttt{b}, so the match is necessarily
absent and both implementations must scan to the end. The naive implementation
reinvents and restarts at every start position, giving $\Theta(n^2)$; the Pike VM
decides in a single linear pass (\texttt{bench/bench\_ablation.py}).
Table~\ref{tab:ablation} reports the cycles. Read per-character cost rather than
absolute figures. Doubling $n$ leaves the Pike VM's per-character cost
essentially unchanged ($-11.1\%$, $-6.4\%$) while the naive path's nearly doubles
($+96.4\%$, $+98.6\%$). At $n = 400$ the naive path is already $90\times$ slower,
and that ratio doubles again with each doubling of $n$. \emph{This is an attack
surface}: the input length for policy compliance is chosen externally---an
agent's reply, a concatenated session---so an adversary who needs only to write a
long reply drives a naive implementation's cost quadratically. The comparison is
meaningful only because the two paths decide identically, which
\texttt{tests/test\_ablation.py} pins down (Python and Rust
\texttt{match\_search} $\equiv$ \texttt{match\_search\_naive}). In the systems we
surveyed, we are not aware of ablation evidence specific to policy-compliance
matching.

\begin{table}[t]
\centering
\caption{Matcher ablation for pattern \texttt{a+b} on input \texttt{a}$^n$.
Cycles are zkVM execution counts. Source:
\texttt{bench/results/ablation.md}.}
\label{tab:ablation}
\tabfont
\begin{tabular}{@{}rrrrr@{}}
\toprule
$n$ & Naive cycles & Pike cycles & Ratio \\
\midrule
100 & 12{,}382{,}720  & 643{,}894   & $19.2\times$ \\
200 & 48{,}638{,}946  & 1{,}145{,}209 & $42.5\times$ \\
400 & 193{,}207{,}130 & 2{,}143{,}617 & $90.1\times$ \\
\bottomrule
\end{tabular}
\end{table}

\paragraph{Proof-path cost.} Re-running the same construction through the
compliance driver (\texttt{bench/results/dos\_proof.md}) reproduces the ratios
on the execution path---$643{,}906$ against $12{,}382{,}732$ cycles at $n=100$,
rising to $90.1\times$ at $n=400$---and then prices the proving path. Proving is
where the quadratic term meets the prover's fixed memory floor: the Pike path at
$n=100$ succeeds at $122.2$\,s and $11{,}695$\,MiB; the naive path at $n=25$
succeeds at $153.2$\,s and $12{,}137$\,MiB; at $n=50$ it is
\textbf{OOM/SIGKILL} after $226.5$\,s at $13{,}772$\,MiB peak. We report the OOM
as a result rather than a discarded run: it is the point at which the
adversarial quadratic term exhausts a $14.6$\,GiB machine, and it is why the
$n\ge100$ naive points and all $n=400$ proving points are left to a cloud
machine. The cycle counts differ from these execution-path figures by about one
part in $10^5$, because they come from a separate driver build; the ratios
agree. These proving runs predate the NFA fix discussed in the main paper and
carry that binary's digest; the fix changes the guest and the verification key,
so the figures are not mixed with any post-fix run.

\begin{table}[t]
\centering
\caption{Proof-path denial of service on a 24-core / 14.6\,GiB machine. The
\texttt{naive} row at $n=50$ is an observed out-of-memory kill, retained as a
result. Source: \texttt{bench/results/dos\_proof.md}.}
\label{tab:dos}
\tabfont
\setlength{\tabcolsep}{3pt}
\begin{tabular}{@{}lrrrl@{}}
\toprule
Point & Cycles & Prove (s) & RSS (MiB) & Result \\
\midrule
pike $n{=}100$ & 643{,}906 & 122.2 & 11{,}695.1 & \checkmark \\
naive $n{=}25$ & --- & 153.2 & 12{,}137.3 & \checkmark \\
naive $n{=}50$ & --- & 226.5 & 13{,}771.6 & OOM/SIGKILL \\
\bottomrule
\end{tabular}
\end{table}

\subsection{Composition: four negative controls}
\label{sec:eval-compose}

The leakage-profile theorem of the main paper prices each half; it does not say
that conjoining two valid halves yields a valid composite. We built the four
cells of Table~\ref{tab:compose} against the composite verifier's binding layer,
with the cryptographic validity of the two sub-proofs covered separately by the
end-to-end tests. Only cell~(a) satisfies the strong form, and only because
$x_U=\mathsf{derive}(T_U)$; the implementation claims only the weak form of the
composition proposition in the main paper and its E-1 remark. Cell~(c) is the
one that matters: the current verifier binds the inference half to
$\mathsf{derive}(T_U)$, not to the input $x_U$ the user submitted, so it accepts
a certificate that proves the weak form
$\exists x: x=\mathsf{derive}(T_U)\wedge T_U=M(x)$ while the strong form
$x=x_U$ remains unproved. This is proof debt E-1, not a test artefact; it is the
concrete reason composition alone does not close the gap.

\begin{table}[t]
\centering
\caption{Composition negative controls. \texttt{ok} is the verifier's binding
check; \texttt{satisfied} is the policy outcome. Source:
\texttt{bench/results/composition\_controls.md}.}
\label{tab:compose}
\tabfont
\setlength{\tabcolsep}{3pt}
\begin{tabular}{@{}lL{0.34\columnwidth}ccL{0.23\columnwidth}@{}}
\toprule
Cell & Construction & ok & sat. & Reading \\
\midrule
(a) & same $n$; $x_U=\mathsf{derive}(T_U)$; \texttt{delegated} empty &
\checkmark & \checkmark & true composition \\
(b) & the two halves use different nonces & $\times$ & $\times$ &
rejected by the four-way binding check \\
(c) & $x_U\neq\mathsf{derive}(T_U)$; inference binds only the derived input &
\checkmark & \checkmark & \textbf{weak form only}; $x=x_U$ not landed (E-1) \\
(d) & \texttt{delegated} non-empty & \checkmark & $\times$ &
true but unsatisfied; fail closed \\
\bottomrule
\end{tabular}
\end{table}

\subsection{Gateway signing: measured cost, blocked circuit}
\label{sec:eval-gateway}

Table~\ref{tab:gateway} prices the trust boundary the main paper states rather
than hides. Off-circuit Ed25519 verification costs $0.094$, $0.844$ and
$10.165$\,ms per receipt batch at $1$, $10$ and $100$ receipts, with a
per-receipt marginal near $0.08$--$0.10$\,ms; the session-end \texttt{ToolSeal}
is $0.074$\,ms to sign, $0.184$\,ms to verify and $345$\,B. A concatenated
2-of-3 threshold prototype signs in $0.173$\,ms, merges in $0.001$\,ms, verifies
in $0.508$\,ms and produces a $981$\,B certificate ($759$\,B when only the two
required signatures are carried). The threshold semantics are machine-checked:
any two distinct signatures pass, one signature is rejected, and two signatures
under the same key id do not count as two.

\begin{table}[t]
\centering
\caption{Gateway signing cost on a 24-core / 14.6\,GiB machine ($n=300$ per
row). Rows are measured medians. The in-circuit row is \emph{blocked}: the
circuit contains no Ed25519 verifier by design, and no cycle, time or memory
figure is reported. Source: \texttt{bench/results/gateway\_sig.md}.}
\label{tab:gateway}
\tabfont
\begin{tabular}{@{}lrrrr@{}}
\toprule
Configuration & Sign (ms) & Verify (ms) & Cert (B) & $n$ \\
\midrule
Ed25519, 1 receipt & 0.058 & 0.094 & 237.0 & 300 \\
Ed25519, 10 receipts & 0.415 & 0.844 & 288.3 & 300 \\
Ed25519, 100 receipts & 4.885 & 10.165 & 295.2 & 300 \\
\texttt{ToolSeal} (10 receipts) & 0.074 & 0.184 & 345 & 300 \\
1-of-1 & 0.069 & 0.173 & 648 & 300 \\
2-of-3 (3 signatures) & 0.173 & 0.508 & 981 & 300 \\
\midrule
In-circuit Ed25519 & \multicolumn{3}{c}{blocked: no guest, no number} & --- \\
\bottomrule
\end{tabular}
\end{table}

The in-circuit row is the honest blank. Verifying Ed25519 inside the zkVM is not
in the circuit---\texttt{circuits/} deliberately keeps the signature field out
of the circuit view---and building a guest to measure it exceeds this
experiment's file ownership, so cycles, proving time and peak memory are
\emph{null}, not zero. We therefore do not infer a circuit cost from the
off-circuit cost, and the trust root still contains the gateway key. What the
threshold prototype buys is not a cheaper verification but a weaker single point
of trust: one key compromised is no longer enough, at roughly three times the
signing work and certificate size, and with a new obligation to argue that the
gateways agree on the chain tail.

Putting Ed25519 inside the guest is a design option, not a measured one. SP1
6.7.0 has no Ed25519 precompile; the guest change would rotate the policy vkey
and add public values to the outcome, invalidating existing certificates, and we
did not build or time it. The only gateway-side variant measured is the 2-of-3
threshold prototype. It removes the single-key trust but not gateway collusion
or tail truncation. An authenticated gateway roster, HSM key custody and a
trusted time source remain open.

\subsection{Cross-system baselines: 49 cells, 46 populated}
\label{sec:eval-baselines}

Table~\ref{tab:baselines} places PoP against six alternatives across seven
dimensions: 49 cells, of which 46 are populated (measured or cited) and 3 are
unavailable. Every populated cell is tagged $[m]$ (measured in this repository
or this study) or $[c]$ (cited from the named paper). The three measured arms are
PoP, the signed-Merkle log and a simulated trusted server; zkAgent, NiyamAI,
Zaguk and PCO are cited. Zaguk and PCO did not yield comparable performance
numbers in our environment---the former's proceedings PDF was unreachable and
the latter reports no performance table---so those cells are em dashes, not
estimates. The seven dimensions impose different proof obligations; the table is
a statement of \emph{what each system obliges the verifier to trust}, not a
single ranking.

\begin{table*}[t]
\centering
\caption{Cross-system comparison. $[m]$ = measured here; $[c]$ = cited from the
source named in \texttt{bench/results/baselines.md}; em dash = not available,
not guessed. The trust root is per obligation, not one column: PoP's response
and policy proofs trust the vkey, while tool receipts additionally trust the
gateway key.}
\label{tab:baselines}
\tabfont
\setlength{\tabcolsep}{3pt}
\scriptsize
\begin{tabular}{@{}L{0.11\textwidth}L{0.13\textwidth}L{0.13\textwidth}L{0.12\textwidth}L{0.11\textwidth}L{0.10\textwidth}L{0.10\textwidth}L{0.08\textwidth}@{}}
\toprule
Dimension & PoP & signed-Merkle & TEE (sim) & zkAgent & NiyamAI & Zaguk & PCO \\
\midrule
Trust root & vkey + crypto $[m]$; tool-receipt gateway key & Ed25519 gateway $[m]$ &
server / TEE $[m]$ & zkVM + provider $[c]$ & SNARK + judge $[c]$ &
ezkl + commitment $[c]$ & authority + ledger $[c]$ \\
Latency & 40.5--100.3\,s core; 279--344\,s \texttt{gr\-oth16}; 89.84\,ms verify $[m]$ &
0.132\,ms sign; 2.131\,ms append; 0.322\,ms check $[m]$ &
0.139\,ms sign; 0.102\,ms verify $[m]$ &
194.4\,s / 0.42\,s $[c]$ & 1.7\,s / 51\,ms $[c]$ &
not comparable $[c]$ & --- \\
Size & 2.79\,MB core; 1{,}874\,B \texttt{gr\-oth16} $[m]$ & 648\,B + 761\,B $[m]$ &
648\,B $[m]$ & 96\,MiB (Shout) $[c]$ & 18.6\,KB $[c]$ & --- & --- \\
Response privacy & committed $[m]$ & none $[m]$ & none $[m]$ &
committed $[c]$ & not reported $[c]$ & yes $[c]$ & yes $[c]$ \\
Policy privacy & public values salted; proof artefact open $[m]$ & none $[m]$ & none $[m]$ &
n/a $[c]$ & contract commitment $[c]$ & yes $[c]$ & partial $[c]$ \\
Needs $T$ plaintext & no $[m]$ & yes $[m]$ & no $[m]$ & no $[c]$ &
unclear $[c]$ & no $[c]$ & no $[c]$ \\
Needs $\pi$ plaintext & no $[m]$ & yes $[m]$ & no $[m]$ & n/a $[c]$ &
unclear $[c]$ & no $[c]$ & partial $[c]$ \\
\bottomrule
\end{tabular}
\end{table*}

The asymmetry the table exposes is the one the signature baseline already
showed: the cheap arms require the verifier to re-run the judgment on plaintext,
while the private arms move that work into a proof or into a trusted party. PoP
is on the proof side, with a measured cost four to five orders of magnitude
above the signature arm and a trust root that, unlike the TEE arm, contains no
hardware.

\subsection{Issuance: detection, not prevention}
\label{sec:eval-issuance}

A provider that never issues a certificate is indistinguishable, from the
verifier's side, from a provider with nothing to certify; no cryptographic
layer forces issuance. What a transparency log changes is the \emph{detection}
rate, and only under an external obligation to append before delivery. In a
simulation with $N=2000$ deliveries, $m=100$ withheld violations and a sampling
rate $p=0.05$, an audit that samples the \emph{delivered} artifacts detects
$0.9948$ of the withheld set in exact arithmetic ($0.9940$ by Monte Carlo over
500 draws). An audit that samples the \emph{log} detects $0.0000$, because the
withheld deliveries have no log entry at all: the sampling frame is wrong, not
unlucky.

\begin{table}[t]
\centering
\caption{Issuance detection under two sampling frames ($N=2000$, $m=100$,
$p=0.05$, seed 20261101). Exact is $1-\binom{N-m}{s}/\binom{N}{s}$ with $s=100$;
Monte Carlo uses 500 draws. Source: \texttt{bench/results/issuance.md}.}
\label{tab:issuance}
\tabfont
\begin{tabular}{@{}lrr@{}}
\toprule
Sampling frame & Exact & Monte Carlo \\
\midrule
Delivered artifacts (correct frame) & 0.9948 & 0.9940 \\
Log entries (wrong frame) & 0.0000 & 0.0000 \\
\bottomrule
\end{tabular}
\end{table}

Two boundaries travel with these numbers. First, the result quantifies
\emph{detection}, never \emph{prevention}: whether a detected violation is
punished depends on penalties and incentives outside the model. Second, the
sufficiency direction assumes an append-before-delivery obligation, a log that
is append-only and publicly readable, inclusion proofs, and an external
mechanism against split views. The prototype implements append-only semantics,
inclusion proofs, optional tree-head signatures and checkpoints; it does
\emph{not} implement consistency proofs or gossip, so split-view resistance is
an external assumption, not a delivered property. Three premises travel with the
result: (F-1) the append-before-deliver obligation is external to cryptography;
(F-3) the log entry must bind the delivered $T$ (or a receipt), not only
$H(\mathrm{cert})$; and (F-4) the audit sample frame is delivered artifacts, not
log entries. The $1-(1-p)^m$ form assumes independent sampling and
detection-on-hit.

\subsection{Chat-template rerun and refusal proxy}
\label{sec:eval-k-chat}

A chat-template rerun of 20 prompts at $N=10$ returned $k=3.32=\log_2 10$ at
temperature $0.7$ and $k=0$ at temperature $0$; its refusal proxy rose from
$15\%$ to $35\%$ ($7/10$ unsafe prompts) while the response-violation rate fell
from $100\%$ to $90\%$. The $90\%$ is dominated by keyword-policy false
positives---$9/10$ safe responses were flagged---and is not a model harm rate;
XSTest labels are prompt-level and the refusal proxy is a prefix heuristic.

% =====================================================================
\section{Policy-Private Mode Details}
\label{app:private-policy}

\paragraph{Three parties.} ``Three'' means \emph{authority / operator / third
party}; an auditor is a named instance of the third party, not a fourth trust
anchor. Table~\ref{tab:wp4-parties} states the model.

\begin{table}[t]
\centering
\caption{The three-party model of policy-private mode.}
\label{tab:wp4-parties}
\tabfont
\setlength{\tabcolsep}{3pt}
\begin{tabular}{@{}L{0.16\columnwidth}L{0.14\columnwidth}L{0.62\columnwidth}@{}}
\toprule
Party & Honest? & Holds; wants \\
\midrule
Authority $R$ & \textbf{Yes}, by (P-PP) & Holds $\pi$ and the issuing record;
wants $\pi$ actually enforced, and wants to be able to check that it was. \\
Operator $O$ & \textbf{No} (the adversary) & Holds $\pi$, $T$, the proving key;
wants $\pi$ hidden from third parties while still being able to certify. \\
Third party $V$ & Yes & Holds only public values---\texttt{policy\_\allowbreak hash},
\texttt{circuit\_\allowbreak hash}, \texttt{response\_\allowbreak commitment},
\texttt{passed}, the proof; wants to confirm that \emph{some issued} policy
passed, and must \emph{not} learn $\pi$. \\
\emph{(Auditor $V^{*}$)} & Yes & As $V$, plus $\pi$ once $O$ discloses it;
wants ex-post accountability. \\
\bottomrule
\end{tabular}
\end{table}

Two boundaries must be stated with the model. First, secrecy of $\pi$ comes from
a \emph{data-flow fact}---$\pi$ never enters the public values---and not from
zero-knowledge; any rule that writes a fragment of $\pi$ into the public values
breaks secrecy on that fragment. Second, accountability for $V^{*}$ is
\emph{voluntary}: $V^{*}$ can check only when $O$ chooses to disclose $\pi$ and
the salt. The mode offers no forced disclosure; that would require the
release-gating family discussed in \texttt{docs/tdsc-plan.md}~\S1.6(b), which we
do not claim.

\paragraph{What each verifier may conclude.} Table~\ref{tab:wp4-capability} is
the contribution of this subsection: the cryptography is downstream of it, and
the table survives even if the cryptography only partially lands.

\begin{table*}[t]
\centering
\caption{What each verifier may and may not conclude in policy-private mode.}
\label{tab:wp4-capability}
\tabfont
\setlength{\tabcolsep}{4pt}
\begin{tabular}{@{}L{0.13\textwidth}L{0.30\textwidth}L{0.30\textwidth}L{0.19\textwidth}@{}}
\toprule
Verifier & \textbf{May} conclude & \textbf{May not} conclude & Rests on \\
\midrule
Third party $V$ & (i) the commitment binds $T$
(Definition~1 of the main paper); (ii) \emph{some} policy with commitment $h_\pi$
passed; (iii) $h_\pi$ appears in the issuing record & (iv) the \emph{content} of
$\pi$; (v) how many rules $\pi$ has, or of which classes; (vi) ``the operator
enforces the regulation \emph{I} care about''---only ``some issued regulation'' &
(P-PP); Assumption~\ref{as:Z}; the $k$ of Proposition~\ref{thm:t1} \\
End user (also receiver) & As $V$, and additionally: the proved $T$ is the
string it received, since its nonce is bound (Theorem~\ref{thm:t3}, requires
$f>0$)$^{\dagger}$ & Only the string it received; \emph{not} what was delivered
to another downstream agent & $f>0$ \\
Authority $R$ & Everything; it holds $\pi$, the salt and can recompute $h_\pi$ &
--- & --- \\
Auditor $V^{*}$ (after disclosure) & (vii) $\mathsf{verify\_commit}(\pi,s)=
\texttt{policy\_\allowbreak hash}$, i.e.\ what ran was the issued $\pi$;
(viii) the disclosed violation record is consistent & (ix) the undisclosed part
of $T$ & voluntary disclosure \\
On-chain verifier & (i)--(iii) only; it hashes and verifies signatures, and does
not interpret policy semantics & (iv)--(ix) & (P-PP) \\
\bottomrule
\end{tabular}
\end{table*}

The $V$ row is the reason the table exists. A third party does not obtain
``compliance''; it obtains ``some issued policy was satisfied''. Conflating the
two is the sentence a referee would use to reject the mode, and preventing that
sentence is what the table is for. $^{\dagger}$\,The end-user binding is
currently a certificate-layer field; the circuit still computes the
verifier-domain binding, so the end-to-end form needs a circuit change and a new
vkey.

% =====================================================================
\section{Related-Work Positioning}
\label{app:related}

Table~\ref{tab:related} records the slot each lineage occupies---which is also,
read the other way, what the main paper does not have to claim.
Table~\ref{tab:zkagent} places PoP against zkAgent: the two prove different
obligations, so the rows are a statement of orthogonality, not a race.

\begin{table}[t]
\centering
\caption{Positioning against the five prior-art lineages. The right-hand column
names what each lineage already occupies, and therefore what we do not claim.}
\label{tab:related}
\tabfont
\setlength{\tabcolsep}{3pt}
\begin{tabular}{@{}L{0.26\columnwidth}L{0.42\columnwidth}L{0.26\columnwidth}@{}}
\toprule
Lineage & Representative work & Slot already occupied \\
\midrule
Policy $\to$ circuit & zkGen \cite{zkgen}, Circuitree, Prezta & DSL-to-circuit
translation \\
Agent/LLM compliance & NiyamAI, Zaguk, PCO \cite{pco}, WITNESS & policy-level
private-mode kernel \\
Model computation & DeepProve, zkGPT, zkAgent \cite{zkagent} & inference
integrity \\
Attacks / negative results & Hollow-LLM \cite{hollowllm}, zk-attach
\cite{zkattach} & naming the gaps \\
Non-ZK industry baseline & \texttt{@brivora/verify}, VAOL & signed audit logs \\
\bottomrule
\end{tabular}
\end{table}

\begin{table}[t]
\centering
\caption{PoP against zkAgent. The two prove different obligations, so the rows
are a statement of orthogonality, not a race.}
\label{tab:zkagent}
\tabfont
\setlength{\tabcolsep}{3pt}
\begin{tabular}{@{}L{0.22\columnwidth}L{0.38\columnwidth}L{0.30\columnwidth}@{}}
\toprule
& \textbf{zkAgent} \cite{zkagent} & \textbf{PoP} \\
\midrule
Obligation & provider executed the declared model and tools (inference integrity)
& response/trace satisfies policy $\pi$ (compliance) \\
Depends on the model & yes (weights, arithmetisation) & \textbf{no} \\
Cost grows with & model size, sequence length & response length; rule-class
composition (per-rule-set slope), not rule count \\
Content privacy & none & committed, selective disclosure \\
Hardware reported & 32 cores / 512\,GB & 16 cores / 90\,GiB \\
Proving & 194.4\,s (Shout, $T{=}512$) & 40.5--100.3\,s core; 279--344\,s \texttt{gr\-oth16} \\
Proof size & 96\,MiB (Shout); 3.1\,MiB (LogUp) & 2.79\,MB core; 1{,}874\,B \texttt{gr\-oth16} \\
Verification & 0.42\,s (Shout); 0.038\,s (LogUp) & $\approx$0.090\,s \\
\bottomrule
\end{tabular}
\end{table}

We deliberately do not claim that PoP is faster, and the temptation to do so is
worth naming because the raw numbers invite it. The correct statement is
narrower, and we make only it: \emph{for the compliance obligation}, PoP
completes in a proving time of the same order as zkAgent's, with a smaller memory
footprint and a comparable proof size, without touching the model---and it adds
committed content privacy that zkAgent does not provide. Two caveats must travel
with those figures. The hardware is not matched (16 cores and a 90\,GiB cgroup
against 32 cores and 512\,GB), and PoP's ceiling is memory-bound rather than
cycle-bound, so its feasible region is a non-rectangular set of measured points
rather than a curve. The comparison is of orders of magnitude under unmatched
hardware, not a race. Source: \texttt{bench/comparison\_zkagent.md}, which
itself warns against reading the rows as ``PoP is faster.''

% =====================================================================
\section{Reproduction}
\label{app:repro}

The full walkthrough---environment, a single compliance proof, cross-checking
against the reference layer, private mode, certificate issuance, the one-command
demo, and third-party verification---is maintained in \texttt{docs/reproduce.md}
together with the exact commands and their expected output. This section records
only what a reader needs in order to judge the claims of the main paper's
evaluation.

\paragraph{Environment.} Ubuntu or WSL2 on x86-64, no GPU. The Rust/SP1 side
needs Rust stable, \texttt{cargo-prove}~v6.7.0, Go $\ge$\,1.24, and the usual
native toolchain; the reference layer is pure Python~3.10+ and runs on its own.
Memory, not time, is the binding constraint: the CPU prover carries a fixed
floor of approximately 10.5\,GiB on the cloud cgroup (10.15\,GiB under
default local prover options; \texttt{NUM\_CORE\_WORKERS=1},
\texttt{CORE\_BUFFER\_SIZE=1} lowers it to $\approx$9.3\,GiB) regardless of the
trace; memory, not time, sets the ceiling in the proof-cost table of the main
paper.

\paragraph{What reproduces without the prover.} The reference layer, the policy
compiler, and the whole test suite run in seconds and need no Rust. This is
enough to reproduce the non-ZK baseline of the main paper, to inspect the
certificate and inclusion proofs, and to check the cross-layer contract---that
the Python reference implementation and the zkVM guest consume the same
compiled policy bytes.

\paragraph{What needs the prover.} Every timing, size, and memory figure in the
cycle, ablation, denial-of-service, and proof-cost evaluations, and the
end-to-end result, requires a real proof. The measurements were taken on a
16-core CPU cgroup with 90\,GiB of RAM, which is what places the ceiling
reported there: the successes and the out-of-memory configurations are both
properties of that machine.

\paragraph{Where each number comes from.} Every measurement in the main paper's
evaluation is the content of a checked-in result file with a machine-readable
companion: \texttt{bench/results/} holds \texttt{cycles}, \texttt{ablation},
\texttt{scale-cloud}, \texttt{scale-6x2000}--\texttt{scale-6x5000},
\texttt{eu\_ai\_act\_proof}, \texttt{privacy\_surface\_real},
\texttt{k\_entropy\_largeN}, \texttt{k\_entropy\_crossmodel},
\texttt{policy\_corpus\_infer}, \texttt{agent\_e2e\_deepseek},
\texttt{baseline}, \texttt{verify}, \texttt{semantic}, \texttt{streaming},
\texttt{compose}, \texttt{dos\_proof}, \texttt{gateway\_sig},
\texttt{baselines}, \texttt{issuance}, and \texttt{prover\_knobs}, each as a
\texttt{.json} alongside its \texttt{.md}. The reproduction guide records the
commit, the hardware, and the prover binary digest for each run, so that a
reported figure can be traced to one specific execution rather than to a state
of the tree.

\begin{thebibliography}{99}

\bibitem{hollowllm}
C.~Gong, B.~Liu, and M.~Li, ``Hollow-LLM Attack: Computationally Trivial
Weights in Zero-Knowledge Verification of LLM Inference,'' in \emph{IEEE
Symposium on Security and Privacy (S\&P)}, 2026, pp.~4339--4355,
DOI 10.1109/SP63933.2026.00258; arXiv:2607.28884.

\bibitem{zkattach}
A.~S.~Gopalakrishna, ``Zero-Knowledge Predicate Proofs Between AI Agents: A
Measured, Cross-Protocol Gateway and the Source-Integrity Gap,''
arXiv:2608.30083, 2026. (The \texttt{zk-attach/v0} envelope is the system
contribution of this paper.)

\bibitem{zkgen}
J.~Lauinger, J.~Ernstberger, and S.~Steinhorst, ``zkGen: Policy-to-Circuit
Transpiler,'' in \emph{IEEE International Conference on Blockchain and
Cryptocurrency (ICBC)}, 2024, pp.~620--624,
DOI 10.1109/ICBC59979.2024.10634440.

\bibitem{pco}
H.~Avni, S.~Dolev, A.~Yagudaev, and M.~Yung, ``Super-intelligence Survival
Guide: Verification via Proof-Carrying Output,'' IACR Cryptology ePrint
Archive, Report 2026/994, 2026.

\bibitem{zkagent}
L.~Wang, H.~Lou, C.~Li, Y.~Yu, and Y.~Hu, ``zkAgent: Verifiable LLM Agent
Execution via One-Shot Transcript Proofs,'' IACR Cryptology ePrint Archive,
Report 2026/199, 2026.

\bibitem{hiddencred}
J.~E.~Holt, R.~W.~Bradshaw, K.~E.~Seamons, and H.~Orman, ``Hidden
Credentials,'' in \emph{Proc.\ 2003 ACM Workshop on Privacy in the Electronic
Society (WPES)}, 2003, pp.~1--8, DOI 10.1145/1005140.1005142.

\bibitem{handshake}
D.~Balfanz, G.~Durfee, N.~Shankar, D.~K.~Smetters, J.~Staddon, and H.-C.~Wong,
``Secret Handshakes from Pairing-Based Key Agreements,'' in \emph{Proc.\ IEEE
Symposium on Security and Privacy (S\&P)}, 2003, pp.~180--196,
DOI 10.1109/SECPRI.2003.1199336.

\bibitem{xacmlpriv}
C.~A.~Ardagna, S.~De~Capitani~di~Vimercati, G.~Neven, S.~Paraboschi,
F.-S.~Preiss, P.~Samarati, and M.~Verdicchio, ``Enabling Privacy-Preserving
Credential-Based Access Control with XACML and SAML,'' in \emph{Proc.\ 10th
IEEE International Conference on Computer and Information Technology (CIT)},
2010, pp.~1090--1095, DOI 10.1109/CIT.2010.199.

\bibitem{sw05}
A.~Sahai and B.~Waters, ``Fuzzy Identity-Based Encryption,'' in \emph{Advances
in Cryptology -- EUROCRYPT 2005}, LNCS 3494, Springer, 2005, pp.~457--473,
DOI 10.1007/11426639\_27.

\bibitem{ksw08}
J.~Katz, A.~Sahai, and B.~Waters, ``Predicate Encryption Supporting
Disjunctions, Polynomial Equations, and Inner Products,'' in \emph{Advances in
Cryptology -- EUROCRYPT 2008}, Springer, 2008, pp.~146--162,
DOI 10.1007/978-3-540-78967-3\_9.

\bibitem{vdr}
H.~Chabanne, R.~Hugel, and J.~Keuffer, ``Verifiable Document Redacting,'' in
\emph{Computer Security -- ESORICS 2017}, LNCS 10492, Springer, 2017,
pp.~334--351, DOI 10.1007/978-3-319-66402-6\_20.

\end{thebibliography}

\end{document}
